% Aide au système immunitaire — SVTEEHB 3ème — Cours signé OnBuch+
% Programme officiel MINESEC. Compiler avec : tectonic N08-aide-systeme-immunitaire.tex
\newcommand{\DOCMATIERE}{SVTEEHB}
\newcommand{\DOCNIVEAU}{3ème}
\newcommand{\DOCMODULE}{SVTEEHB – classe de 3ème}
\newcommand{\DOCLECON}{N08}
\newcommand{\DOCTITRE}{Aide au système immunitaire}
% OnBuch+ — préambule « cours signés » ENRICHI (compilé avec tectonic / XeLaTeX)
% Système visuel « Pop » d'OnBuch+ : fond crème (jamais blanc), bordures encre
% épaisses, ombres « dures » décalées, titres Archivo Black en capitales.
% Version MATHÉMATIQUES : graphes (pgfplots/tikz), boîtes pédagogiques
% (définition, propriété, méthode, attention, le savais-tu, exercice résolu,
% exercice, TP, bilan), page de garde et signature OnBuch+ ancrée en bas.
%
% Macros attendues dans main.tex AVANT \input{preamble} :
%   \DOCMATIERE  \DOCNIVEAU  \DOCMODULE  \DOCLECON (numéro)  \DOCTITRE
\documentclass[11pt]{article}

\usepackage{fontspec}
\usepackage[french]{babel}
\usepackage[a4paper,margin=2cm,top=3.4cm,bottom=2.8cm,headheight=1.4cm,footskip=1.3cm]{geometry}
\usepackage{amsmath,amssymb,amsfonts}
\usepackage{mathtools} % \xmapsto, \coloneqq... souvent utilisés par les modèles
\usepackage{cancel} % simplifications barrées (bilans de forces)
% \abs{z} (module, valeur absolue) et \norm{u} (norme), utilisés par les modèles
\providecommand{\abs}{}\let\abs\relax\DeclarePairedDelimiter{\abs}{\lvert}{\rvert}
\providecommand{\norm}{}\let\norm\relax\DeclarePairedDelimiter{\norm}{\lVert}{\rVert}
\usepackage[table]{xcolor} % [table] : fournit aussi \rowcolors (alternance de lignes)
\usepackage{pagecolor}
\usepackage{tikz}
\usetikzlibrary{arrows.meta,positioning,calc,shapes.geometric,shapes.misc,shapes.arrows,decorations.pathmorphing,decorations.pathreplacing,decorations.markings,patterns,shadows,mindmap,trees,fit,backgrounds,matrix,intersections,angles,quotes}
\usepackage{pgfplots}
\usepackage[version=4]{mhchem} % équations nucléaires \ce{^{235}_{92}U}
\usepackage[european,siunitx]{circuitikz} % schémas électriques
\pgfplotsset{compat=1.18}
\usepgfplotslibrary{fillbetween,groupplots} % aires entre courbes (calcul intégral)
\usepackage{enumitem}
\usepackage{fancyhdr}
% [most] charge toutes les bibliothèques tcolorbox (listings, minted, raster,
% documentation...) et gonfle inutilement la mémoire TeX sur un document long
% (~300 boîtes) : on ne charge que ce qui est réellement utilisé.
\usepackage[skins,breakable]{tcolorbox}
\usepackage{graphicx}
\usepackage{colortbl}
\usepackage{booktabs}
\usepackage{tabularx}
\usepackage{multirow}
\usepackage{diagbox} % case d'en-tête diagonale des tableaux à double entrée
% Colonnes extensibles centrée / gauche / droite (C, L, R), souvent utilisées par les modèles
\newcolumntype{C}{>{\centering\arraybackslash}X}
\newcolumntype{L}{>{\raggedright\arraybackslash}X}
\newcolumntype{R}{>{\raggedleft\arraybackslash}X}
\usepackage{array}
\usepackage{multicol}
\usepackage{siunitx}
% parse-numbers=false : accepte aussi \SI{2,5 \times 10^{-3}}{...}
\sisetup{locale=FR,output-decimal-marker={,},group-minimum-digits=4,parse-numbers=false}
\DeclareSIUnit\ohm{\ensuremath{\symup{\Omega}}} % U+2126 absent de la police
\DeclareSIUnit\division{div} % réglages d'oscilloscope (ms/div, V/div)
\DeclareSIUnit\tr{tr} % tours (vitesse de rotation, ex. \SI{1500}{\tr\per\minute})
\DeclareSIUnit\molar{M} % concentration molaire (ex. \SI{3}{\molar}, \SI{1}{\micro\molar})
\DeclareSIUnit\an{an} % année (le modèle confond parfois avec \year, un compteur TeX) — ex. \SI{5}{\kilo\meter\per\an}
\DeclareSIUnit\FCFA{FCFA} % franc CFA (ex. \SI{500}{\FCFA\per\kilo\gram})
\DeclareSIUnit\degreeBrix{\textdegree Brix} % teneur en sucre d'un sirop/jus (réfractométrie)
\DeclareSIUnit\month{mois} % le modèle utilise parfois \month, un compteur TeX, comme unité de durée
\DeclareSIUnit\microliter{\micro\liter} % le modèle écrit parfois \microliter en un seul token
\DeclareSIUnit\million{millions} % dénombrement (ex. \SI{15}{\million\per\mL} spermatozoïdes)
\usepackage{qrcode}
% \degree : commande fréquemment utilisée par les modèles pour un angle ou une
% température en dehors de \SI{}{} (non fournie par les paquets chargés).
\providecommand{\degree}{^{\circ}}
% \qed (fin de démonstration), utilisé par les modèles sans amsthm
\providecommand{\qed}{\hfill$\square$}
% \barre : double barre des valeurs interdites dans un tableau de variations (tkz-tab)
\providecommand{\barre}{\ensuremath{\|}}
% environnement proof (amsthm non chargé) utilisé par les modèles
\ifdefined\proof\else\newenvironment{proof}[1][Démonstration]{\par\noindent\textit{#1.}\ }{\qed\par}\fi
\usepackage{lastpage}
\usepackage{hyperref}
\hypersetup{hidelinks}

% ============================================================
% Palette « Pop » OnBuch+ — mêmes hex que la classe P (plus_ui.dart)
% ============================================================
\definecolor{poporange}{HTML}{F59321}
\definecolor{popdark}{HTML}{E07A0C}
\definecolor{popcream}{HTML}{FFF6E8}
\definecolor{popink}{HTML}{1C1714}
\definecolor{poppurple}{HTML}{7A5AE0}
\definecolor{popgreen}{HTML}{2E9E6A}
\definecolor{poppink}{HTML}{E0457B}
\definecolor{popblue}{HTML}{2F7DE1}
\definecolor{popgold}{HTML}{C98A12}
\definecolor{popmuted}{HTML}{8A7F76}
\definecolor{popline}{HTML}{EADFD2}
\definecolor{popgray}{HTML}{8A7F76}
\colorlet{popgrey}{popgray}
% Alias tolérants (le modèle nomme parfois les couleurs différemment)
\colorlet{popwhite}{white}
\colorlet{popblack}{popink}
\colorlet{popred}{poppink}
\colorlet{popyellow}{popgold}
% Teintes claires (fonds de tableaux/encadrés) — le modèle nomme parfois
% les couleurs avec un suffixe L (ex. popblueL) sans les avoir définies.
\colorlet{poporangeL}{poporange!20}
\colorlet{popdarkL}{popdark!20}
\colorlet{poppurpleL}{poppurple!20}
\colorlet{popgreenL}{popgreen!20}
\colorlet{poppinkL}{poppink!20}
\colorlet{popblueL}{popblue!20}
\colorlet{popgoldL}{popgold!20}
\colorlet{popredL}{popred!20}
\colorlet{popgrayL}{popgray!20}
% Teintes claires pour fonds de tableaux / zones
\colorlet{popblueL}{popblue!10!white}
\colorlet{poporangeL}{poporange!14!white}
\colorlet{popgreenL}{popgreen!12!white}
\colorlet{poppurpleL}{poppurple!12!white}
\colorlet{poppinkL}{poppink!12!white}
\colorlet{popgoldL}{popgold!14!white}

\pagecolor{popcream}

% ============================================================
% Typographie
% ============================================================
\defaultfontfeatures{Ligatures=TeX}
\setmainfont{PlusJakartaSans-Medium.ttf}[Path=../../../fonts/,BoldFont=PlusJakartaSans-Bold.ttf]
% Maths en sans-serif (Fira Math) avec chiffres et lettres droites de Plus
% Jakarta, pour que les formules se fondent dans le texte.
\usepackage[math-style=ISO,bold-style=ISO]{unicode-math}
\setmathfont{FiraMath-Regular.otf}[Path=../../../fonts/]
\setmathfont{PlusJakartaSans-Medium.ttf}[Path=../../../fonts/,range={up/num,up/{Latin,latin},\mathup}]
\usepackage{icomma} % virgule décimale sans espace en mode math
% Caractères mathématiques Unicode absents des polices de texte (Plus Jakarta,
% Archivo Black) : rendus par leur équivalent mathématique, en texte comme en
% formule — sinon ils s'affichent en carré vide (ex. « Arithmétique dans ℤ »).
\usepackage{newunicodechar}
\newunicodechar{ℝ}{\ensuremath{\mathbb{R}}}
\newunicodechar{ℤ}{\ensuremath{\mathbb{Z}}}
\newunicodechar{ℂ}{\ensuremath{\mathbb{C}}}
\newunicodechar{ℕ}{\ensuremath{\mathbb{N}}}
\newunicodechar{ℚ}{\ensuremath{\mathbb{Q}}}
\newunicodechar{𝔻}{\ensuremath{\mathbb{D}}}
\newunicodechar{α}{\ensuremath{\alpha}}
\newunicodechar{β}{\ensuremath{\beta}}
\newunicodechar{γ}{\ensuremath{\gamma}}
\newunicodechar{δ}{\ensuremath{\delta}}
\newunicodechar{ε}{\ensuremath{\varepsilon}}
\newunicodechar{θ}{\ensuremath{\theta}}
\newunicodechar{λ}{\ensuremath{\lambda}}
\newunicodechar{μ}{\ensuremath{\mu}}
\newunicodechar{σ}{\ensuremath{\sigma}}
\newunicodechar{τ}{\ensuremath{\tau}}
\newunicodechar{φ}{\ensuremath{\varphi}}
\newunicodechar{ψ}{\ensuremath{\psi}}
\newunicodechar{ω}{\ensuremath{\omega}}
\newunicodechar{Σ}{\ensuremath{\Sigma}}
% majuscules produites par \MakeUppercase dans les titres (α-aminés -> Α)
\newunicodechar{Α}{\ensuremath{\alpha}}
\newunicodechar{Β}{\ensuremath{\beta}}
\newunicodechar{Γ}{\ensuremath{\Gamma}}
\newunicodechar{Δ}{\ensuremath{\Delta}}
\newunicodechar{Ε}{\ensuremath{\varepsilon}}
\newunicodechar{Θ}{\ensuremath{\Theta}}
\newunicodechar{Λ}{\ensuremath{\Lambda}}
\newunicodechar{Μ}{\ensuremath{\mu}}
\newunicodechar{Τ}{\ensuremath{\tau}}
\newunicodechar{Φ}{\ensuremath{\Phi}}
\newunicodechar{Ψ}{\ensuremath{\Psi}}
\newunicodechar{Ω}{\ensuremath{\Omega}}
\newunicodechar{↦}{\ensuremath{\mapsto}}
\newunicodechar{∈}{\ensuremath{\in}}
\newunicodechar{∉}{\ensuremath{\notin}}
\newunicodechar{∧}{\ensuremath{\wedge}}
\newunicodechar{⇔}{\ensuremath{\Leftrightarrow}}
\newunicodechar{⟺}{\ensuremath{\Longleftrightarrow}}
\newunicodechar{⇒}{\ensuremath{\Rightarrow}}
\newunicodechar{⟹}{\ensuremath{\Longrightarrow}}
\newunicodechar{∘}{\ensuremath{\circ}}
\newunicodechar{∪}{\ensuremath{\cup}}
\newunicodechar{∩}{\ensuremath{\cap}}
\newunicodechar{≡}{\ensuremath{\equiv}}
\newunicodechar{⊂}{\ensuremath{\subset}}
\newunicodechar{⊥}{\ensuremath{\perp}}
\newunicodechar{∃}{\ensuremath{\exists}}
\newunicodechar{∀}{\ensuremath{\forall}}
\newunicodechar{∛}{\ensuremath{\sqrt[3]{\,}}}
\newunicodechar{∜}{\ensuremath{\sqrt[4]{\,}}}
\newunicodechar{⃗}{\ensuremath{{}^{\to}}}
\newunicodechar{ⁿ}{\ifmmode^{n}\else\textsuperscript{\ensuremath{n}}\fi}
\newunicodechar{ˣ}{\ifmmode^{x}\else\textsuperscript{\ensuremath{x}}\fi}
\newunicodechar{ᵇ}{\ifmmode^{b}\else\textsuperscript{\ensuremath{b}}\fi}
\newunicodechar{ʳ}{\ifmmode^{r}\else\textsuperscript{\ensuremath{r}}\fi}
\newunicodechar{ᵘ}{\ifmmode^{u}\else\textsuperscript{\ensuremath{u}}\fi}
\newunicodechar{ʸ}{\ifmmode^{y}\else\textsuperscript{\ensuremath{y}}\fi}
\newunicodechar{ᵃ}{\ifmmode^{a}\else\textsuperscript{\ensuremath{a}}\fi}
\newunicodechar{ᵉ}{\ifmmode^{e}\else\textsuperscript{\ensuremath{e}}\fi}
\newunicodechar{ᵏ}{\ifmmode^{k}\else\textsuperscript{\ensuremath{k}}\fi}
\newunicodechar{ᵖ}{\ifmmode^{p}\else\textsuperscript{\ensuremath{p}}\fi}
\newunicodechar{ᴺ}{\ifmmode^{N}\else\textsuperscript{\ensuremath{N}}\fi}
\newunicodechar{ᶜ}{\ifmmode^{c}\else\textsuperscript{\ensuremath{c}}\fi}
\newunicodechar{ʰ}{\ifmmode^{h}\else\textsuperscript{\ensuremath{h}}\fi}
\newunicodechar{ᵠ}{\ifmmode^{q}\else\textsuperscript{\ensuremath{q}}\fi}
\newunicodechar{ₙ}{\ifmmode_{n}\else\textsubscript{\ensuremath{n}}\fi}
\newunicodechar{ₐ}{\ifmmode_{a}\else\textsubscript{\ensuremath{a}}\fi}
\newunicodechar{ᵢ}{\ifmmode_{i}\else\textsubscript{\ensuremath{i}}\fi}
\newunicodechar{ᵣ}{\ifmmode_{r}\else\textsubscript{\ensuremath{r}}\fi}
\newunicodechar{ₖ}{\ifmmode_{k}\else\textsubscript{\ensuremath{k}}\fi}
\newunicodechar{ᵦ}{\ifmmode_{\beta}\else\textsubscript{\ensuremath{\beta}}\fi}
\newunicodechar{ₓ}{\ifmmode_{x}\else\textsubscript{\ensuremath{x}}\fi}
\newfontfamily\popdisplay{ArchivoBlack-Regular.ttf}[Path=../../../fonts/]
\newfontfamily\poplogo{SpaceGrotesk-Bold.ttf}[Path=../../../fonts/]
\newfontfamily\popsemi{PlusJakartaSans-SemiBold.ttf}[Path=../../../fonts/]
\newfontfamily\popxbold{PlusJakartaSans-ExtraBold.ttf}[Path=../../../fonts/]
\newfontfamily\popxboldls{PlusJakartaSans-ExtraBold.ttf}[Path=../../../fonts/,LetterSpace=3.0]

\setlist[enumerate]{leftmargin=1.7em,itemsep=5pt,topsep=4pt}
\setlist[itemize]{leftmargin=1.7em,itemsep=4pt,topsep=4pt,label={\color{poporange}\textbullet}}
\setlength{\parindent}{0pt}
\setlength{\parskip}{5pt}
\renewcommand{\arraystretch}{1.25}

% pgfplots : style Pop par défaut
\pgfplotsset{
  every axis/.append style={axis line style={popink,line width=1pt},tick style={popink},
    grid style={popline,line width=0.5pt},label style={font=\small},tick label style={font=\footnotesize}},
  popcurve/.style={poporange,line width=1.8pt,no markers,smooth},
}
% Même style disponible dans un \draw TikZ simple (hors pgfplots)
\tikzset{popcurve/.style={poporange,line width=1.8pt}}
\tikzset{popgrid/.style={popline,very thin}}
\colorlet{popcurve}{poporange} % le nom du style est parfois utilisé comme couleur

% ============================================================
% Pill / squiggle
% ============================================================
\newcommand{\poppill}[3]{%
  \tikz[baseline=-0.55ex]\node[draw=popink,line width=1.2pt,rounded corners=2.6pt,%
    fill=#2,inner xsep=6.5pt,inner ysep=3pt]{%
    \popxboldls\fontsize{7}{7}\selectfont\color{#3}\MakeUppercase{#1}};%
}
\newcommand{\popsquiggle}[1]{%
  \tikz[baseline=-0.4ex]{
    \draw[#1,line width=2.6pt,line cap=round]
      plot [smooth] coordinates {(0,0.09) (0.34,0.01) (0.68,0.21) (1.02,0.01) (1.36,0.10)};
  }%
}
% Mot-clé surligné (marqueur orange sous le texte)
\newcommand{\cle}[1]{{\popxbold\color{popdark}#1}}

% ============================================================
% En-tête / pied
% ============================================================
\pagestyle{fancy}
\fancyhf{}
\renewcommand{\headrulewidth}{0pt}
\renewcommand{\footrulewidth}{0pt}
\lhead{\raisebox{-0.15ex}{\poplogo\fontsize{13}{13}\selectfont\color{popink}OnBuch\color{poporange}+}}
\rhead{\poppill{\DOCMATIERE\ \textbullet\ \DOCNIVEAU\ \textbullet\ Leçon \DOCLECON}{white}{popink}}
\lfoot{\popxbold\fontsize{7.6}{7.6}\selectfont\color{popmuted}\MakeUppercase{Cours signé OnBuch+}\ \textbullet\ {\popsemi\DOCTITRE}}
\rfoot{\tikz[baseline=-0.6ex]\node[rounded corners=8.5pt,fill=popink,minimum height=17pt,minimum width=17pt,inner xsep=5pt,inner sep=0]{\popxbold\fontsize{8}{8}\selectfont\color{popcream}\thepage\,/\,\pageref*{LastPage}};}

% ============================================================
% Titres
% ============================================================
\newcommand{\chaptitre}[2]{%
  \begin{tcolorbox}[enhanced,colback=poporange,colframe=popink,boxrule=2.6pt,arc=11pt,
      drop shadow=popink,left=15pt,right=15pt,top=12pt,bottom=11pt,before skip=3pt,after skip=18pt]
    \poppill{#2}{popink}{popcream}\par\vspace{9pt}
    {\raggedright\hyphenpenalty=10000\exhyphenpenalty=10000\popdisplay\fontsize{22}{24}\selectfont\color{popcream}\MakeUppercase{#1}\par}
  \end{tcolorbox}
}
% Section numérotée : pastille orange avec le numéro + titre Archivo
\newcounter{popsec}
\newcommand{\coursec}[1]{%
  \stepcounter{popsec}\setcounter{popsub}{0}%
  \par\vspace{16pt}\noindent
  \begin{minipage}{\linewidth}
  \tikz[baseline=-0.7ex]\node[draw=popink,line width=1.6pt,fill=poporange,rounded corners=4pt,minimum size=22pt,inner sep=0]{\popdisplay\fontsize{12}{12}\selectfont\color{popcream}\thepopsec};%
  \hspace{8pt}{\popdisplay\fontsize{15}{17}\selectfont\color{popink}\MakeUppercase{#1}}\par
  \vspace{4pt}\noindent{\color{poporange}\rule{36pt}{3.6pt}}
  \end{minipage}\par\nopagebreak\vspace{8pt}
}
\newcounter{popsub}
\newcommand{\courssub}[1]{%
  \stepcounter{popsub}%
  \par\vspace{9pt}\noindent{\popxbold\fontsize{11.5}{13}\selectfont\color{popink}{\color{poporange}\thepopsec.\thepopsub}\hspace{6pt}#1}\par\nopagebreak\vspace{4pt}
}

% ============================================================
% Boîtes pédagogiques — même recette : fond blanc, bordure épaisse colorée,
% ombre dure encre, badge plein. Titre optionnel [#1].
% ============================================================
\tcbset{popbase/.style={breakable,enhanced,colback=white,boxrule=2.3pt,arc=9pt,
  shadow={3pt}{-3pt}{0pt}{popink},left=14pt,right=14pt,top=11pt,bottom=11pt,before skip=11pt,after skip=15pt,
  fontupper=\popsemi\fontsize{10.6}{14.5}\selectfont\color{popink},
  fontlower=\popsemi\fontsize{10.6}{14.5}\selectfont\color{popink}}}
% Coche dessinée (le glyphe ✓ n'existe pas dans les polices du document)
\newcommand{\popcheck}[1]{\tikz[baseline=-0.5ex]\draw[#1,line width=1.6pt,line cap=round,line join=round](-0.13,0.0)--(-0.04,-0.09)--(0.13,0.1);}
\providecommand{\coche}{\popcheck{popgreen}}
\providecommand{\resultat}[1]{\par\smallskip\noindent\fcolorbox{popgreen}{popgreenL}{\parbox{\dimexpr\linewidth-2\fboxsep-2\fboxrule\relax}{\textbf{Résultat.} #1}}\par\smallskip}
\newcommand{\popboxhead}[3]{\poppill{#1}{#2}{white}\ifx\relax#3\relax\else\hspace{6pt}{\popxbold\fontsize{10.6}{12}\selectfont\color{popink}#3}\fi\par\vspace{7pt}}

\newtcolorbox{aretenir}[1][]{popbase,colframe=popgreen,before upper={\popboxhead{À retenir}{popgreen}{#1}}}
\newtcolorbox{exemplebox}[1][]{popbase,colframe=poporange,before upper={\popboxhead{Exemple}{poporange}{#1}}}
\newtcolorbox{definition}[1][]{popbase,colframe=popblue,before upper={\popboxhead{Définition}{popblue}{#1}}}
\newtcolorbox{propriete}[1][]{popbase,colframe=poppurple,before upper={\popboxhead{Propriété}{poppurple}{#1}}}
\newtcolorbox{methode}[1][]{popbase,colframe=popgold,before upper={\popboxhead{Méthode}{popgold}{#1}}}
\newtcolorbox{attention}[1][]{popbase,colframe=poppink,before upper={\popboxhead{Attention}{poppink}{#1}}}
\newtcolorbox{experience}[1][]{popbase,colframe=popblue,colback=popblueL,before upper={\popboxhead{Expérience}{popblue}{#1}}}
% « Le savais-tu ? » : carte violette pleine (contexte, culture, Cameroun)
\newtcolorbox{savaistu}[1][]{popbase,colback=poppurple,colframe=popink,
  fontupper=\popsemi\fontsize{10.4}{14.2}\selectfont\color{white},
  before upper={\poppill{Le savais-tu\,?}{white}{poppurple}\ifx\relax#1\relax\else\hspace{6pt}{\popxbold\color{white}#1}\fi\par\vspace{7pt}}}
% Exercice résolu : énoncé en haut, \tcblower puis la solution sur fond crème
\newcounter{popexo}\newcounter{popexr}
\newtcolorbox{exoresolu}[1][]{popbase,colframe=poporange,colbacklower=popcream,
  segmentation style={popink,line width=1.2pt,dashed},
  before upper={\stepcounter{popexr}\popboxhead{Exercice résolu \thepopexr}{poporange}{#1}},
  before lower={\poppill{Solution}{popink}{popcream}\par\vspace{6pt}}}
% Exercice (non résolu en place) — la correction est dans la partie Corrigés
\newtcolorbox{exercice}[1][]{popbase,colframe=popink,
  before upper={\stepcounter{popexo}\popboxhead{Exercice \thepopexo}{popink}{#1}}}
% Corrigé
\newtcolorbox{corrige}[1][]{popbase,colframe=popgreen,colback=popgreenL,
  before upper={\popboxhead{Corrigé}{popgreen}{#1}}}
% Objectifs / prérequis / bilan
\newtcolorbox{objectifs}{popbase,colframe=popink,colback=poporangeL,
  before upper={\popboxhead{Objectifs de la leçon}{popink}{}}}
\newtcolorbox{prerequis}{popbase,colframe=popink,colback=popblueL,
  before upper={\popboxhead{Prérequis}{popblue}{}}}
\newtcolorbox{bilan}{popbase,colframe=popink,colback=white,boxrule=2.8pt,
  before upper={\popboxhead{Fiche bilan}{poporange}{L'essentiel à connaître pour l'examen}}}
% Figure encadrée avec légende
\newtcolorbox{popfigure}[1][]{enhanced,breakable=false,colback=white,colframe=popline,boxrule=1.4pt,arc=8pt,
  left=10pt,right=10pt,top=10pt,bottom=8pt,before skip=10pt,after skip=12pt,halign=center,
  lower separated=false,halign lower=center,
  fontlower=\popsemi\fontsize{9}{11.5}\selectfont\color{popmuted},
  #1}
\newcommand{\legende}[1]{\par\vspace{6pt}{\popsemi\fontsize{9}{11.5}\selectfont\color{popmuted}#1}}

% ============================================================
% Signature OnBuch+
% ============================================================
\newcommand{\poplogoinline}[1]{{\poplogo\fontsize{#1}{#1}\selectfont\color{popink}OnBuch\color{poporange}+}}
\newcommand{\signatureonbuch}{%
  \begin{tcolorbox}[enhanced,colback=popink,colframe=popink,boxrule=2.2pt,arc=10pt,
      drop shadow=poporange,left=14pt,right=14pt,top=10pt,bottom=10pt,before skip=0pt,after skip=0pt]
    \tikz[baseline=-0.7ex]\node[circle,fill=popgreen,draw=popcream,line width=1.3pt,minimum size=22pt,inner sep=0]{\popcheck{white}};%
    \hspace{9pt}\begin{minipage}[c]{0.62\linewidth}
      {\popxbold\fontsize{12}{13}\selectfont\color{popcream}Signé\ \textbullet\ {\poplogo OnBuch\color{poporange}+}}\par\vspace{2pt}
      {\raggedright\popsemi\fontsize{8}{10.5}\selectfont\color{popline}\MakeUppercase{Équipe pédagogique OnBuch+}\\\MakeUppercase{Conforme au programme officiel MINESEC}\par}
    \end{minipage}\hfill
    {\poplogo\fontsize{22}{22}\selectfont\color{popcream}OnBuch\color{poporange}+}
  \end{tcolorbox}%
}

% ============================================================
% Page de garde
% #1 titre · #2 matière · #3 niveau · #4 module · #5 n° leçon · #6 durée ·
% #7 description / objectifs (texte ou itemize)
% ============================================================
\newcommand{\pagegarde}[7]{%
  \thispagestyle{empty}
  \newgeometry{a4paper,margin=1.7cm,top=1.6cm,bottom=1.4cm}%
  \begin{tikzpicture}[remember picture,overlay]
    \fill[poporange!18] ([xshift=-3.2cm,yshift=-3.6cm]current page.north east) circle (4.6cm);
    \fill[poppurple!14] ([xshift=2.4cm,yshift=7.6cm]current page.south west) circle (3.8cm);
  \end{tikzpicture}%
  {\poplogo\fontsize{30}{30}\selectfont\color{popink}OnBuch\color{poporange}+}\hfill
  \raisebox{0.6ex}{\poppill{Cours signé \textbullet\ Premium}{popink}{popcream}}\par
  \vspace{1pt}\popsquiggle{poppurple}\par
  \vspace{8pt}
  \begin{tcolorbox}[enhanced,colback=poporange,colframe=popink,boxrule=2.8pt,arc=13pt,
      drop shadow=popink,left=18pt,right=18pt,top=12pt,bottom=13pt,after skip=10pt]
    \poppill{#2 \textbullet\ #3}{popink}{popcream}\hspace{5pt}\poppill{#4}{white}{popink}\par\vspace{8pt}
    {\popdisplay\fontsize{14}{14}\selectfont\color{popink}LEÇON #5}\par\vspace{4pt}
    {\raggedright\hyphenpenalty=10000\exhyphenpenalty=10000\popdisplay\fontsize{25}{27}\selectfont\color{popcream}\MakeUppercase{#1}\par}
    \vspace{8pt}
    {\popxbold\fontsize{9.5}{9.5}\selectfont\color{popink}\MakeUppercase{Durée conseillée : #6}}
  \end{tcolorbox}
  \begin{tcolorbox}[enhanced,colback=white,colframe=popink,boxrule=2.2pt,arc=10pt,
      drop shadow=popink,left=16pt,right=16pt,top=10pt,bottom=10pt,after skip=8pt]
    \poppill{Dans cette leçon}{poppurple}{white}\par\vspace{5pt}
    \popsemi\fontsize{9.3}{12.6}\selectfont\color{popink}#7
  \end{tcolorbox}
  \noindent\begin{minipage}[t]{0.487\linewidth}
    \begin{tcolorbox}[enhanced,colback=popgreen,colframe=popink,boxrule=2.2pt,arc=10pt,
        drop shadow=popink,left=11pt,right=11pt,top=10pt,bottom=10pt,height=3.1cm,valign=center]
      \tcbox[colback=white,colframe=popink,boxrule=1.3pt,arc=5pt,boxsep=0pt,nobeforeafter,
        left=3pt,right=3pt,top=3pt,bottom=3pt,valign=center]{\qrcode[height=1.9cm]{https://play.google.com/store/apps/details?id=com.luvvix.onbuch&hl=fr}}%
      \hspace{8pt}\begin{minipage}[c]{0.5\linewidth}
        {\popdisplay\fontsize{11}{12.5}\selectfont\color{white}\MakeUppercase{Télécharge OnBuch}}\par\vspace{4pt}
        {\raggedright\popsemi\fontsize{8.4}{11}\selectfont\color{white}Gratuit \textbullet\ Google Play\\Tuteur IA, annales, concours, résultats\par}
      \end{minipage}
    \end{tcolorbox}
  \end{minipage}\hfill
  \begin{minipage}[t]{0.487\linewidth}
    \begin{tcolorbox}[enhanced,colback=poppurple,colframe=popink,boxrule=2.2pt,arc=10pt,
        drop shadow=popink,left=11pt,right=11pt,top=10pt,bottom=10pt,height=3.1cm,valign=center]
      \tikz[baseline=-0.7ex]\node[circle,fill=white,draw=popink,line width=1.3pt,minimum size=1.6cm,inner sep=0]{\popdisplay\fontsize{24}{24}\selectfont\color{poppurple}+};%
      \hspace{8pt}\begin{minipage}[c]{0.6\linewidth}
        {\popdisplay\fontsize{11}{12.5}\selectfont\color{white}\MakeUppercase{Passe à OnBuch+}}\par\vspace{4pt}
        {\raggedright\popsemi\fontsize{8.4}{11}\selectfont\color{white}Tous les cours signés, Tuteur IA illimité, fiches PDF — onglet OnBuch+ de l'app\par}
      \end{minipage}
    \end{tcolorbox}
  \end{minipage}
  \vspace{8pt}
  \popcontenu
  \vfill
  \signatureonbuch
  \restoregeometry
  \newpage
}

% Rangée « Ce cours contient » (page de garde)
\newcommand{\poptile}[3]{% #1 couleur #2 gros texte #3 légende
  \begin{tcolorbox}[enhanced,colback=#1,colframe=popink,boxrule=2pt,arc=9pt,shadow={2.5pt}{-2.5pt}{0pt}{popink},
      left=6pt,right=6pt,top=9pt,bottom=9pt,halign=center,valign=center,height=1.75cm,width=\linewidth,nobeforeafter]
    {\popdisplay\fontsize{12.5}{13}\selectfont\color{white}\MakeUppercase{#2}}\par\vspace{3pt}
    {\centering\popsemi\fontsize{7.8}{9.5}\selectfont\color{white}#3\par}
  \end{tcolorbox}}
\newcommand{\popcontenu}{%
  \par\noindent{\popxboldls\fontsize{7.5}{7.5}\selectfont\color{popmuted}CE COURS SIGNÉ CONTIENT}\par\vspace{7pt}
  \noindent\begin{minipage}[t]{0.235\linewidth}\poptile{popblue}{Cours}{complet, illustré, pas à pas}\end{minipage}\hfill
  \begin{minipage}[t]{0.235\linewidth}\poptile{poporange}{Méthodes}{et exercices résolus}\end{minipage}\hfill
  \begin{minipage}[t]{0.235\linewidth}\poptile{poppink}{Exercices}{type examen + corrigés}\end{minipage}\hfill
  \begin{minipage}[t]{0.235\linewidth}\poptile{popgreen}{Fiche bilan}{l'essentiel à retenir}\end{minipage}\par}

% Sommaire
\newtcolorbox{sommaire}{enhanced,colback=white,colframe=popink,boxrule=2.2pt,arc=10pt,
  drop shadow=popink,left=18pt,right=18pt,top=13pt,bottom=13pt,before skip=6pt,after skip=20pt,
  before upper={\poppill{Sommaire}{poppurple}{white}\par\vspace{9pt}\popsemi\fontsize{10.6}{14.5}\selectfont\color{popink}}}

% Signature de fin de cours — poussée tout en bas de la dernière page
\newcommand{\findecours}{%
  \par\vfill\nopagebreak
  \begin{center}\popsquiggle{poporange}\end{center}\vspace{4pt}
  \signatureonbuch
}
\AtBeginDocument{\def\llbracket{\mathopen{[\mkern-3.5mu[}}\def\rrbracket{\mathclose{]\mkern-3.5mu]}}} % intervalles d'entiers (glyphe ⟦ absent de la police)
% Glyphes absents de Fira Math : redéfinis à partir de glyphes disponibles
\AtBeginDocument{%
  \def\setminus{\mathbin{\backslash}}\let\smallsetminus\setminus
  \def\vdots{\mathord{\vbox{\baselineskip3.2pt\lineskiplimit0pt\kern4pt\hbox{.}\hbox{.}\hbox{.}}}}%
  \def\ddots{\mathinner{\mkern1mu\raise6.5pt\hbox{.}\mkern2mu\raise3.6pt\hbox{.}\mkern2mu\raise0.7pt\hbox{.}\mkern1mu}}%
  \def\triangle{\mathord{\scalebox{0.9}{$\symup{\Delta}$}}}%
  \def\nabla{\mathord{\raisebox{\depth}{\scalebox{1}[-1]{$\symup{\Delta}$}}}}%
  \def\checkmark{\popcheck{popdark}}%
  \def\star{\mathbin{\ast}}\def\bigstar{\mathord{\ast}}%
  \def\diamond{\mathbin{\scalebox{0.6}{$\Diamond$}}}%
  \def\bigcup{\mathop{\mathchoice{\raisebox{-0.3em}{\scalebox{1.7}{$\cup$}}}{\scalebox{1.25}{$\cup$}}{\cup}{\cup}}\displaylimits}%
  \def\bigcap{\mathop{\mathchoice{\raisebox{-0.3em}{\scalebox{1.7}{$\cap$}}}{\scalebox{1.25}{$\cap$}}{\cap}{\cap}}\displaylimits}%
  \def\lesseqqgtr{\mathrel{\lessgtr}}\def\bigoplus{\mathop{\oplus}\displaylimits}%
}
\providecommand{\ssi}{si et seulement si} % abréviation fréquente
\AtBeginDocument{\providecommand{\wideparen}{\overparen}} % arc de cercle

%% FIGFIX-BEGIN v4 — correctifs globaux des figures (docs/onbuch-generation/tools/figfix)
\usepackage{adjustbox}\usepackage{etoolbox}
% toute figure TikZ/pgfplots plus large que la ligne est réduite à la largeur disponible
\BeforeBeginEnvironment{tikzpicture}{\begin{adjustbox}{max width=\linewidth}}
\AfterEndEnvironment{tikzpicture}{\end{adjustbox}}
% légendes pgfplots lisibles par-dessus les courbes
\pgfplotsset{every axis/.append style={legend style={fill=white,fill opacity=0.92,text opacity=1,draw=black!15,inner sep=3pt}}}
% étiquettes d'arêtes (syntaxe edge["..."]) avec fond blanc
\tikzset{every edge quotes/.append style={fill=white,inner sep=1.2pt}}
%% FIGFIX-END
\begin{document}

\pagegarde{Aide au système immunitaire}{SVTEEHB}{3ème}{SVTEEHB – classe de 3ème}{N08}{3 séances (fiche de progression harmonisée)}{Cette leçon présente les différentes façons d'aider le système immunitaire à défendre l'organisme : l'antibiothérapie, la sérothérapie, la vaccinothérapie et la séro-vaccinothérapie. À partir d'expériences historiques (Pasteur, Fleming, Behring) et de documents à interpréter, l'élève apprend à distinguer ces méthodes et à les appliquer à des situations concrètes de santé au Cameroun.
\vspace{4pt}
\begin{multicols}{2}\small
\begin{itemize}
\item À la fin de cette leçon, je sais définir l'antibiothérapie et expliquer son principe d'action.
\item À la fin de cette leçon, je sais interpréter l'expérience de Fleming pour établir le mode d'action des antibiotiques.
\item À la fin de cette leçon, je sais définir la sérothérapie et expliquer son principe à partir de l'expérience de Behring.
\item À la fin de cette leçon, je sais définir la vaccinothérapie et expliquer son principe à partir des expériences de Pasteur.
\item À la fin de cette leçon, je sais définir la séro-vaccinothérapie et justifier son utilisation en cas d'urgence.
\item À la fin de cette leçon, je sais comparer les quatre méthodes d'aide à l'immunité (rapidité, durée, type d'immunité).
\end{itemize}\end{multicols}}

\begin{sommaire}
\begin{multicols}{2}
\begin{enumerate}
\item I. L'antibiothérapie
\item II. La sérothérapie
\item III. La vaccinothérapie
\item IV. La séro-vaccinothérapie
\item V. Comparaison et synthèse des méthodes d'aide à l'immunité
\item Activité d'intégration
\item Méthodes et exercices résolus
\item Exercices
\item Corrigés des exercices
\item Fiche bilan
\end{enumerate}
\end{multicols}
\end{sommaire}

\begin{objectifs}
\begin{itemize}
\item À la fin de cette leçon, je sais définir l'antibiothérapie et expliquer son principe d'action.
\item À la fin de cette leçon, je sais interpréter l'expérience de Fleming pour établir le mode d'action des antibiotiques.
\item À la fin de cette leçon, je sais définir la sérothérapie et expliquer son principe à partir de l'expérience de Behring.
\item À la fin de cette leçon, je sais définir la vaccinothérapie et expliquer son principe à partir des expériences de Pasteur.
\item À la fin de cette leçon, je sais définir la séro-vaccinothérapie et justifier son utilisation en cas d'urgence.
\item À la fin de cette leçon, je sais comparer les quatre méthodes d'aide à l'immunité (rapidité, durée, type d'immunité).
\item À la fin de cette leçon, je sais appliquer ces notions à des situations concrètes de la vie courante (morsure de serpent, tétanos, campagne de vaccination).
\item À la fin de cette leçon, je sais rédiger et justifier une démarche, une réponse ou une conclusion à partir de documents.
\end{itemize}
\end{objectifs}

\begin{prerequis}
\begin{itemize}
\item Les agents pathogènes : microbes, virus, bactéries (notions de 4ème).
\item Les défenses de l'organisme : barrières naturelles, phagocytose (leçons précédentes de 3ème).
\item La réaction immunitaire : antigènes, anticorps, lymphocytes (leçon sur les réactions immunitaires).
\item L'immunité naturelle et l'immunité acquise.
\item La lecture et l'interprétation de tableaux de résultats expérimentaux et de courbes.
\end{itemize}
\end{prerequis}

\newpage

\coursec{I. L'antibiothérapie}

\courssub{Situation d'accroche et problématique}

À l'hôpital de district de Bafoussam, dans la région de l'Ouest du Cameroun, le docteur Ngo Bassa reçoit ce matin deux enfants atteints de la même infection bactérienne : une angine à streptocoques. Au premier, elle prescrit un médicament à prendre pendant plusieurs jours : un antibiotique. Plus tard dans la journée, un troisième enfant arrive, blessé au pied par un clou rouillé : il reçoit une injection d'un sérum antitétanique. Le lendemain, toute l'école du quartier est convoquée pour une campagne de vaccination contre la fièvre jaune.

Trois situations, trois « aides » différentes données au système immunitaire : un médicament, un sérum, un vaccin. Pourquoi le médecin choisit-il tantôt l'un, tantôt l'autre ? Ces trois méthodes fonctionnent-elles de la même manière ? Agissent-elles avant ou après la maladie ? C'est ce que nous allons découvrir dans cette leçon, en commençant par la plus ancienne et la plus connue de ces aides : l'\cle{antibiothérapie}.

\begin{aretenir}[Problématique de la leçon]
Comment l'être humain peut-il aider son système immunitaire à lutter contre les microbes ? Nous étudierons successivement quatre méthodes : l'antibiothérapie, la sérothérapie, la vaccinothérapie et la séro-vaccinothérapie.
\end{aretenir}

\courssub{Rappel : les bactéries, des microbes responsables de maladies}

Au cours des leçons précédentes, nous avons appris que de nombreuses maladies sont causées par des \cle{microorganismes} : des êtres vivants invisibles à l'œil nu. Parmi eux, les \cle{bactéries} sont des êtres unicellulaires capables de se multiplier très rapidement dans l'organisme. Elles sont responsables de nombreuses maladies :

\begin{itemize}
\item l'\textbf{angine} (bactérie : le streptocoque), fréquente chez les élèves ;
\item la \textbf{tuberculose} (bactérie : le bacille de Koch), qui touche les poumons ;
\item le \textbf{choléra} (bactérie : le vibrion cholérique), qui provoque encore des épidémies dans certaines localités du Cameroun lorsque l'eau de boisson est souillée.
\end{itemize}

Face à ces bactéries, notre organisme possède des défenses naturelles : la peau, les phagocytes, les anticorps. Mais ces défenses ne suffisent pas toujours : certaines bactéries se multiplient si vite, ou sont si agressives, que le système immunitaire est débordé. Autrefois, une simple angine ou une plaie infectée pouvait tuer. Il fallait trouver une aide extérieure. Cette aide est venue d'une observation extraordinaire, faite par hasard dans un laboratoire de Londres.

\courssub{L'observation de Fleming (1928) : la naissance des antibiotiques}

\begin{experience}[L'observation de Fleming : une moisissure qui tue les bactéries]
\textbf{Contexte.} En 1928, le médecin et bactériologiste écossais Alexander \textbf{Fleming} cultive des bactéries, les \textbf{staphylocoques}, dans des boîtes de Pétri contenant un milieu nutritif (gélose).

\textbf{Protocole.} Avant de partir en vacances, Fleming laisse une boîte de Pétri ensemencée de staphylocoques sur sa paillasse, sans couvercle bien fermé. À son retour, il observe la boîte avant de la jeter.

\textbf{Observations.}
\begin{enumerate}
\item La surface de la gélose est couverte de colonies de staphylocoques (taches blanchâtres).
\item Une \textbf{moisissure} verte, le champignon microscopique \emph{Penicillium}, s'est développée au milieu de la boîte (contamination accidentelle).
\item Autour de la moisissure, sur une certaine distance, les colonies de staphylocoques ont \textbf{disparu} : on observe une \textbf{zone claire}, appelée \cle{zone d'inhibition}.
\end{enumerate}

\textbf{Hypothèse de Fleming.} Le \emph{Penicillium} libère dans la gélose une substance qui diffuse autour de lui et qui tue les bactéries ou empêche leur multiplication.

\textbf{Vérification.} Fleming cultive le \emph{Penicillium}, récupère le liquide de culture filtré (sans le champignon) et le dépose sur des cultures de bactéries : les bactéries meurent. L'hypothèse est confirmée.

\textbf{Conclusion.} Le \emph{Penicillium} produit une substance antibactérienne que Fleming nomme \textbf{pénicilline}. C'est le premier \cle{antibiotique} découvert.
\end{experience}

\begin{popfigure}
\begin{center}
\begin{tikzpicture}[scale=1.0]
% Boîte de Pétri
\draw[thick, popdark, fill=popcream] (0,0) ellipse (4.2 and 1.9);
\draw[thick, popdark] (-4.2,0) -- (-4.2,0.45) arc[start angle=180, end angle=0, x radius=4.2, y radius=1.9] -- (4.2,0);
\draw[thick, popdark] (0,0.45) ellipse (4.2 and 1.9);
% Zone d'inhibition (claire) autour du Penicillium
\fill[white, opacity=0.95] (0.6,0.35) ellipse (1.5 and 0.75);
\draw[dashed, poporange, thick] (0.6,0.35) ellipse (1.5 and 0.75);
% Colonies de staphylocoques (petits points)
\foreach \p in {(-3.4,0.3),(-3.0,-0.6),(-2.6,0.9),(-2.2,-0.2),(-1.8,0.5),(-3.6,-0.9),(-2.9,1.2),(-1.5,-0.9),(2.4,0.8),(3.0,-0.4),(3.5,0.5),(2.8,-1.0),(3.6,-0.8),(2.2,-0.3),(3.2,1.1),(-3.8,0.8),(2.0,1.3),(-2.4,-1.2),(3.8,0.1),(-1.2,1.3)}{
  \fill[popgold] \p circle (0.09);}
% Penicillium
\fill[popgreen] (0.6,0.35) ellipse (0.45 and 0.25);
\draw[popgreen!60!black, thick] (0.6,0.35) ellipse (0.45 and 0.25);
\node[popgreen!50!black, font=\small] at (0.6,0.35) {\emph{Penicillium}};
% Légendes fléchées
\draw[->, poporange, thick] (1.9,-0.1) -- (3.1,-1.6) node[below, font=\small, align=center, text=popdark]{zone claire d'inhibition\\ (bactéries détruites)};
\draw[->, popgold!70!black, thick] (-3.0,0.9) -- (-3.4,1.7) node[above, font=\small, text=popdark]{colonies de staphylocoques};
\draw[->, popgreen!60!black, thick] (0.6,0.6) -- (0.6,1.7) node[above, font=\small, text=popdark]{moisissure};
\end{tikzpicture}
\end{center}
\legende{Boîte de Pétri de Fleming (1928). La gélose est ensemencée de staphylocoques (points jaunes). Une moisissure, le \emph{Penicillium} (tache verte), s'est développée accidentellement. Autour d'elle, une zone claire (cerclée de pointillés orange) montre que les bactéries ont disparu : le \emph{Penicillium} libère une substance antibactérienne, la pénicilline.}
\end{popfigure}

\begin{savaistu}[Une découverte due au hasard... et à un bon observateur]
Alexander Fleming a découvert la pénicilline par hasard en 1928, grâce à une boîte de culture oubliée. Mais le hasard ne suffit pas : beaucoup de chercheurs auraient jeté cette boîte « ratée » sans y regarder de plus près. Fleming, lui, s'est posé la bonne question : \emph{pourquoi les bactéries ont-elles disparu autour de la moisissure ?} La pénicilline, produite industriellement à partir de 1941, a sauvé des millions de vies pendant la Seconde Guerre mondiale et depuis. Fleming a reçu le prix Nobel de médecine en 1945, avec Florey et Chain qui avaient permis la fabrication du médicament.
\end{savaistu}

\courssub{Définition et principe de l'antibiothérapie}

\begin{definition}[L'antibiothérapie]
L'\cle{antibiothérapie} est le \textbf{traitement d'une maladie microbienne à l'aide d'antibiotiques}. Un \cle{antibiotique} est une substance, produite par un microorganisme (champignon ou bactérie) ou fabriquée en laboratoire, qui \textbf{tue les bactéries} ou \textbf{bloque leur multiplication}.
\end{definition}

Le principe de l'antibiothérapie est simple à comprendre :

\begin{itemize}
\item l'antibiotique, absorbé par le malade (comprimés, sirop ou injection), se répand dans l'organisme et atteint le lieu de l'infection ;
\item il \textbf{détruit les bactéries} ou les \textbf{empêche de se multiplier} ;
\item le nombre de bactéries diminue fortement, ce qui permet au \textbf{système immunitaire} (phagocytes, anticorps) d'achever l'élimination de l'infection.
\end{itemize}

L'antibiotique ne remplace donc pas le système immunitaire : il l'\textbf{aide} en réduisant l'armée des bactéries à une taille que nos défenses peuvent vaincre. C'est bien une « aide au système immunitaire ».

\begin{exemplebox}[Exemples d'antibiotiques courants]
La \textbf{pénicilline} et ses dérivés (comme l'\textbf{amoxicilline}) sont utilisés contre l'angine, les infections des plaies, certaines pneumonies. D'autres antibiotiques existent, chacun étant actif contre certaines bactéries seulement : c'est le médecin qui choisit le bon antibiotique selon la maladie.
\end{exemplebox}

\courssub{Mise en évidence expérimentale : l'antibiogramme}

Toutes les bactéries ne sont pas sensibles à tous les antibiotiques. Comment le laboratoire choisit-il l'antibiotique efficace contre la bactérie d'un malade ? Grâce à un test appelé \cle{antibiogramme}.

\begin{experience}[Réalisation d'un antibiogramme]
\textbf{Matériel.} Boîte de Pétri contenant une gélose nutritive, culture de la bactérie prélevée chez le malade, petits disques de papier imprégnés chacun d'un antibiotique différent (A, B, C, D), étuve à \SI{37}{\celsius}.

\textbf{Protocole.}
\begin{enumerate}
\item On ensemence uniformément la gélose avec la bactérie du malade.
\item On dépose à la surface plusieurs disques, chacun imbibé d'un antibiotique différent.
\item On place la boîte à l'étuve pendant 24 à 48 heures.
\end{enumerate}

\textbf{Observations.} Les bactéries ont poussé partout sur la gélose (voile trouble), \textbf{sauf autour de certains disques}, où l'on observe des zones claires circulaires appelées \textbf{auréoles d'inhibition}. Ces auréoles ont des diamètres différents selon les antibiotiques ; autour de certains disques, il n'y a aucune auréole.

\textbf{Interprétation.}
\begin{itemize}
\item Une \textbf{auréole d'inhibition} autour d'un disque signifie que l'antibiotique a empêché la croissance des bactéries : la bactérie est \textbf{sensible} à cet antibiotique.
\item \textbf{Plus l'auréole est grande, plus l'antibiotique est efficace} contre cette bactérie.
\item L'\textbf{absence d'auréole} signifie que la bactérie est \textbf{résistante} à cet antibiotique : il ne faut pas l'utiliser pour soigner le malade.
\end{itemize}

\textbf{Conclusion.} L'antibiogramme permet de choisir, parmi plusieurs antibiotiques, celui qui sera le plus efficace pour traiter l'infection du malade.
\end{experience}

\begin{popfigure}
\begin{center}
\begin{tikzpicture}[scale=1.0]
% Boîte de Pétri
\draw[thick, popdark, fill=popcream] (0,0) circle (3.4);
\draw[thick, popdark] (0,0) circle (3.55);
% Voile bactérien (léger)
\fill[popgold!25, opacity=0.5] (0,0) circle (3.35);
% Disque A : grande auréole
\fill[white] (-1.7,1.2) circle (1.15);
\draw[dashed, poporange, thick] (-1.7,1.2) circle (1.15);
\fill[popblue] (-1.7,1.2) circle (0.22);
\node[white, font=\footnotesize\bfseries] at (-1.7,1.2) {A};
% Disque B : auréole moyenne
\fill[white] (1.8,1.0) circle (0.7);
\draw[dashed, poporange, thick] (1.8,1.0) circle (0.7);
\fill[popblue] (1.8,1.0) circle (0.22);
\node[white, font=\footnotesize\bfseries] at (1.8,1.0) {B};
% Disque C : petite auréole
\fill[white] (-1.4,-1.6) circle (0.42);
\draw[dashed, poporange, thick] (-1.4,-1.6) circle (0.42);
\fill[popblue] (-1.4,-1.6) circle (0.22);
\node[white, font=\footnotesize\bfseries] at (-1.4,-1.6) {C};
% Disque D : pas d'auréole
\fill[popblue] (1.7,-1.5) circle (0.22);
\node[white, font=\footnotesize\bfseries] at (1.7,-1.5) {D};
% Légendes
\draw[->, popdark, thick] (-2.85,1.2) -- (-4.1,2.3) node[above left, font=\small, align=center, text=popdark]{grande auréole :\\ bactérie très sensible};
\draw[->, popdark, thick] (1.7,-1.72) -- (2.6,-2.9) node[below, font=\small, align=center, text=popdark]{pas d'auréole :\\ bactérie résistante};
\draw[->, popdark, thick] (2.5,1.35) -- (3.9,2.4) node[above right, font=\small, align=center, text=popdark]{auréole moyenne :\\ sensibilité moyenne};
\node[font=\small, text=popdark] at (0,-3.95) {voile de culture bactérienne (fond jaunâtre)};
\end{tikzpicture}
\end{center}
\legende{Antibiogramme. Quatre disques (A, B, C, D) imprégnés d'antibiotiques différents sont déposés sur une culture de la bactérie du malade. Autour des disques A, B et C apparaissent des auréoles d'inhibition de tailles différentes : l'antibiotique A est le plus efficace. Autour du disque D, aucune auréole : la bactérie est résistante à l'antibiotique D.}
\end{popfigure}

\begin{methode}[Comment interpréter un antibiogramme]
\begin{enumerate}
\item Repérer chaque disque et l'antibiotique qu'il contient.
\item Observer s'il existe une \textbf{auréole d'inhibition} (zone claire) autour du disque.
\item \textbf{Comparer les diamètres} des auréoles : plus l'auréole est grande, plus l'antibiotique est efficace contre la bactérie.
\item Conclure : l'antibiotique à prescrire est celui qui donne la \textbf{plus grande auréole} ; un disque sans auréole correspond à un antibiotique \textbf{inefficace} (bactérie résistante).
\end{enumerate}
\end{methode}

\begin{exoresolu}[Interpréter un antibiogramme]
\textbf{Énoncé.} Un laboratoire de l'hôpital de Bafoussam réalise l'antibiogramme d'une bactérie prélevée sur un malade souffrant d'une infection urinaire. Les diamètres des auréoles d'inhibition mesurés sont : antibiotique A : \SI{28}{mm} ; antibiotique B : \SI{12}{mm} ; antibiotique C : \SI{0}{mm} (aucune auréole). Quel antibiotique le médecin doit-il prescrire ? Justifier la réponse.
\tcblower
\textbf{Solution détaillée.}
\begin{enumerate}
\item L'antibiotique C ne produit aucune auréole d'inhibition (\SI{0}{mm}) : la bactérie se multiplie normalement en sa présence. La bactérie est donc \textbf{résistante} à l'antibiotique C, qui ne doit pas être prescrit.
\item Les antibiotiques A et B produisent des auréoles : la bactérie leur est \textbf{sensible}.
\item On compare les diamètres : \SI{28}{mm} pour A contre \SI{12}{mm} pour B. Or, plus l'auréole est grande, plus l'antibiotique est efficace.
\end{enumerate}
\textbf{Conclusion.} Le médecin doit prescrire l'\textbf{antibiotique A}, car c'est celui qui donne la plus grande auréole d'inhibition : c'est le plus efficace contre la bactérie du malade.
\end{exoresolu}

\courssub{Les limites de l'antibiothérapie}

\textbf{1. Les antibiotiques sont inefficaces contre les virus.} Les antibiotiques n'agissent que sur les \textbf{bactéries}. Ils sont sans aucun effet sur les maladies causées par des \textbf{virus}, comme la grippe, la rougeole ou le SIDA. Prendre un antibiotique contre une grippe est donc inutile : le virus n'est pas une bactérie et n'a pas les structures sur lesquelles l'antibiotique agit.

\textbf{2. L'usage abusif crée des bactéries résistantes.} Dans une population de bactéries, il existe parfois quelques individus naturellement capables de résister à un antibiotique. Si l'on utilise mal l'antibiotique (dose trop faible, traitement arrêté trop tôt, antibiotique pris sans raison), les bactéries sensibles meurent mais les bactéries résistantes survivent et se multiplient. On obtient alors une souche \textbf{résistante} : c'est l'\cle{antibiorésistance}. Les infections causées par ces bactéries deviennent très difficiles à soigner.

\begin{attention}[Erreur fréquente à éviter]
Ne crois pas que les antibiotiques guérissent toutes les maladies ! Ils sont \textbf{sans effet sur les maladies virales} (grippe, rougeole, SIDA, rhume). De même, ce n'est pas parce qu'un antibiotique a soigné ton voisin qu'il soignera ta maladie : chaque infection exige un antibiotique adapté, choisi par le médecin, parfois grâce à un antibiogramme.
\end{attention}

\begin{exemplebox}[L'automédication au Cameroun : un danger réel]
Au Cameroun, on peut malheureusement acheter de l'amoxicilline dans la rue ou au marché, sans ordonnance. Beaucoup de personnes en prennent dès le premier mal de tête ou pour un simple rhume (une maladie virale !), puis arrêtent le traitement dès qu'elles se sentent mieux. Ces mauvaises habitudes favorisent l'\textbf{antibiorésistance}, devenue un grave problème de santé publique : certaines infections autrefois faciles à soigner ne répondent plus aux antibiotiques courants.
\end{exemplebox}

\courssub{Le bon usage des antibiotiques}

Pour que les antibiotiques restent efficaces, chacun doit adopter les bons comportements :

\begin{itemize}
\item ne prendre un antibiotique que sur \textbf{prescription du médecin} ;
\item respecter la \textbf{dose} prescrite (nombre de comprimés par jour) ;
\item respecter la \textbf{durée} du traitement, même si l'on se sent guéri avant la fin, afin d'éliminer toutes les bactéries et d'éviter l'apparition de résistants ;
\item ne jamais \textbf{s'automédiquer}, ne jamais utiliser les restes d'un ancien traitement ni l'ordonnance d'un autre malade ;
\item ne pas demander d'antibiotique pour une maladie virale.
\end{itemize}

\begin{aretenir}[L'essentiel sur l'antibiothérapie]
\begin{itemize}
\item L'\textbf{antibiothérapie} est le traitement d'une maladie microbienne à l'aide d'antibiotiques.
\item Un \textbf{antibiotique} tue les bactéries ou bloque leur multiplication, aidant ainsi le système immunitaire à éliminer l'infection.
\item Le premier antibiotique, la \textbf{pénicilline}, a été découvert par Fleming en 1928 grâce à l'observation d'une zone d'inhibition autour d'un \emph{Penicillium}.
\item L'\textbf{antibiogramme} permet de choisir l'antibiotique efficace : plus l'auréole d'inhibition est grande, plus l'antibiotique est actif.
\item Les antibiotiques sont \textbf{inefficaces contre les virus} ; leur mauvais usage provoque l'\textbf{antibiorésistance}.
\end{itemize}
\end{aretenir}

\coursec{II. La sérothérapie}

\courssub{Situation d'accroche}

Revenons à l'enfant de Bafoussam blessé par un clou rouillé. Le médecin ne lui a pas donné d'antibiotique : il lui a injecté un \textbf{sérum antitétanique}. Pourquoi ? Parce que le danger, dans le tétanos, n'est pas seulement la bactérie elle-même, mais surtout la \textbf{toxine} (poison) qu'elle fabrique et qui attaque le système nerveux. Il faut neutraliser ce poison \emph{immédiatement}, sans attendre que l'organisme fabrique lui-même ses anticorps, ce qui prendrait plusieurs jours. D'où vient ce sérum si rapide ?

\courssub{Principe : fournir des anticorps tout prêts}

Nous avons vu dans les leçons précédentes que, lors d'une infection, les lymphocytes fabriquent des \cle{anticorps}, des protéines capables de neutraliser les microbes ou leurs toxines. Mais cette fabrication demande du temps (plusieurs jours). Dans certaines maladies très rapides et très graves (tétanos, morsure de serpent venimeux, rage), ce délai peut être fatal.

L'idée de la sérothérapie est alors la suivante : \textbf{injecter au malade des anticorps déjà fabriqués}, prélevés dans le sang d'un individu (souvent un animal, comme le cheval) qui a été préalablement immunisé contre le microbe ou la toxine. Le liquide injecté, obtenu à partir du sang débarrassé de ses cellules, s'appelle un \cle{sérum}.

\begin{definition}[La sérothérapie]
La \cle{sérothérapie} est le \textbf{traitement d'une maladie par l'injection d'un sérum contenant des anticorps tout prêts}, fabriqués par un autre individu (homme ou animal) immunisé contre le microbe ou la toxine responsable de la maladie.
\end{definition}

\begin{experience}[Comment obtient-on un sérum thérapeutique ? (principe)]
\textbf{Protocole.}
\begin{enumerate}
\item On injecte à un animal (par exemple un cheval) de petites doses, de plus en plus fortes, de la toxine du tétanos rendue inoffensive.
\item Le système immunitaire du cheval réagit en fabriquant de grandes quantités d'\textbf{anticorps antitétaniques}.
\item On prélève du sang du cheval, on élimine les cellules sanguines et on purifie le liquide : on obtient le \textbf{sérum antitétanique}, riche en anticorps.
\item Ce sérum est injecté au malade : les anticorps qu'il contient \textbf{neutralisent immédiatement} la toxine présente dans son organisme.
\end{enumerate}
\textbf{Conclusion.} Le sérum apporte au malade des anticorps tout prêts, qui agissent sans délai.
\end{experience}

\courssub{Caractères de la sérothérapie}

\begin{itemize}
\item \textbf{Action immédiate} : les anticorps injectés agissent dès l'injection, ce qui est précieux en cas d'urgence.
\item \textbf{Action de courte durée} : les anticorps reçus sont progressivement détruits et éliminés ; la protection ne dure que quelques semaines.
\item \textbf{Immunité dite « passive »} : l'organisme du malade ne fabrique rien lui-même ; il reçoit les anticorps tout faits. Il ne garde \textbf{aucun souvenir} de l'agression : si le microbe revient plus tard, il n'est pas mieux protégé qu'avant.
\item La sérothérapie est surtout \textbf{curative} : on l'utilise quand la maladie (ou le danger) est déjà là.
\end{itemize}

\begin{exemplebox}[Exemples d'utilisation des sérums]
\begin{itemize}
\item le \textbf{sérum antitétanique}, injecté en urgence après une blessure souillée (clou rouillé, plaie de terre) chez une personne non vaccinée ;
\item le \textbf{sérum antivenimeux}, utilisé dans les centres de santé camerounais après une morsure de serpent venimeux (vipère, cobra) : il contient des anticorps qui neutralisent le venin ;
\item le \textbf{sérum antirabique}, contre la rage après morsure d'un chien enragé.
\end{itemize}
\end{exemplebox}

\begin{aretenir}[L'essentiel sur la sérothérapie]
La sérothérapie consiste à injecter au malade un \textbf{sérum contenant des anticorps tout prêts}. Son action est \textbf{immédiate mais de courte durée} : c'est une immunité \textbf{passive}, sans mémoire. Elle sert surtout à \textbf{soigner} (urgences : tétanos, venin de serpent, rage).
\end{aretenir}

\coursec{III. La vaccinothérapie}

\courssub{Situation d'accroche}

Troisième scène de notre situation d'accroche : toute l'école est convoquée pour une campagne de \textbf{vaccination} contre la fièvre jaune. Ici, personne n'est malade ! Pourquoi injecter quelque chose à des enfants en bonne santé ? Parce que le vaccin prépare le système immunitaire \emph{avant} la rencontre avec le microbe. C'est l'idée géniale de la vaccination, née des travaux de Jenner et de Pasteur.

\courssub{Les expériences historiques : de Jenner à Pasteur}

\begin{experience}[L'expérience de Jenner (1796) : la première vaccination]
\textbf{Observation.} Le médecin anglais Edward \textbf{Jenner} remarque que les fermières qui ont attrapé la « vaccine », une maladie bénigne des vaches (pustules sur les pis), ne sont \textbf{jamais atteintes par la variole}, une maladie humaine grave et souvent mortelle.

\textbf{Hypothèse.} La vaccine, maladie voisine mais bénigne de la variole, protégerait contre la variole.

\textbf{Expérience.} En 1796, Jenner prélève du liquide d'une pustule de vaccine sur la main d'une fermière et l'inocule (l'introduit par une petite scarification) à un jeune garçon de 8 ans, James Phipps. L'enfant développe une maladie légère, puis guérit. Quelques semaines plus tard, Jenner lui inocule du liquide de pustules de \textbf{variole} : l'enfant \textbf{ne développe pas la maladie}.

\textbf{Interprétation.} Le contact préalable avec le microbe de la vaccine a rendu l'organisme de l'enfant \textbf{capable de se défendre} contre la variole : il est devenu \textbf{immunisé}.

\textbf{Conclusion.} On peut protéger un individu contre une maladie grave en lui inoculant un microbe voisin, inoffensif ou affaibli. Jenner appelle cette méthode \textbf{vaccination} (du latin \emph{vacca}, la vache).
\end{experience}

\begin{savaistu}[Pasteur et les vaccins à microbes atténués]
À la fin du XIX\textsuperscript{e} siècle, le Français Louis \textbf{Pasteur} généralise l'idée de Jenner. En travaillant sur le choléra des poules, il découvre par hasard que des cultures de bactéries \textbf{vieillies} (affaiblies) ne rendent plus les poules malades, mais les \textbf{protègent} contre les cultures fraîches et virulentes. Il applique ensuite cette méthode des \textbf{microbes atténués} au charbon des moutons, puis, en 1885, à la \textbf{rage} : il sauve le jeune Joseph Meister, mordu par un chien enragé. La variole, première maladie combattue par la vaccination, a été totalement \textbf{éliminée de la planète} en 1980 grâce aux campagnes mondiales de vaccination : c'est l'un des plus grands succès de la médecine.
\end{savaistu}

\courssub{Définition et principe de la vaccinothérapie}

\begin{definition}[La vaccinothérapie]
La \cle{vaccinothérapie} (ou \textbf{vaccination}) est la \textbf{prévention d'une maladie microbienne par l'administration d'un vaccin}. Un \cle{vaccin} est une préparation contenant le microbe (ou une partie du microbe, ou sa toxine) \textbf{tué, atténué ou modifié}, c'est-à-dire rendu \textbf{incapable de provoquer la maladie}, mais capable de stimuler le système immunitaire.
\end{definition}

Le principe de la vaccination repose sur deux propriétés du système immunitaire étudiées dans les leçons précédentes : la \textbf{spécificité} et la \textbf{mémoire}.

\begin{enumerate}
\item Le vaccin introduit dans l'organisme le microbe \textbf{affaibli ou tué} (ou sa toxine inactivée) : il ne peut pas causer la maladie, mais il porte encore les marques (antigènes) que le système immunitaire reconnaît comme étrangères.
\item Les lymphocytes réagissent : ils fabriquent des \textbf{anticorps} spécifiques contre ce microbe, exactement comme lors d'une vraie infection.
\item Surtout, l'organisme garde en mémoire ce microbe grâce à des \textbf{lymphocytes mémoires} qui persistent longtemps dans le sang.
\item Si, des mois ou des années plus tard, le \textbf{vrai microbe} agresse l'organisme, la réponse immunitaire est \textbf{plus rapide et plus forte} : le microbe est détruit avant de provoquer la maladie. La personne est \textbf{immunisée}.
\end{enumerate}

\begin{popfigure}
\begin{center}
\begin{tikzpicture}[scale=0.92, every node/.style={font=\small}]
% Étape 1
\node[draw=popblue, thick, fill=popblueL, rounded corners, align=center, text width=3.1cm] (e1) at (0,0) {\textbf{1. Vaccination}\\ injection du microbe tué ou atténué};
% Étape 2
\node[draw=poppurple, thick, fill=poppurpleL, rounded corners, align=center, text width=3.1cm] (e2) at (4.6,0) {\textbf{2. Réaction}\\ fabrication d'anticorps et de lymphocytes mémoires};
% Étape 3
\node[draw=popgreen, thick, fill=popgreenL, rounded corners, align=center, text width=3.1cm] (e3) at (9.2,0) {\textbf{3. Mémoire}\\ l'organisme garde le souvenir du microbe};
% Étape 4
\node[draw=poporange, thick, fill=poporangeL, rounded corners, align=center, text width=3.4cm] (e4) at (4.6,-2.6) {\textbf{4. Contact avec le vrai microbe}\\ défense rapide et forte : pas de maladie};
\draw[->, very thick, popdark] (e1) -- (e2);
\draw[->, very thick, popdark] (e2) -- (e3);
\draw[->, very thick, popdark] (e3.south) |- (e4.east);
\draw[->, very thick, popdark] (e1.south) |- (e4.west);
\end{tikzpicture}
\end{center}
\legende{Principe de la vaccination. Le vaccin (microbe tué ou atténué) déclenche une réaction immunitaire avec fabrication d'anticorps et de lymphocytes mémoires. Lors d'un contact ultérieur avec le vrai microbe, la défense est rapide et efficace : la maladie est évitée.}
\end{popfigure}

\courssub{Caractères de la vaccinothérapie}

\begin{itemize}
\item \textbf{Action retardée} : la protection n'est pas immédiate ; il faut quelques jours à quelques semaines pour que l'organisme fabrique ses anticorps. C'est pourquoi on vaccine \emph{avant} le contact avec le microbe.
\item \textbf{Action durable} : grâce aux lymphocytes mémoires, la protection dure des \textbf{années}, parfois toute la vie. Des \textbf{rappels} (nouvelles injections) sont parfois nécessaires pour entretenir la mémoire immunitaire.
\item \textbf{Immunité dite « active »} : c'est l'organisme du vacciné qui \textbf{fabrique lui-même} ses anticorps.
\item La vaccination est une méthode \textbf{préventive} : elle se fait sur des personnes \textbf{en bonne santé}, pour éviter la maladie.
\end{itemize}

\begin{exemplebox}[Quelques vaccins du calendrier vaccinal au Cameroun]
Le Programme Élargi de Vaccination (PEV) du Cameroun offre gratuitement aux enfants des vaccins contre : la \textbf{tuberculose} (BCG, à la naissance), la \textbf{poliomyélite}, la \textbf{diphtérie}, le \textbf{tétanos}, la \textbf{coqueluche}, la \textbf{rougeole}, la \textbf{fièvre jaune} (à 9 mois), l'\textbf{hépatite B}. Grâce à ces vaccinations, la poliomyélite a presque disparu du pays et de nombreuses vies d'enfants sont sauvées chaque année.
\end{exemplebox}

\begin{attention}[Ne pas confondre vaccin et sérum]
Le \textbf{vaccin} contient le \textbf{microbe tué ou atténué} (ou sa toxine modifiée) : il pousse l'organisme à fabriquer \emph{lui-même} ses anticorps (immunité active, lente mais durable). Le \textbf{sérum} contient des \textbf{anticorps tout prêts} fabriqués par un autre individu (immunité passive, immédiate mais courte). Ne dis jamais « vaccin » pour désigner un sérum, et inversement !
\end{attention}

\begin{aretenir}[L'essentiel sur la vaccinothérapie]
La vaccination consiste à administrer un \textbf{vaccin} (microbe tué, atténué ou toxine modifiée) qui stimule le système immunitaire \textbf{sans provoquer la maladie}. L'organisme fabrique ses propres anticorps et garde une \textbf{mémoire} du microbe : c'est une immunité \textbf{active}, \textbf{lente à s'installer mais durable}. La vaccination est une méthode de \textbf{prévention}.
\end{aretenir}

\coursec{IV. La séro-vaccinothérapie}

\courssub{Situation d'accroche}

Un enfant non vacciné est mordu par un chien suspect de rage à Yaoundé. La rage est une maladie mortelle dès que les signes apparaissent, mais le virus met plusieurs jours à atteindre le système nerveux. Le médecin dispose donc d'un court délai pour agir. Que faire ? Le sérum seul protège immédiatement mais quelques semaines seulement ; le vaccin seul protège durablement mais trop lentement. La solution : \textbf{associer les deux}. C'est la \cle{séro-vaccinothérapie}.

\courssub{Définition et principe}

\begin{definition}[La séro-vaccinothérapie]
La \cle{séro-vaccinothérapie} est la méthode qui consiste à \textbf{associer, chez un même malade, l'injection d'un sérum et l'administration d'un vaccin} contre le même microbe (ou la même toxine).
\end{definition}

Le principe est de combiner les avantages des deux méthodes précédentes :

\begin{itemize}
\item le \textbf{sérum} fournit des \textbf{anticorps tout prêts} qui neutralisent \textbf{immédiatement} le microbe ou sa toxine : le malade est protégé \emph{tout de suite} ;
\item le \textbf{vaccin} stimule \textbf{en même temps} le système immunitaire du malade, qui fabrique ses propres anticorps et des lymphocytes mémoires : la protection devient \emph{durable}.
\end{itemize}

Ainsi, quand les anticorps du sérum disparaissent (au bout de quelques semaines), ceux fabriqués par l'organisme grâce au vaccin prennent le relais. La séro-vaccinothérapie assure donc une protection \textbf{à la fois immédiate et durable}.

\begin{exemplebox}[Exemples d'utilisation de la séro-vaccinothérapie]
\begin{itemize}
\item \textbf{La rage} : après la morsure par un animal enragé, on injecte au blessé un sérum antirabique (action immédiate) et on commence en même temps une série d'injections de vaccin antirabique (protection durable).
\item \textbf{Le tétanos} : chez une personne blessée non vaccinée (ou mal vaccinée), on associe le sérum antitétanique et le vaccin antitétanique, avec des rappels ultérieurs.
\end{itemize}
\end{exemplebox}

\begin{aretenir}[L'essentiel sur la séro-vaccinothérapie]
La séro-vaccinothérapie associe un \textbf{sérum} (anticorps tout prêts, action \textbf{immédiate}) et un \textbf{vaccin} (stimulation de l'immunité propre, action \textbf{durable}). Elle est utilisée en urgence, quand le malade a déjà été exposé à un microbe dangereux (rage, tétanos) et qu'il faut le protéger à la fois tout de suite et longtemps.
\end{aretenir}

\coursec{V. Comparaison et synthèse des méthodes d'aide à l'immunité}

\courssub{Tableau comparatif}

\begin{center}
\begin{tabularx}{\linewidth}{|l|X|X|X|X|}
\hline
\rowcolor{popblueL}
\textbf{Critère} & \textbf{Antibiothérapie} & \textbf{Sérothérapie} & \textbf{Vaccinothérapie} & \textbf{Séro-vaccinothérapie} \\
\hline
Substance utilisée & Antibiotique & Sérum (anticorps tout prêts) & Vaccin (microbe tué ou atténué, toxine modifiée) & Sérum + vaccin \\
\hline
Cible & Les bactéries (jamais les virus) & Le microbe ou sa toxine & Le système immunitaire du vacciné & Le microbe et le système immunitaire \\
\hline
Début d'action & Rapide & Immédiate & Lente (jours à semaines) & Immédiate (sérum) puis durable (vaccin) \\
\hline
Durée de la protection & Le temps du traitement & Courte (quelques semaines) & Longue (années, parfois à vie) & Longue \\
\hline
But & Soigner (curatif) & Soigner (curatif, urgence) & Prévenir (préventif) & Soigner et prévenir (urgence) \\
\hline
Exemples & Angine, tuberculose, choléra & Tétanos, morsure de serpent & Rougeole, fièvre jaune, poliomyélite & Rage, tétanos \\
\hline
\end{tabularx}
\end{center}

\courssub{Ce qu'il faut retenir pour choisir la bonne méthode}

\begin{methode}[Choisir la méthode adaptée à une situation]
\begin{enumerate}
\item \textbf{Identifier l'agresseur} : s'agit-il d'une bactérie, d'un virus, d'une toxine, d'un venin ?
\item \textbf{Identifier le moment} : la personne est-elle déjà malade ou en danger (il faut \emph{soigner}) ou est-elle en bonne santé (il faut \emph{prévenir}) ?
\item \textbf{Choisir} :
\begin{itemize}
\item infection bactérienne déclarée $\Rightarrow$ \textbf{antibiothérapie} (éventuellement guidée par un antibiogramme) ;
\item urgence avec toxine ou venin déjà présent $\Rightarrow$ \textbf{sérothérapie} ;
\item personne saine à protéger à long terme $\Rightarrow$ \textbf{vaccinothérapie} ;
\item urgence chez une personne exposée à un microbe très dangereux (rage, tétanos) $\Rightarrow$ \textbf{séro-vaccinothérapie}.
\end{itemize}
\end{enumerate}
\end{methode}

\begin{attention}[Pièges fréquents au BEPC]
\begin{itemize}
\item Ne dis pas qu'un antibiotique « stimule » ou « renforce » le système immunitaire : il \textbf{tue les bactéries ou bloque leur multiplication}, ce qui \emph{aide} les défenses naturelles.
\item Ne dis pas que le vaccin contient des anticorps : il contient le \textbf{microbe tué ou atténué} ; ce sont les anticorps que l'organisme \emph{fabrique ensuite}.
\item Ne dis pas que la sérothérapie protège longtemps : ses anticorps disparaissent en quelques semaines (immunité passive, sans mémoire).
\item Précise toujours si la méthode est \textbf{curative} (elle soigne) ou \textbf{préventive} (elle empêche la maladie).
\end{itemize}
\end{attention}

\begin{exoresolu}[Problème de synthèse : choisir la bonne aide à l'immunité]
\textbf{Énoncé.} Pour chacune des situations suivantes, indique la méthode d'aide au système immunitaire la plus adaptée (antibiothérapie, sérothérapie, vaccinothérapie ou séro-vaccinothérapie) et justifie ton choix.
\begin{enumerate}
\item Un élève de 3ème souffre d'une angine à streptocoques confirmée.
\item Un nourrisson de 9 mois, en bonne santé, est amené au centre de santé pour être protégé contre la fièvre jaune.
\item Un cultivateur de Foumbot, jamais vacciné, se blesse profondément au pied avec une houe souillée de terre ; on craint le tétanos.
\item Une femme enceinte est mordue par un chien errant suspect de rage.
\end{enumerate}
\tcblower
\textbf{Solution détaillée.}
\begin{enumerate}
\item \textbf{Antibiothérapie.} L'angine est une infection \textbf{bactérienne} (streptocoque) déjà déclarée : il faut la soigner. L'antibiotique tuera les bactéries ou bloquera leur multiplication, aidant le système immunitaire à éliminer l'infection.
\item \textbf{Vaccinothérapie.} Le nourrisson est \textbf{en bonne santé} : il s'agit de \textbf{prévenir} la fièvre jaune. Le vaccin (virus atténué) stimulera son système immunitaire, qui fabriquera ses propres anticorps et gardera une mémoire durable du virus.
\item \textbf{Séro-vaccinothérapie.} Le cultivateur est exposé à un danger grave et immédiat (la toxine tétanique) et n'est pas vacciné. Le \textbf{sérum} antitétanique neutralisera immédiatement la toxine, tandis que le \textbf{vaccin} antitétanique déclenchera la fabrication d'anticorps par son organisme pour une protection durable (avec des rappels).
\item \textbf{Séro-vaccinothérapie.} La rage est mortelle et le virus agit vite : il faut une protection \textbf{immédiate} (sérum antirabique, anticorps tout prêts) \textbf{et durable} (vaccin antirabique, qui stimule l'immunité propre de la victime).
\end{enumerate}
\end{exoresolu}

\begin{exercice}[Pour vérifier tes connaissances]
\begin{enumerate}
\item Définis l'antibiothérapie et la sérothérapie, en précisant pour chacune la substance utilisée.
\item Raconte en quelques phrases l'observation de Fleming et explique sa conclusion.
\item Quelle différence fondamentale existe-t-il entre un vaccin et un sérum ? Pourquoi dit-on que la vaccination confère une immunité « active » et la sérothérapie une immunité « passive » ?
\item Sur un antibiogramme, le disque de l'antibiotique X est entouré d'une auréole de \SI{25}{mm} et celui de l'antibiotique Y d'une auréole de \SI{8}{mm}. Lequel choisir ? Pourquoi ?
\item Ton petit frère a la grippe. Ta mère veut lui donner l'amoxicilline qui reste dans l'armoire à pharmacie. Quels conseils lui donnes-tu ? Justifie.
\item Un enfant non vacciné est mordu par un chien enragé. Le médecin lui injecte un sérum antirabique puis lui administre le vaccin antirabique. Nomme cette méthode et explique pourquoi elle est préférable au sérum seul ou au vaccin seul.
\end{enumerate}
\end{exercice}

\begin{corrige}[Exercice : Pour vérifier tes connaissances]
\begin{enumerate}
\item L'antibiothérapie est le traitement d'une maladie microbienne à l'aide d'antibiotiques, substances qui tuent les bactéries ou bloquent leur multiplication. La sérothérapie est le traitement d'une maladie par l'injection d'un sérum contenant des anticorps tout prêts, fabriqués par un autre individu immunisé.
\item En 1928, Fleming observe qu'une boîte de Pétri ensemencée de staphylocoques a été contaminée par une moisissure, le \emph{Penicillium}. Autour de la moisissure, les bactéries ont disparu (zone claire d'inhibition). Il en conclut que le \emph{Penicillium} libère une substance, la pénicilline, qui tue les bactéries ou empêche leur multiplication.
\item Le vaccin contient le microbe tué ou atténué (ou sa toxine modifiée) : il pousse l'organisme à fabriquer lui-même ses anticorps et des lymphocytes mémoires, d'où une immunité « active », lente mais durable. Le sérum contient des anticorps tout prêts fabriqués par un autre individu : l'organisme ne fabrique rien et ne garde aucune mémoire, d'où une immunité « passive », immédiate mais de courte durée.
\item Il faut choisir l'antibiotique X, car son auréole d'inhibition (\SI{25}{mm}) est plus grande que celle de Y (\SI{8}{mm}) : or, plus l'auréole est grande, plus l'antibiotique est efficace contre la bactérie.
\item Je conseille à ma mère de ne pas lui donner cet antibiotique, pour trois raisons : la grippe est une maladie \textbf{virale}, or les antibiotiques sont inefficaces contre les virus ; il ne faut jamais s'automédiquer ni utiliser les restes d'un traitement ; cet usage abusif favorise l'apparition de bactéries résistantes (antibiorésistance). Il faut consulter un médecin si l'état de l'enfant s'aggrave.
\item Cette méthode est la \textbf{séro-vaccinothérapie}. Le sérum seul protège immédiatement mais quelques semaines seulement (immunité passive sans mémoire) ; le vaccin seul protège durablement mais trop lentement, alors que le virus de la rage agit vite. L'association des deux assure une protection à la fois immédiate (grâce aux anticorps du sérum) et durable (grâce aux anticorps et aux lymphocytes mémoires fabriqués après la vaccination).
\end{enumerate}
\end{corrige}

\coursec{II. La sérothérapie}

Dans la section précédente, nous avons vu que l'antibiothérapie consiste à utiliser des substances chimiques, les antibiotiques, pour détruire ou bloquer les bactéries responsables d'une maladie. Mais les antibiotiques sont impuissants contre les \cle{toxines} (poisons fabriqués par certains microbes ou injectés par des animaux venimeux) : aucun antibiotique ne peut neutraliser le venin d'un serpent ou la toxine du tétanos. Face à ces situations d'urgence, la médecine dispose d'une autre arme, découverte à la fin du XIX\textsuperscript{e} siècle : la \cle{sérothérapie}. C'est elle que nous allons étudier maintenant.

\courssub{Situation d'introduction : une urgence au centre de santé}

\begin{exemplebox}[Deux situations concrètes au Cameroun]
\textbf{Situation 1 :} Dans un village de la région de l'Extrême-Nord, un cultivateur est mordu au pied par une vipère alors qu'il travaille dans son champ de mil. Sa jambe enfle rapidement, des signes de coagulation sanguine apparaissent. Transporté d'urgence au centre de santé, il reçoit immédiatement une injection de \textbf{sérum antivenimeux polyvalent}. En quelques heures, l'œdème régresse et son état se stabilise.

\textbf{Situation 2 :} Une élève de 3ème se blesse profondément au pied avec un clou rouillé en rentrant de l'école. À l'hôpital de district, le médecin constate qu'elle n'a jamais reçu le vaccin antitétanique (ou que son calendrier vaccinal n'est pas à jour). Il lui injecte en urgence du \textbf{sérum antitétanique} (immunoglobulines tétaniques) pour la protéger de la maladie.
\end{exemplebox}

\textbf{Problème scientifique :} Dans les deux cas, le malade n'a pas le temps de fabriquer lui-même ses anticorps : le venin ou la toxine agit en quelques heures, alors que la réponse immunitaire active (production d'anticorps par les lymphocytes B) nécessite plusieurs jours. D'où viennent alors les anticorps qui neutralisent le poison et sauvent le malade ? Comment les obtient-on ? C'est à ces questions que répond la sérothérapie.

\courssub{Le fait expérimental fondateur : l'expérience de Behring et Kitasato (1890)}

\begin{experience}[L'expérience historique de Behring et Kitasato (1890) -- Démonstration du principe de la sérothérapie]
En 1890, le médecin allemand Émile von Behring et le bactériologiste japonais Shibasaburo Kitasato réalisent l'expérience décisive suivante :

\textbf{Protocole :}
\begin{enumerate}
\item \textbf{Immunisation d'un animal donneur :} On injecte à un cheval (ou un mouton) des doses \textbf{croissantes} de toxine tétanique \textbf{atténuée} (toxinoïde : toxine traitée par le formol pour perdre sa dangerosité mais garder son pouvoir antigénique), à plusieurs reprises, pendant plusieurs semaines.
\item \textbf{Réponse de l'animal :} L'animal ne tombe pas malade. Son système immunitaire reconnaît la toxine atténuée comme un antigène et fabrique massivement des \textbf{anticorps spécifiques} (immunoglobulines) dirigés contre cette toxine. On dit que l'animal est \textbf{immunisé} (immunité active artificielle).
\item \textbf{Prélèvement du sérum :} On prélève du sang sur cet animal immunisé, on laisse coaguler (le fibrine et les cellules forment le caillot) et on récupère le \textbf{sérum} (liquide surnageant, dépourvu de cellules et de fibrine, mais riche en anticorps).
\item \textbf{Transfert passif chez un animal récepteur :} On injecte ce sérum immune à une souris saine (non immunisée). Immédiatement après, on inocule à cette souris une dose mortelle de toxine tétanique \textbf{native} (active).
\item \textbf{Témoin :} Une souris témoin reçoit la même dose mortelle de toxine, mais \textbf{sans} injection préalable de sérum.
\end{enumerate}

\textbf{Observations :}
\begin{itemize}
\item La souris ayant reçu le sérum de l'animal immunisé \textbf{survit} : elle est protégée contre l'effet létal de la toxine.
\item La souris témoin, qui n'a pas reçu de sérum, \textbf{meurt} rapidement des suites du tétanos.
\end{itemize}

\textbf{Interprétation :} La protection de la souris réceptrice ne provient pas de son propre système immunitaire (elle n'a jamais été en contact avec l'antigène auparavant, et le délai est trop court pour une synthèse active). Elle provient exclusivement des \textbf{anticorps apportés par le sérum} injecté. Ces anticorps se sont fixés sur la toxine (antigène) pour former des complexes antigène-anticorps inoffensifs : c'est la \textbf{neutralisation}.

\textbf{Conclusion :} On peut protéger un individu menacé ou malade en lui injectant le sérum d'un individu (ou animal) immunisé contre le même agresseur. C'est le principe de la \textbf{sérothérapie} (ou immunothérapie passive).
\end{experience}

\begin{popfigure}
\begin{tikzpicture}[
  box/.style={draw=popblue, fill=popblueL, rounded corners=4pt, align=center, minimum height=1.4cm, text width=3.2cm, font=\small, inner sep=3pt},
  etape/.style={font=\small\bfseries, text=popdark, yshift=8pt},
  fleche/.style={-{Stealth[length=3mm]}, very thick, poporange},
  node distance=1.2cm and 0.5cm]
% Ligne 1 : Immunisation
\node[box, align=center] (animal){Cheval (ou mouton)\\ recevant des doses croissantes\\ de toxine tétanique atténuée (toxinoïde)};
\node[etape, above=of animal, align=center]{1. Immunisation active\\ de l'animal donneur};
% Ligne 2 : Prélèvement
\node[box, fill=poppurpleL, draw=poppurple, right=of animal, align=center] (serum){Prélèvement sanguin,\\ coagulation, récupération\\ du sérum riche en anticorps anti-tétaniques};
\node[etape, above=of serum, align=center]{2. Préparation du sérum\\ thérapeutique};
% Ligne 3 : Injection au malade
\node[box, fill=popgreenL, draw=popgreen, right=of serum, align=center] (malade){Injection du sérum immune\\ à l'homme (ou animal) exposé\\ au danger (toxine, venin)};
\node[etape, above=of malade, align=center]{3. Injection thérapeutique\\ (urgence)};
% Ligne 4 : Résultat
\node[box, fill=popgoldL, draw=popgold, below=of malade, align=center] (gueri){Les anticorps injectés se fixent\\ sur la toxine/venin :\\ \textbf{neutralisation immédiate}\\ $\rightarrow$ Guérison / Protection};
\node[etape, below=of gueri, align=center]{4. Effet protecteur\\ (immunité passive)};
% Flèches
\draw[fleche] (animal) -- (serum);
\draw[fleche] (serum) -- (malade);
\draw[fleche] (malade) -- (gueri);
\end{tikzpicture}
\legende{Schéma du protocole de Behring et Kitasato (1890) : 1. Immunisation active d'un animal donneur par un toxinoïde ; 2. Prélèvement et purification du sérum contenant les anticorps ; 3. Injection de ce sérum à un sujet en danger ; 4. Neutralisation immédiate de l'agresseur par les anticorps apportés (immunité passive).}
\end{popfigure}

\courssub{Définition et principe de la sérothérapie}

\begin{definition}[La sérothérapie]
La \cle{sérothérapie} est une méthode thérapeutique (curative ou préventive d'urgence) qui consiste à injecter à un malade un \textbf{sérum} (dérivé du sang) contenant des \textbf{anticorps spécifiques (immunoglobulines) déjà fabriqués} par un autre individu (homme ou animal) immunisé contre le même agent agresseur (toxine, venin, virus).
\end{definition}

Le \textbf{principe} est simple à retenir : au lieu d'attendre que l'organisme malade fabrique lui-même ses anticorps (réponse active lente, 5 à 10 jours), on lui apporte des \textbf{anticorps « tout faits »}, opérationnels immédiatement. Ces anticorps se fixent sans délai sur l'antigène (toxine, venin, particule virale) et le neutralisent par agglutination, précipitation ou blocage de ses sites actifs.

\begin{propriete}[Caractères de l'immunité conférée par la sérothérapie (admis)]
La sérothérapie confère une \textbf{immunité passive artificielle} caractérisée par :
\begin{itemize}
\item une action \textbf{immédiate} : les anticorps sont fonctionnels dès leur arrivée dans la circulation sanguine (précieux en urgence vitale) ;
\item une durée \textbf{limitée (temporaire)} : les anticorps injectés (protéines étrangères) sont progressivement catabolisés (détruits) par l'organisme hôte sans être renouvelés. La demi-vie des IgG est d'environ 21 à 28 jours. La protection disparaît en \textbf{quelques semaines (3 à 6 semaines)}.
\item l'\textbf{absence de mémoire immunologique} : l'organisme du receveur n'a pas été stimulé par l'antigène, ses lymphocytes B n'ont pas été activés, il ne garde \textbf{aucun souvenir} de l'agression. Une nouvelle exposition ultérieure le retrouvera sans défense.
\end{itemize}
On parle d'immunité \emph{passive} car le receveur \textbf{ne produit pas lui-même} ses anticorps : il les reçoit « clés en main ».
\end{propriete}

\begin{popfigure}
\begin{tikzpicture}
\begin{axis}[
  width=0.9\linewidth, height=6.5cm,
  xlabel={Temps après injection (jours)}, ylabel={Taux d'anticorps circulants (unités arbitraires)},
  xmin=0, xmax=45, ymin=0, ymax=110,
  xtick={0,7,14,21,28,35,42},
  grid=both, grid style={popline!40},
  axis line style={popdark, thick},
  label style={font=\small},
  tick label style={font=\small},
  legend style={font=\small, at={(0.98,0.98)}, anchor=north east, fill=white, draw=popline}]
\addplot[popcurve, popblue, very thick, domain=0:45, samples=300] {100*exp(-ln(2)/23*x)};
\addlegendentry{Anticorps apportés (IgG, demi-vie $\approx$ 23 j)}
\draw[dashed, poporange, thick] (axis cs:0,30) -- (axis cs:45,30) node[pos=0.5, above, font=\small, text=poporange]{Seuil de protection théorique};
\draw[popred, thick, <->] (axis cs:0,100) -- (axis cs:0,0) node[midway, left, font=\small, text=popred]{Pic immédiat};
\draw[popred, thick, <->] (axis cs:23,50) -- (axis cs:23,0) node[midway, left, font=\small, text=popred]{Demi-vie};
\end{axis}
\end{tikzpicture}
\legende{Évolution du taux d'anticorps après une injection de sérum (immunité passive). Le taux est maximal dès l'injection (action immédiate), puis décroît exponentiellement (demi-vie $\approx$ 3 semaines) car les anticorps ne sont pas renouvelés. La protection cesse quand le taux passe sous le seuil protecteur.}
\end{popfigure}

\courssub{Principaux sérums thérapeutiques et indications}

La sérothérapie est indiquée chaque fois qu'il faut neutraliser \textbf{rapidement} une toxine, un venin ou un virus, avant que l'immunité active du malade ne puisse se mettre en place. Les principaux sérums (ou immunoglobulines purifiées) utilisés en pratique médicale sont :

\begin{center}
\begin{tabularx}{\linewidth}{|l|X|X|}
\hline
\rowcolor{popblueL} \textbf{Sérum (Immunoglobulines)} & \textbf{Agresseur neutralisé} & \textbf{Indication principale (Urgence)} \\
\hline
Sérum \textbf{antitétanique} (SAT / Ig antitetaniques) & Toxine tétanique (\emph{Clostridium tetani}) & Blessure profonde, souillée, chez un sujet non vacciné ou vaccination incomplète \\
\hline
Sérum \textbf{antivenimeux} (polyvalent ou monovalent) & Venins de serpents (vipères, cobras, mambas\ldots) et scorpions & Morsure de serpent venimeux / piqûre de scorpion grave (signes généraux) \\
\hline
Immunoglobulines \textbf{antirabiques} (IAR) & Virus de la rage (Lyssavirus) & Morsure/griffure par animal suspecté enragé (associées au vaccin) \\
\hline
Sérum \textbf{antidiphtérique} & Toxine diphtérique (\emph{Corynebacterium diphtheriae}) & Diphérie déclarée (traitement curatif d'urgence) \\
\hline
Immunoglobulines \textbf{anti-hépatite B} & Virus de l'hépatite B (VHB) & Exposition accidentelle au sang (piqûre d'aiguille) chez non vacciné / Nourrisson de mère porteuse \\
\hline
\end{tabularx}
\end{center}

\begin{exemplebox}[Le sérum antivenimeux : un enjeu de santé publique au Cameroun]
Au Cameroun, les envenimations par morsures de serpents (vipères, cobras, mambas) sont fréquentes en milieu rural (Extrême-Nord, Nord, Adamaoua, Est, Sud), particulièrement pendant la saison des pluies et les travaux agricoles. Le \textbf{sérum antivenimeux polyvalent} (efficace contre les venins des principales espèces locales), conservé à la chaîne du froid ($+2^\circ\text{C}$ à $+8^\circ\text{C}$) dans les formations sanitaires, est le \textbf{seul traitement spécifique efficace}. Son injection précoce (idéalement $< 4$ h) sauve des vies et prévient les séquelles (nécroses, amputations, insuffisance rénale).
\begin{attention}[Conduite à tenir]
En cas de morsure : \textbf{immobiliser le membre} (garrot \emph{interdit} sauf cas extrême), \textbf{ne pas inciser, ne pas sucer}, \textbf{transporter la victime allongée} au centre de santé le plus proche \textbf{sans délai}. Le tradipraticien ne dispose pas de sérum antivenimeux.
\end{attention}
\end{exemplebox}

\courssub{Précautions d'emploi et effets indésirables}

\begin{attention}[Risque majeur : le choc anaphylactique (choc allergique systémique)]
Les sérums thérapeutiques classiques sont souvent des \textbf{sérums hétérologues} (produits chez l'animal : cheval, mouton). Ils contiennent donc des protéines animales (globulines sériques) reconnues comme \textbf{étrangères} par le système immunitaire du receveur humain.
\begin{itemize}
\item Chez un sujet sensibilisé (antécédent d'injection de sérum animal) ou parfois dès la première injection, cela peut déclencher une réaction d'hypersensibilité immédiate de type I (mediée par IgE) : \textbf{choc anaphylactique} (chute brutale de la tension artérielle, œdème du larynx, bronchospasme, urticaire, troubles digestifs). C'est une \textbf{urgence vitale absolue}.
\item Même les immunoglobulines humaines (homologues) comportent un risque allergique faible mais non nul.
\end{itemize}
\textbf{Règle d'or :} Le sérum doit \textbf{toujours être administré en milieu médical équipé} (hôpital, centre de santé, CMA), par du personnel formé, avec surveillance clinique pendant au moins 1 heure post-injection (adrenaline, corticoïdes, antihistaminiques, oxygène, perfusion disponibles). \textbf{Jamais d'auto-injection.}
\end{attention}

\begin{attention}[Test de sensibilité (pratique ancienne, moins utilisé avec Ig purifiées)]
Autrefois, on réalisait un test cutané (injection intradermique de sérum dilué) pour détecter une allergie avant l'injection complète. Avec les immunoglobulines purifiées modernes (moins de protéines étrangères), ce test est moins prédictif et souvent omis, mais la \textbf{surveillance médicale stricte pendant et après l'injection reste obligatoire}.
\end{attention}

\begin{attention}[Erreur fréquente au BEPC : Confondre SÉRUM et VACCIN]
C'est l'erreur n°1 dans les copies. Grave cette distinction dans ta mémoire :
\begin{center}
\begin{tabular}{|l|c|c|}
\hline
\rowcolor{popblueL} \textbf{Caractère} & \textbf{SÉRUM (Sérothérapie)} & \textbf{VACCIN (Vaccinothérapie)} \\
\hline
\textbf{Substance injectée} & \textbf{Anticorps} (protéines, IgG) & \textbf{Antigènes} (microbe tué/atténué, toxinoïde, fragment) \\
\hline
\textbf{Origine} & Individu/Animal \textbf{immunisé} (donneur) & Laboratoire (cultures, ingénierie) \\
\hline
\textbf{Rôle de l'organisme receveur} & \textbf{Passif} : reçoit les armes toutes faites & \textbf{Actif} : doit fabriquer ses propres armes \\
\hline
\textbf{Délai de protection} & \textbf{Immédiat} (heures) & \textbf{Différé} (1 à 3 semaines) \\
\hline
\textbf{Durée de protection} & \textbf{Courte} (semaines) & \textbf{Longue} (années, vie) \\
\hline
\textbf{Mémoire immunitaire} & \textbf{Aucune} & \textbf{Oui} (lymphocytes B mémoire) \\
\hline
\textbf{Indication type} & \textbf{Urgence} (venin, toxine, exposition virale) & \textbf{Prévention programmée} (calendrier vaccinal) \\
\hline
\end{tabular}
\end{center}
\emph{Astuce mnémotechnique :} \textbf{Sérum} = \textbf{S}ecours \textbf{S}oudain (anticorps \textbf{S}tockés) ; \textbf{Vaccin} = \textbf{V}ictoire \textbf{V}olontaire (l'organisme \textbf{V}abrique).
\end{attention}

\begin{savaistu}[Histoire des sciences : Le premier Nobel de Médecine]
Émile von Behring a reçu en 1901 le \textbf{tout premier prix Nobel de physiologie ou médecine} de l'histoire, « pour ses travaux sur la sérothérapie, spécialement son application contre la diphtérie, par lesquels il a ouvert une nouvelle route dans le domaine de la thérapeutique médicale et mis entre les mains du médecin une arme victorieuse contre la maladie et la mort ». Avant 1890, la diphtérie (angine fausse, « le croup ») tuait \textbf{50 \% des enfants atteints} par étouffement (membrane fibreuse dans la gorge). Le sérum antidiphtérique a fait chuter la mortalité à \textbf{< 10 \%}. C'est la première grande victoire de l'immunothérapie.
\end{savaistu}

\courssub{Exercice résolu : Choisir le bon traitement d'urgence}

\begin{exoresolu}[Sérum ou Antibiotique ? -- Cas cliniques au centre de santé d'Ekondo-Titi (Sud-Ouest)]
\textbf{Énoncé.} Trois malades arrivent simultanément au centre de santé. Pour chacun, indique le traitement d'urgence de première intention (antibiotique ou sérum) et justifie ton choix en te basant sur la nature de l'agresseur et le mécanisme d'action du traitement.
\begin{enumerate}
\item \textbf{Malade A :} Jeune cultivateur, morsure de serpent (vipère) au pied il y a 2 heures. Gonflement rapide de la cheville, saignement au point de morsure, nausées.
\item \textbf{Malade B :} Enfant de 5 ans, fièvre à 39°C, angine érythémateuse avec ganglions cervicaux douloureux. Test rapide streptocoque (TDR) positif.
\item \textbf{Malade C :} Adulte, blessure profonde et sale (clou rouillé) au pied il y a 6 heures. Carte de vaccination introuvable, pas de rappel antitétanique depuis plus de 10 ans.
\end{enumerate}

\tcblower
\textbf{Solution détaillée.}

\begin{itemize}
\item \textbf{Malade A : Traitement = Sérum antivenimeux (polyvalent).}
\textbf{Justification :} L'agresseur est le \textbf{venin} (mélange de toxines enzymatiques et neurotoxiques), pas une bactérie vivante. Les antibiotiques n'ont \textbf{aucune action} sur les toxines déjà injectées. Seule l'injection immédiate d'anticorps spécifiques (sérum) peut neutraliser le venin circulante et arrêter l'envenimation. L'antibiotique pourra être associé secondairement pour prévenir la surinfection bactérienne de la plaie, mais il ne traite pas l'urgence vitale.
\item \textbf{Malade B : Traitement = Antibiotique (ex: Amoxicilline ou Pénicilline V).}
\textbf{Justification :} L'agresseur est une \textbf{bactérie vivante} (\emph{Streptococcus pyogenes}, streptocoque du groupe A) qui se multiplie dans l'amygdale. L'antibiotique (bactéricide ou bactériostatique) va détruire ou bloquer la multiplication de cette bactérie, stoppant l'infection. Il n'y a pas de toxine préformée massive à neutraliser en urgence (la toxine streptococcique est produite \emph{in situ} par la bactérie vivante ; tuer la bactérie arrête la production).
\item \textbf{Malade C : Traitement = Sérum antitétanique (Immunoglobulines tétaniques humaines si possible).}
\textbf{Justification :} Le risque est le tétanos par contamination de la plaie par des spores de \emph{Clostridium tetani}. Le malade n'est \textbf{pas protégé} (pas de vaccination à jour $\rightarrow$ pas d'anticorps, pas de mémoire). L'incubation du tétanos est courte (3 à 14 jours en moyenne). Il n'a \textbf{pas le temps} de fabriquer ses anticorps s'on lui donnait un vaccin (toxinoïde) seul (réponse active en 10-15 jours). On doit lui apporter \textbf{immédiatement} des anticorps tout faits (sérum) pour neutraliser la toxine tétanique qui serait produite par les bactéries dans la plaie. \emph{Note : On pratiquera aussi le nettoyage chirurgical de la plaie et on démarrera une vaccination antitétanique (immunité active) en même temps, sur un autre site d'injection.}
\end{itemize}

\textbf{Bilan :} \textbf{Antibiotique} $\rightarrow$ cible la \textbf{bactérie vivante} (infection). \textbf{Sérum} $\rightarrow$ cible la \textbf{toxine/venin/virus} (intoxication/envenimation) par apport d'\textbf{anticorps prêts à l'emploi}.
\end{exoresolu}

\begin{aretenir}[L'essentiel sur la sérothérapie -- Fiche de révision BEPC]
\begin{itemize}
\item \textbf{Définition :} Injection d'un sérum contenant des \textbf{anticorps tout faits} (produits par un donneur immunisé) pour neutraliser un agresseur (toxine, venin, virus).
\item \textbf{Origine expérimentale :} Behring et Kitasato (1890) : transfert de sérum d'animal immunisé $\rightarrow$ protection d'animal naïf.
\item \textbf{Type d'immunité :} \textbf{Passive artificielle} (le receveur ne fabrique pas, ne mémorise pas).
\item \textbf{Caractères :} Action \textbf{immédiate} (urgence) mais \textbf{brève} (quelques semaines, pas de mémoire).
\item \textbf{Principaux sérums :} \textbf{Antitétanique, Antivenimeux, Antirabique, Antidiphtérique, Anti-Hépatite B}.
\item \textbf{Risque majeur :} \textbf{Choc anaphylactique} (allergie aux protéines animales du sérum) $\rightarrow$ Injection \textbf{uniquement en milieu médical} sous surveillance.
\item \textbf{Distinction clé :} \textbf{Sérum = Anticorps (passif, court)} ; \textbf{Vaccin = Antigènes (actif, long, mémoire)}.
\end{itemize}
\end{aretenir}

La sérothérapie protège donc \textbf{vite}, mais \textbf{pas longtemps}. Elle ne laisse aucun souvenir immunologique. Comment obtenir une protection \textbf{durable}, qui laisse à l'organisme le « souvenir » de l'agresseur pour se défendre seul lors des rencontres futures ? C'est précisément le rôle de la \textbf{vaccinothérapie} (immunité active artificielle), que nous allons étudier dans la section suivante.

\coursec{III. La vaccinothérapie}

Après avoir étudié la sérothérapie, qui apporte une protection \emph{immédiate mais temporaire} par l'injection d'anticorps tout faits, nous allons maintenant découvrir une tout autre stratégie : \textbf{apprendre à l'organisme à se défendre lui-même}. C'est le principe de la vaccinothérapie, une méthode \emph{préventive} qui a permis d'éradiquer la variole et de contrôler de nombreuses maladies mortelles.

\courssub{Situation d'introduction : les campagnes de vaccination au Cameroun}

Chaque année, dans les écoles, les centres de santé et les postes de vaccination avancés du Cameroun, des milliers d'enfants reçoivent des gouttes ou des injections dans le cadre du \textbf{Programme Élargi de Vaccination (PEV)}. Tu as sûrement conservé ton \emph{carnet de vaccination} : il mentionne le BCG (contre la tuberculose), le vaccin antiamaril (contre la fièvre jaune, obligatoire pour tout voyageur), le DTC (diphtérie, tétanos, coqueluche), le vaccin antipoliomyélitique (VPO ou VIP), et plus récemment le vaccin contre la rougeole, la fièvre jaune, l'hépatite B, ou le méningocoque.

\begin{savaistu}[Le carnet de vaccination, un document officiel]
Au Cameroun, le carnet de vaccination est un document légal. Il est exigé à l'inscription à l'école, pour l'obtention du passeport (certificat de vaccination antiamarile obligatoire pour sortir du territoire) et parfois pour l'embauche. Il témoigne de la politique de santé publique visant à instaurer une \textbf{immunité collective} : si une proportion suffisante de la population est vaccinée (généralement > 80-90 \%), la circulation du microbe est bloquée, protégeant ainsi même les non-vaccinés (nourrissons trop jeunes, immunodéprimés).
\end{savaistu}

Mais comment une simple injection peut-elle protéger pendant des années, voire toute une vie, contre une maladie que l'on n'a jamais eue ? L'histoire de cette découverte commence dans le laboratoire de Louis Pasteur.

\courssub{Fondements historiques : les expériences décisives de Pasteur}

\courssub{L'expérience du choléra des poules (1880) : la découverte de l'atténuation}

En 1880, Pasteur étudie le choléra des poules, une maladie bactérienne mortelle. Il observe qu'une culture de la bactérie (\emph{Pasteurella multocida}) laissée au repos pendant plusieurs semaines perd sa virulence : elle devient \textbf{atténuée}.

\begin{experience}[L'expérience de Pasteur sur le choléra des poules (1880)]
\textbf{Matériel :} Poules, cultures de \emph{Pasteurella multocida} (fraîche = virulente ; vieille de 3 mois = atténuée).
\textbf{Protocole :}
\begin{enumerate}
    \item \textbf{Groupe A (témoins) :} Inoculation de culture virulente fraîche $\rightarrow$ toutes meurent.
    \item \textbf{Groupe B :} Inoculation de culture atténuée (vieillie) $\rightarrow$ poules légèrement malades, puis guérison complète.
    \item \textbf{Groupe C (étape cruciale) :} Les poules guéries du groupe B reçoivent, après guérison, une injection de culture \emph{virulente fraîche}.
\end{enumerate}
\textbf{Observations :} Les poules du groupe C \textbf{ne meurent pas} ; elles résistent à la dose mortelle.
\textbf{Interprétation :} Le contact avec le microbe atténué a \textbf{protégé} les poules contre le microbe virulent. L'organisme a « appris » à se défendre.
\end{experience}

\courssub{La vaccination antirabique (1885) : du laboratoire à l'homme}

Fort de ce principe, Pasteur atténue le virus de la rage (en séchant la moelle épinière de lapins infectés). Le 6 juillet 1885, Joseph Meister, un enfant de 9 ans, est amené à Pasteur : il a été mordu 14 fois par un chien enragé. La mort est certaine sans intervention. Pasteur, bien que n'étant pas médecin, décide d'appliquer son protocole : une série d'injections de virus de plus en plus virulents (du plus atténué au plus frais) sur 10 jours.

\begin{savaistu}[Joseph Meister, le premier vacciné contre la rage]
Joseph Meister survit. C'est le triomphe de la \textbf{vaccinothérapie}. Pasteur a prouvé qu'on peut \emph{prévenir} une maladie mortelle en injectant l'agent pathogène sous une forme inoffensive. Joseph Meister deviendra plus tard gardien de l'Institut Pasteur. En 1940, pour éviter de livrer les secrets de l'Institut aux nazis, il se suicida. Son nom reste gravé dans l'histoire de la médecine.
\end{savaistu}

\courssub{Interprétation et définition de la vaccinothérapie}

\courssub{Ce qu'est un vaccin}

Un \cle{vaccin} est une préparation biologique qui contient l'agent pathogène (bactérie, virus, toxine) sous une forme \textbf{perdue de sa dangerosité} (pathogénicité nulle ou très faible) mais \textbf{conservant ses propriétés antigéniques} (capacité à être reconnu par le système immunitaire).

\begin{center}
\begin{tabularx}{\linewidth}{|l|X|X|}\hline
\rowcolor{popblueL}
\textbf{Type de vaccin} & \textbf{Principe d'obtention} & \textbf{Exemples (PEV Cameroun)} \\ \hline
\cle{Vaccin à microbe vivant atténué} & Culture dans des conditions qui affaiblissent le microbe (passages sur milieu inhabituel, température, mutations) & BCG (tuberculose), VPO (polio oral), Rougeole, Fièvre jaune (antiamaril) \\ \hline
\cle{Vaccin à microbe tué (inactivé)} & Traitement chimique (formol) ou physique (chaleur, rayonnements) qui tue le microbe & VIP (polio injectable), Coqueluche (partie du DTC), Rage (humain) \\ \hline
\cle{Vaccin à sous-unités / anatoxines} & Purification d'antigènes de surface ou toxines inactivées (anatoxines) & Diphtérie, Tétanos (anatoxines du DTC), Hépatite B (protéine de surface HBsAg) \\ \hline
\end{tabularx}
\end{center}

\courssub{Définition et principe}

\begin{definition}[Vaccinothérapie]
La \textbf{vaccinothérapie} est une méthode \textbf{préventive} qui consiste à administrer à un individu \textbf{sain} un \textbf{vaccin} (microbe tué, atténué, ou fragment de microbe) afin de \textbf{stimuler son propre système immunitaire} pour qu'il fabrique des anticorps et des cellules mémoires spécifiques, conférant ainsi une protection durable contre la maladie ciblée.
\end{definition}

\begin{propriete}[Immunité active vaccinale]
La vaccination confère une \textbf{immunité active} : c'est l'organisme du vacciné qui \textbf{produit lui-même} ses anticorps et ses lymphocytes mémoires.
\begin{itemize}
    \item \textbf{Lente à s'installer :} 2 à 3 semaines pour la primo-vaccination (réponse primaire).
    \item \textbf{Durable :} Plusieurs années, souvent à vie (grâce à la mémoire immunitaire).
    \item \textbf{Spécifique :} Ne protège que contre l'agent ciblé (ou ses variants proches).
\end{itemize}
\end{propriete}

\courssub{Schéma du principe de la vaccination}

\begin{popfigure}
\begin{tikzpicture}[
    node distance=1.1cm and 1.4cm,
    box/.style={rectangle, draw=popdark, thick, fill=popcream, rounded corners, minimum height=1.1cm, text width=3.0cm, align=center, font=\small},
    arrow/.style={-Stealth, thick, popblue},
    arrowr/.style={-Stealth, thick, poporange},
]
% Nodes
\node[box, align=center] (vaccin){Vaccin\\ (Ag atténué/tué/fragment)};
\node[box, right=of vaccin, align=center] (apc){Cellule Présentatrice\\ d'Antigène (APC)\\(Macrophage, Dendritique)};
\node[box, above right=of apc, align=center] (lymphT){Lymphocyte T4\\ (Coordinateur)};
\node[box, below right=of apc, align=center] (lymphB){Lymphocyte B\\ (Récepteur Ag spécifique)};
\node[box, right=of lymphT, align=center] (cytok){Cytokines\\ (IL-2, IL-4...)};
\node[box, right=of lymphB, align=center] (plasma){Plasmocyte\\ (Usine à Anticorps)};
\node[box, below=of plasma, align=center] (memB){Lymphocyte B\\ Mémoire};
\node[box, above=of memB, align=center] (memT){Lymphocyte T4\\ Mémoire};
\node[box, right=of plasma, fill=popgreenL, text width=3.5cm] (protect) {PROTECTION DURABLE\\ Anticorps circulants\\ + Cellules Mémoire (LBm, LTm)};

% Arrows
\draw[arrow] (vaccin) -- node[above, font=\tiny, sloped] {Phagocytose \& Présentation (MHC II)} (apc);
\draw[arrow] (apc) -- node[above left, font=\tiny] {Présentation Ag + MHC II} (lymphT);
\draw[arrow] (apc) -- node[below left, font=\tiny] {Reconnaissance Ag natif} (lymphB);
\draw[arrow] (lymphT) -- node[above, font=\tiny] {Activation (Cytokines)} (lymphB);
\draw[arrow] (lymphB) -- node[above, font=\tiny] {Différenciation} (plasma);
\draw[arrow] (lymphB) -- node[right, font=\tiny] {Différenciation} (memB);
\draw[arrow] (lymphT) -- node[left, font=\tiny] {Différenciation} (memT);
\draw[arrowr] (plasma) -- node[above, font=\tiny] {Sécrétion massive d'Ac} (protect);
\draw[arrowr] (memB) -- node[right, font=\tiny] {Persistance : années} (protect);
\draw[arrowr] (memT) -- node[right, font=\tiny] {Persistance : années} (protect);

% Legend
\node[below=0.6cm of vaccin, font=\footnotesize, text width=13cm, align=center, popmuted] {\textbf{Légende :} 1. Injection du vaccin. 2. Prise en charge et présentation de l'antigène par les APC. 3. Activation des lymphocytes T4 et B spécifiques (sélection clonale). 4. Clonage et différenciation en plasmocytes (producteurs d'anticorps) et cellules mémoires. 5. Installation de l'immunité active durable.};
\end{tikzpicture}
\legende{Principes de la vaccinothérapie : le vaccin déclenche une réponse immunitaire adaptative aboutissant à la production d'anticorps et à la mise en place d'une mémoire immunologique.}
\end{popfigure}

\courssub{La réponse immunitaire lors de la vaccination : primo-vaccination et rappel}

Lorsqu'on vaccine un individu pour la première fois (primo-vaccination), la réponse immunitaire suit la cinétique d'une \textbf{réponse primaire}. Si on réinjecte le même vaccin plus tard (rappel), on observe une \textbf{réponse secondaire} (ou anamnestique), bien plus performante.

\begin{experience}[Étude de la cinétique des anticorps après vaccination (Document d'analyse)]
On suit le taux d'anticorps sériques (titre sérique) chez un individu vacciné contre le tétanos (anatoxine tétanique).
\begin{itemize}
    \item \textbf{Jour 0 :} Injection 1 (Primo-vaccination).
    \item \textbf{Jour 30 :} Injection 2 (Rappel).
    \item Mesure régulière du taux d'anticorps (en UI/mL, échelle logarithmique).
\end{itemize}
\end{experience}

\begin{popfigure}
\begin{tikzpicture}
\begin{axis}[
    width=\linewidth,
    height=10cm,
    xlabel={Temps (jours)},
    ylabel={Taux d'anticorps (UI/mL) -- échelle logarithmique},
    ymode=log,
    ymin=0.1, ymax=10000,
    xmin=0, xmax=120,
    xtick={0,14,30,44,60,90,120},
    ytick={0.1,1,10,100,1000,10000},
    grid=major,
    grid style={popline},
    legend pos=north west,
    legend style={fill=popcream, draw=popdark, font=\small},
    tick label style={font=\small},
    label style={font=\small},
    title style={font=\normalsize},
    title={Cinétique de la réponse anticorps : Primo-vaccination et Rappel},
    clip=false
]
% Primary response curve (approximate logistic growth + decay)
\addplot[popcurve, popblue, very thick, domain=0:30, samples=200] 
    ({x}, {0.1 + 800/(1+exp(-0.6*(x-12))) * exp(-0.03*(x-12))});
\addlegendentry{Réponse primaire (1\textsuperscript{re} injection)}

% Secondary response curve (booster at day 30)
\addplot[popcurve, poporange, very thick, domain=30:120, samples=200] 
    ({x}, {100 + 8000/(1+exp(-1.5*(x-35))) * exp(-0.01*(x-35))});
\addlegendentry{Réponse secondaire (Rappel)}

% Annotations
\draw[popblue, dashed, thick] (12,0.1) -- (12,800) node[above, pos=0.9, font=\tiny, popblue] {Pic primaire $\approx$ J12-14};
\draw[poporange, dashed, thick] (35,0.1) -- (35,8000) node[above, pos=0.95, font=\tiny, poporange] {Pic secondaire $\approx$ J5 post-rappel};
\draw[popdark, dashed, thick] (30,0.1) -- (30,10000) node[above, font=\tiny, popdark, rotate=90, anchor=south] {Injection de rappel (J30)};
\end{axis}
\end{tikzpicture}
\legende{Évolution du taux d'anticorps après primo-vaccination (courbe bleue) et après rappel (courbe orange). La réponse secondaire est plus précoce, plus intense (pic 10 à 100 fois plus haut) et plus durable. Échelle logarithmique sur l'ordonnée.}
\end{popfigure}

\begin{methode}[Analyser une courbe de vaccination (réponse primaire vs secondaire)]
Pour exploiter un document montrant l'évolution du taux d'anticorps après vaccination :
\begin{enumerate}
    \item \textbf{Identifier les injections :} Repérer les flèches ou dates des injections (J0, J30...).
    \item \textbf{Caractériser la réponse primaire (après 1\textsuperscript{re} inj.) :} Noter le \textbf{délai de latence} (quelques jours sans anticorps détectables), la \textbf{vitesse de montée}, le \textbf{taux maximal (pic)} et la \textbf{durée de vie} des anticorps (baisse progressive).
    \item \textbf{Caractériser la réponse secondaire (après rappel) :} Noter l'absence (ou brièveté) de la phase de latence, la \textbf{vitesse de montée beaucoup plus grande}, le \textbf{taux maximal 10 à 100 fois plus élevé}, et la \textbf{pente de descente plus faible} (persistance longue).
    \item \textbf{Conclure :} Le rappel \textbf{réactive les cellules mémoires} $\rightarrow$ réponse plus rapide, plus forte, plus durable = \textbf{protection efficace}.
\end{enumerate}
\end{methode}

\begin{exoresolu}[Analyse d'une courbe de vaccination anti-hépatite B]
\textbf{Énoncé :} Le document ci-dessus (courbe) montre l'évolution du taux d'anti-HBs (anticorps anti-hépatite B) chez un infirmier vacciné selon le schéma 0, 1, 6 mois. On considère ici les deux premières injections (M0 et M1).
\begin{enumerate}
    \item Préciser la nature de l'immunité acquise (active/passive ? naturelle/artificielle ?).
    \item Comparer la cinétique des deux pics d'anticorps (vitesse, intensité).
    \item Expliquer la base biologique de la différence observée.
    \item Pourquoi une 3\textsuperscript{e} injection à M6 est-elle recommandée ?
\end{enumerate}
\tcblower
\textbf{Solution détaillée :}
\begin{enumerate}
    \item C'est une \textbf{immunité active artificielle} : l'organisme fabrique ses propres anticorps suite à l'injection d'un antigène (vaccin recombinant HBsAg).
    \item \textbf{Pic 1 (M1, soit $\sim$J30) :} Montée lente (détection $\sim$J14), pic modéré ($\sim$ 100-500 UI/mL), puis baisse. \textbf{Pic 2 (M2, rappel à M1) :} Montée très rapide (détection $\sim$J3-5), pic très élevé ($\sim$ 5000-10000 UI/mL), baisse très lente.
    \item La différence s'explique par la présence de \textbf{lymphocytes B mémoires} et de \textbf{lymphocytes T4 mémoires} générés lors de la primo-vaccination. Lors du rappel, ces cellules sont déjà présentes en grand nombre, spécifiques de l'antigène, et ne nécessitent plus la phase lente de sélection clonale et de différenciation initiale. Elles se transforment immédiatement en plasmocytes hyper-producteurs.
    \item La 3\textsuperscript{e} injection (rappel à 6 mois) vise à \textbf{stabiliser le taux d'anticorps à un niveau protecteur très élevé (> 100 UI/mL)} sur le très long terme (immunité > 20 ans) et à augmenter le pool de cellules mémoires à longue durée de vie (cellules plasmatiques médullaires).
\end{enumerate}
\end{exoresolu}

\begin{attention}[Piège majeur : Vaccin $\neq$ Médicament curatif]
\textbf{Erreur fréquente :} « On vaccine un malade pour le guérir. »
\textbf{Réalité :} Le vaccin est \textbf{strictement préventif}. Il faut 2 à 3 semaines pour que l'immunité active s'installe. Si le patient est \emph{déjà} incubant la maladie ou malade, le vaccin n'aura \textbf{aucun effet thérapeutique} (sauf exceptions rares comme la rage post-exposition où le vaccin + immunoglobulines agissent pendant la longue incubation). Pour traiter une maladie en cours, on utilise des \textbf{antibiotiques} (bactéries) ou des \textbf{antiviraux}, voire de la \textbf{sérothérapie} (anticorps tout faits) pour une action immédiate.
\end{attention}

\courssub{Principaux vaccins du calendrier vaccinal camerounais (PEV)}

Le Cameroun suit les recommandations de l'OMS adaptées à son épidémiologie. Voici les vaccins clés que tu dois connaître pour le BEPC :

\begin{exemplebox}[Les vaccins du PEV au Cameroun (Rappel utile pour le BEPC)]
\begin{itemize}
    \item \textbf{BCG (Bacille de Calmette et Guérin) :} Vivant atténué (\emph{Mycobacterium bovis} atténué). Protège contre les formes graves de la \textbf{tuberculose} (méningite, miliaire). Injecté à la naissance (voie intradermique) $\rightarrow$ laisse une petite cicatrice caractéristique.
    \item \textbf{Vaccin Antiamaril (Fièvre Jaune) :} Vivant atténué (souche 17D). \textbf{Obligatoire} pour l'entrée/sortie du territoire (Certificat international de vaccination). Une dose $\rightarrow$ immunité à vie (depuis 2014, OMS).
    \item \textbf{DTC (Diphtérie - Tétanos - Coqueluche) :} Vaccin combiné. Anatoxines diphtérique et tétanique + coqueluche (acellulaire ou entière). 3 doses (6, 10, 14 sem) + rappels. Protège contre le tétanos (toxine tétanique), la diphtérie (toxine diphtérique), la coqueluche (\emph{Bordetella pertussis}).
    \item \textbf{VPO (Vaccin Polio Oral) / VIP (Vaccin Polio Inactivé) :} VPO = virus vivants atténués (Sabin), administration orale (gouttes), induit une immunité digestive (IgA) $\rightarrow$ bloque transmission. VIP = virus tués (Salk), injection, pas de risque de VAPP (polio vaccinale). Stratégie actuelle : séquentiel VIP puis VPO ou VIP seul.
    \item \textbf{Rougeole / ROR (Rougeole-Oreillons-Rubéole) :} Vivants atténués. 2 doses (9 mois et 15-18 mois).
    \item \textbf{Hépatite B :} Vaccin recombinant (protéine HBsAg produite par levure). 3 doses (naissance, 6 sem, 14 sem). Prévention du cancer du foie (carcinome hépatocellulaire).
    \item \textbf{Autres :} \emph{Haemophilus influenzae} type b (Hib), Pneumocoque (PCV), Rotavirus, HPV (filles 9-14 ans, cancer col utérus), Méningocoque A (MenAfriVac, ceinture de la méningite).
\end{itemize}
\end{exemplebox}

\begin{aretenir}[Synthèse : La Vaccinothérapie]
\begin{itemize}
    \item \textbf{But :} Prévenir la maladie chez le sujet sain.
    \item \textbf{Agent :} Vaccin = Ag tué, atténué ou fragment (non pathogène, immunogène).
    \item \textbf{Mécanisme :} Déclenche une \textbf{immunité active} : l'organisme produit \textbf{ses propres anticorps} et des \textbf{cellules mémoires (LBm, LTm)}.
    \item \textbf{Cinétique :} Primo-vaccination = réponse primaire (lente, faible, courte). Rappel = réponse secondaire (rapide, forte, durable).
    \item \textbf{Intérêt majeur :} Protection individuelle durable + \textbf{immunité collective} (arrêt de la circulation du pathogène).
    \item \textbf{Distinction clé :} Vaccinothérapie (préventif, actif, lent) $\neq$ Sérothérapie (curatif/urgence, passif, immédiat, court).
\end{itemize}
\end{aretenir}

\coursec{IV. La séro-vaccinothérapie}

Nous venons de voir que la \cle{vaccinothérapie} protège durablement l'organisme, mais qu'elle demande du temps : il faut plusieurs jours, parfois plusieurs semaines, pour que le système immunitaire fabrique ses propres anticorps après l'injection du vaccin. Que faire alors lorsque le danger est \emph{immédiat} ? C'est toute la question de cette dernière partie.

\courssub{Situation-problème : la morsure par un chien suspect de rage}

\begin{exemplebox}[Une urgence au Centre Pasteur du Cameroun]
Un enfant de 12 ans, habitant un quartier de Yaoundé, est mordu à la jambe par un chien errant qui présente un comportement anormal (agressivité, écume à la gueule) : le chien est \cle{suspect de rage}. L'enfant n'a \textbf{jamais été vacciné} contre la rage. Ses parents le conduisent en urgence au Centre Pasteur du Cameroun.
\end{exemplebox}

La rage est une maladie mortelle dès que les premiers signes apparaissent. Le médecin se trouve devant un dilemme :

\begin{itemize}
\item S'il injecte \textbf{uniquement le vaccin} antirabique : l'organisme de l'enfant mettra environ 2 à 4 semaines pour fabriquer suffisamment d'anticorps. Or le virus de la rage peut atteindre le système nerveux avant ce délai. \textbf{Le vaccin seul serait trop lent.}
\item S'il injecte \textbf{uniquement le sérum} antirabique (anticorps prêts à l'emploi) : l'enfant est protégé immédiatement, mais ces anticorps étrangers sont détruits en quelques semaines. \textbf{Le sérum seul ne protège pas durablement.}
\end{itemize}

La solution adoptée consiste à injecter \textbf{les deux en même temps} : c'est la \cle{séro-vaccinothérapie}.

\courssub{Définition}

\begin{definition}[Séro-vaccinothérapie]
La \cle{séro-vaccinothérapie} est une méthode d'aide au système immunitaire qui consiste à injecter \textbf{simultanément (ou successivement)} à un individu :
\begin{itemize}
\item un \cle{sérum}, c'est-à-dire des \textbf{anticorps prêts à l'emploi} (sérothérapie), qui assurent une protection \textbf{immédiate} ;
\item un \cle{vaccin}, c'est-à-dire des \textbf{antigènes} du microbe (vaccinothérapie), qui déclenchent la fabrication d'anticorps par l'organisme lui-même, assurant une protection \textbf{durable}.
\end{itemize}
\end{definition}

\begin{aretenir}[L'idée essentielle]
La séro-vaccinothérapie associe les \textbf{avantages} des deux méthodes précédentes et supprime leurs \textbf{inconvénients} :
\begin{itemize}
\item le sérum = protection \textbf{immédiate} mais \textbf{temporaire} (\cle{immunité passive}) ;
\item le vaccin = protection \textbf{lente} à s'installer mais \textbf{durable} (\cle{immunité active}).
\end{itemize}
Le sérum « fait le pont » pendant que le vaccin « prépare l'avenir ».
\end{aretenir}

\courssub{Principe : deux protections qui se relaient}

Le principe de la séro-vaccinothérapie repose sur la \textbf{complémentarité dans le temps} des deux types d'immunité :

\begin{enumerate}
\item \textbf{Juste après l'injection} : les anticorps du sérum sont déjà présents en grande quantité dans le sang. Ils neutralisent immédiatement le microbe ou sa toxine. C'est l'\cle{immunité passive} : l'organisme reçoit des anticorps tout faits, il n'a rien fabriqué.
\item \textbf{Au cours des semaines suivantes} : les anticorps du sérum sont progressivement détruits et éliminés (leur taux diminue). Mais pendant ce temps, le vaccin a stimulé les lymphocytes de l'organisme, qui fabriquent leurs propres anticorps : le taux d'anticorps « maison » augmente. C'est l'\cle{immunité active}.
\item \textbf{À long terme} : quand les anticorps du sérum ont disparu, ceux fabriqués grâce au vaccin sont là pour prendre le relais, et l'organisme garde une \textbf{mémoire immunitaire} durable.
\end{enumerate}

\courssub{Justification à partir des courbes}

Pour comprendre rigoureusement pourquoi l'association sérum + vaccin est efficace, étudions l'évolution du \textbf{taux d'anticorps} dans le sang en fonction du temps, dans les deux cas.

\begin{popfigure}
\begin{center}
\begin{tikzpicture}
\begin{axis}[
width=0.95\linewidth,
height=7cm,
xlabel={Temps (semaines)},
ylabel={Taux d'anticorps (unités arbitraires)},
xmin=0, xmax=10,
ymin=0, ymax=11,
xtick={0,2,4,6,8,10},
ytick={0,2,4,6,8,10},
legend style={at={(0.97,0.55)},anchor=east,font=\small},
grid=both,
grid style={popline!40},
axis line style={popdark},
]
% Anticorps du sérum : décroissance exponentielle
\addplot[popcurve,poporange,very thick,domain=0:10,samples=100]{9*exp(-0.45*x)};
\addlegendentry{Anticorps apportés par le sérum}
% Anticorps produits grâce au vaccin : montée sigmoïdale simplifiée
\addplot[popcurve,popblue,very thick,domain=0:10,samples=100]{8*(1-exp(-0.5*(x-1)))*(x>1)};
\addlegendentry{Anticorps fabriqués grâce au vaccin}
\end{axis}
\end{tikzpicture}
\end{center}
\legende{Évolution comparée du taux d'anticorps dans le sang après une séro-vaccinothérapie. En orange : les anticorps du sérum, abondants dès l'injection, diminuent rapidement et deviennent négligeables en quelques semaines. En bleu : les anticorps fabriqués par l'organisme grâce au vaccin apparaissent lentement puis se maintiennent durablement. Les deux courbes \textbf{se complètent} : à tout moment, le taux total d'anticorps reste suffisant pour protéger l'individu.}
\end{popfigure}

\textbf{Interprétation de la courbe.} On observe que :

\begin{itemize}
\item pendant les \textbf{deux premières semaines}, seuls les anticorps du sérum (courbe orange) assurent la protection, car le vaccin n'a pas encore provoqué de réponse détectable ;
\item vers la \textbf{3\textsuperscript{e}--4\textsuperscript{e} semaine}, les anticorps du sérum chutent fortement, mais ceux produits grâce au vaccin (courbe bleue) sont déjà abondants : \textbf{le relais est assuré} ;
\item à \textbf{aucun moment} le taux total d'anticorps ne tombe à zéro : la protection est \textbf{continue}.
\end{itemize}

\begin{methode}[Justifier une séro-vaccinothérapie]
Pour justifier scientifiquement le recours à une séro-vaccinothérapie, suis ces étapes :
\begin{enumerate}
\item Identifier la situation : contamination récente chez un sujet \textbf{non immunisé} (non vacciné) et maladie à évolution rapide.
\item Montrer que le \textbf{vaccin seul est insuffisant} : le délai de fabrication des anticorps (2 à 4 semaines) est supérieur au délai d'apparition de la maladie.
\item Montrer que le \textbf{sérum seul est insuffisant} : ses anticorps sont détruits en quelques semaines, sans mémoire immunitaire.
\item Conclure : le sérum \textbf{couvre le délai} nécessaire à l'organisme pour répondre au vaccin ; la protection est donc à la fois \textbf{immédiate et durable}.
\end{enumerate}
\end{methode}

\begin{attention}[Ne pas mélanger sérum et vaccin !]
Erreur fréquente : croire qu'on peut injecter le sérum et le vaccin \textbf{au même endroit} ou les mélanger dans la même seringue. C'est à proscrire : les anticorps du sérum risqueraient de \textbf{neutraliser les antigènes du vaccin} avant que l'organisme n'ait pu réagir, et le vaccin deviendrait inefficace. En pratique, on injecte le sérum et le vaccin \textbf{en deux points différents} du corps (par exemple un dans chaque bras, ou le sérum autour de la plaie et le vaccin dans le bras).
\end{attention}

\courssub{Exemples d'application}

\begin{exemplebox}[La séro-vaccination antirabique au Cameroun]
Après une morsure par un chien suspect de rage à Yaoundé, le Centre Pasteur du Cameroun pratique en urgence la \textbf{séro-vaccination antirabique} :
\begin{itemize}
\item injection du \textbf{sérum antirabique} (anticorps antirabiques prêts à l'emploi), de préférence autour de la plaie, pour neutraliser immédiatement le virus ;
\item injection du \textbf{vaccin antirabique} dans le bras, renouvelée selon un calendrier précis (plusieurs doses sur environ un mois), pour que l'organisme fabrique ses propres anticorps.
\end{itemize}
Cette méthode, appliquée à temps, sauve la vie des personnes mordues.
\end{exemplebox}

\begin{exemplebox}[La séro-vaccination antitétanique]
Un cultivateur se blesse profondément au pied avec une houe rouillée (\textbf{plaie souillée} de terre). Il n'a jamais été vacciné contre le tétanos. Le centre de santé lui administre :
\begin{itemize}
\item le \textbf{sérum antitétanique}, qui neutralise immédiatement la toxine tétanique déjà présente ;
\item le \textbf{vaccin antitétanique}, qui déclenche la fabrication durable d'anticorps et la mémoire immunitaire.
\end{itemize}
\end{exemplebox}

\begin{popfigure}
\begin{center}
\begin{tikzpicture}[
node distance=0.6cm,
box/.style={rectangle,rounded corners=4pt,draw=popdark,thick,align=center,minimum height=1cm,font=\small},
fleche/.style={-{Stealth[length=3mm]},thick,popdark}
]
\node[box,fill=poppinkL,text width=3.2cm] (contam) {\textbf{Contamination}\\ morsure de chien\\ enragé, plaie souillée};
\node[box,fill=poporangeL,text width=3.4cm,right=1.1cm of contam] (inj) {\textbf{Injection}\\ sérum + vaccin\\ (en 2 points différents)};
\node[box,fill=popgoldL,text width=3.4cm,right=1.1cm of inj,yshift=1.3cm] (passive) {\textbf{Protection immédiate}\\ anticorps du sérum\\ (immunité passive)};
\node[box,fill=popgreenL,text width=3.4cm,right=1.1cm of inj,yshift=-1.3cm] (active) {\textbf{Protection durable}\\ anticorps fabriqués\\ grâce au vaccin\\ (immunité active)};
\draw[fleche] (contam) -- (inj);
\draw[fleche] (inj.east) -- ++(0.35,0) |- (passive.west);
\draw[fleche] (inj.east) -- ++(0.35,0) |- (active.west);
\end{tikzpicture}
\end{center}
\legende{Schéma récapitulatif de la séro-vaccinothérapie : après une contamination, l'injection simultanée du sérum et du vaccin assure une double protection — immédiate grâce aux anticorps du sérum, durable grâce aux anticorps fabriqués par l'organisme en réponse au vaccin.}
\end{popfigure}

\courssub{Condition d'efficacité}

\begin{propriete}[Agir le plus tôt possible]
La séro-vaccinothérapie n'est efficace que si l'injection du sérum et du vaccin est réalisée \textbf{le plus tôt possible après la contamination}, avant que le microbe ou sa toxine ne se soit propagé dans l'organisme. Passé un certain délai (apparition des premiers signes de la maladie), il est souvent \textbf{trop tard} : c'est pourquoi toute morsure d'animal suspect ou toute plaie profonde souillée constitue une \textbf{urgence médicale}.
\end{propriete}

\begin{savaistu}[La rage au Cameroun]
La rage existe encore au Cameroun, transmise surtout par les chiens errants. Chaque année, des milliers de personnes mordues reçoivent la séro-vaccination antirabique dans des centres spécialisés comme le Centre Pasteur du Cameroun à Yaoundé. La rage est une maladie \textbf{mortelle à près de 100\,\%} une fois les signes apparus, mais \textbf{évitable à 100\,\%} si la séro-vaccinothérapie est pratiquée rapidement après la morsure. La vaccination des chiens reste le moyen le plus efficace de lutter contre cette maladie.
\end{savaistu}

\begin{exoresolu}[Justifier un protocole d'urgence]
\textbf{Énoncé.} Un élève de 3\textsuperscript{e} se fait mordre par un chien errant présentant des signes de rage. Non vacciné contre la rage, il est conduit au centre de santé où le médecin lui injecte à la fois du sérum antirabique et du vaccin antirabique. Un camarade s'étonne : « Pourquoi deux injections ? Le vaccin ne suffit-il pas ? » Rédige une réponse argumentée pour expliquer ce protocole.
\tcblower
\textbf{Solution détaillée.}
\begin{enumerate}
\item \textbf{Le vaccin seul est insuffisant} : après l'injection du vaccin, l'organisme met environ 2 à 4 semaines pour fabriquer une quantité suffisante d'anticorps. Or le virus de la rage évolue vite et peut atteindre le système nerveux avant ce délai : la maladie serait alors mortelle.
\item \textbf{Le sérum seul est insuffisant} : les anticorps du sérum agissent immédiatement, mais ils sont détruits en quelques semaines et ne laissent aucune mémoire immunitaire : la protection disparaîtrait.
\item \textbf{L'association des deux est la solution} : le sérum antirabique fournit des anticorps prêts à l'emploi qui neutralisent le virus \textbf{immédiatement} (immunité passive), pendant que le vaccin antirabique stimule l'organisme qui fabrique ses \textbf{propres anticorps} de façon durable (immunité active).
\item \textbf{Conclusion} : la séro-vaccinothérapie assure une protection à la fois immédiate et durable, à condition d'être pratiquée le plus tôt possible après la morsure, en injectant le sérum et le vaccin en deux points différents du corps.
\end{enumerate}
\end{exoresolu}

\begin{exercice}[À toi de jouer]
Un agriculteur de l'ouest du Cameroun reçoit une piqûre profonde d'un clou rouillé en travaillant dans son champ de maïs. Il n'a jamais reçu le vaccin antitétanique.
\begin{enumerate}
\item Quel traitement le personnel du centre de santé doit-il lui administrer ? Justifie.
\item Pourquoi faut-il agir rapidement ?
\item Rappelle pourquoi le sérum et le vaccin ne doivent pas être injectés au même endroit.
\end{enumerate}
\end{exercice}

\begin{corrige}[Exercice « À toi de jouer »]
\begin{enumerate}
\item On doit pratiquer une \textbf{séro-vaccinothérapie antitétanique} : injection du sérum antitétanique (anticorps prêts à l'emploi qui neutralisent immédiatement la toxine tétanique) et du vaccin antitétanique (qui déclenche la fabrication durable d'anticorps par l'organisme). Le vaccin seul serait trop lent et le sérum seul ne protégerait pas durablement.
\item Il faut agir rapidement car la séro-vaccinothérapie n'est efficace que si elle est réalisée \textbf{avant que la toxine tétanique ne se propage} dans l'organisme ; une fois les signes du tétanos apparus, il est souvent trop tard.
\item Le sérum et le vaccin ne doivent pas être injectés au même endroit (ni mélangés) car les anticorps du sérum neutraliseraient les antigènes du vaccin, rendant celui-ci inefficace.
\end{enumerate}
\end{corrige}

\coursec{V. Comparaison et synthèse des méthodes d'aide à l'immunité}

Nous venons d'étudier la \cle{séro-vaccinothérapie}, qui associe un sérum et un vaccin pour protéger à la fois immédiatement et durablement une personne exposée à un danger infectieux, comme dans le cas du tétanos. Nous disposons donc maintenant de quatre méthodes d'aide à l'organisme face aux maladies : l'\cle{antibiothérapie}, la \cle{sérothérapie}, la \cle{vaccinothérapie} et la \cle{séro-vaccinothérapie}. Mais face à une situation concrète — un enfant souffrant d'une angine bactérienne, un élève mordu par un chien errant, une campagne de vaccination contre la rougeole au Cameroun — laquelle choisir ? C'est toute l'utilité de cette synthèse : apprendre à comparer ces méthodes et à raisonner comme un médecin.

\courssub{Tableau comparatif des quatre méthodes}

Pour comparer rigoureusement, il faut d'abord se poser les bonnes questions, celles qu'un médecin se pose devant son patient :

\begin{itemize}
\item \textbf{Quelle substance} administre-t-on ? Un produit chimique, des anticorps, des antigènes, ou les deux ?
\item \textbf{Sur quoi agit-elle} ? Directement sur le microbe, ou en aidant le système immunitaire ?
\item \textbf{Agit-elle vite} ? La protection est-elle immédiate ou faut-il attendre ?
\item \textbf{Combien de temps} dure la protection ?
\item \textbf{S'agit-il d'une immunité active} (l'organisme fabrique lui-même ses anticorps) \textbf{ou passive} (on reçoit des anticorps tout faits) ?
\item \textbf{Est-ce un traitement curatif} (on soigne une maladie déclarée) \textbf{ou préventif} (on empêche la maladie d'apparaître) ?
\end{itemize}

\begin{aretenir}[Les quatre méthodes en un coup d'œil]
\begin{center}
\rowcolors{2}{popblueL}{white}
\begin{tabularx}{\linewidth}{|l|X|X|X|X|}
\hline
\rowcolor{popblue!40} \textbf{Critère} & \textbf{Antibiothérapie} & \textbf{Sérothérapie} & \textbf{Vaccinothérapie} & \textbf{Séro-vaccinothérapie} \\
\hline
Substance injectée & Antibiotique (substance chimique) & Sérum contenant des anticorps & Vaccin contenant des antigènes & Sérum (anticorps) + vaccin (antigènes) \\
\hline
Mécanisme & Détruit ou bloque la bactérie ; n'agit pas sur le système immunitaire & Les anticorps reçus neutralisent immédiatement l'antigène & L'organisme fabrique lui-même ses anticorps et des cellules mémoires & Les anticorps agissent tout de suite, puis le vaccin déclenche la fabrication d'anticorps \\
\hline
Rapidité d'action & Rapide (quelques jours) & Immédiate & Lente (2 à 3~semaines) & Immédiate puis durable \\
\hline
Durée de protection & Aucune après le traitement & Courte (environ 3~semaines) & Longue (plusieurs années) & Courte puis longue \\
\hline
Type d'immunité & Aucune (n'immunise pas) & Passive & Active & Passive puis active \\
\hline
Caractère & Curatif & Curatif, d'urgence & Préventif & Préventif d'urgence \\
\hline
Exemples & Angine, infections urinaires & Morsure de serpent, rage & BCG, antitétanique, antirougeoleux & Tétanos (plaie souillée), rage \\
\hline
\end{tabularx}
\end{center}
\end{aretenir}

\begin{popfigure}
\begin{tikzpicture}[
  box/.style={draw=popdark, rounded corners=2pt, fill=popblueL, minimum width=3.1cm, minimum height=0.85cm, align=center, font=\small},
  head/.style={draw=popdark, rounded corners=2pt, fill=popblue!55, minimum width=3.1cm, minimum height=0.85cm, align=center, font=\small\bfseries},
  crit/.style={draw=popdark, rounded corners=2pt, fill=popgoldL, minimum width=2.6cm, minimum height=0.85cm, align=center, font=\small\bfseries},
]
% En-têtes de colonnes
\node[crit] at (0,0) {Critère};
\node[head, align=center] at (3.5,0){Antibio-\\thérapie};
\node[head, align=center] at (7,0){Séro-\\thérapie};
\node[head, align=center] at (10.5,0){Vaccino-\\thérapie};
\node[head, align=center] at (14,0){Séro-vaccino-\\thérapie};
% Lignes
\node[crit] at (0,-1.2) {Substance};
\node[box, align=center] at (3.5,-1.2){Produit\\chimique};
\node[box, fill=popgreenL] at (7,-1.2) {Anticorps};
\node[box, fill=poporangeL] at (10.5,-1.2) {Antigènes};
\node[box, fill=poppinkL, align=center] at (14,-1.2){Anticorps\\+ antigènes};

\node[crit] at (0,-2.4) {Rapidité};
\node[box] at (3.5,-2.4) {Rapide};
\node[box, fill=popgreenL] at (7,-2.4) {Immédiate};
\node[box, fill=poporangeL] at (10.5,-2.4) {Lente};
\node[box, fill=poppinkL] at (14,-2.4) {Immédiate};

\node[crit] at (0,-3.6) {Durée};
\node[box] at (3.5,-3.6) {Nulle};
\node[box, fill=popgreenL, align=center] at (7,-3.6){Courte\\($\approx$ 3 sem.)};
\node[box, fill=poporangeL, align=center] at (10.5,-3.6){Longue\\(années)};
\node[box, fill=poppinkL, align=center] at (14,-3.6){Courte puis\\longue};

\node[crit] at (0,-4.8) {Immunité};
\node[box] at (3.5,-4.8) {Aucune};
\node[box, fill=popgreenL] at (7,-4.8) {Passive};
\node[box, fill=poporangeL] at (10.5,-4.8) {Active};
\node[box, fill=poppinkL, align=center] at (14,-4.8){Passive puis\\active};

\node[crit] at (0,-6) {Rôle};
\node[box] at (3.5,-6) {Curatif};
\node[box, fill=popgreenL, align=center] at (7,-6){Curatif\\d'urgence};
\node[box, fill=poporangeL] at (10.5,-6) {Préventif};
\node[box, fill=poppinkL, align=center] at (14,-6){Préventif\\d'urgence};
\end{tikzpicture}
\legende{Tableau comparatif à double entrée des quatre méthodes d'aide à l'immunité. Chaque colonne correspond à une méthode, chaque ligne à un critère de comparaison.}
\end{popfigure}

\courssub{Ce qu'il faut retenir de chaque méthode}

\textbf{L'antibiothérapie : le seul traitement qui n'immunise pas.}\par
L'antibiotique est une \cle{substance chimique} qui agit directement sur la \cle{bactérie} : il la tue ou l'empêche de se multiplier. Point essentiel : il \textbf{n'agit pas sur le système immunitaire} et ne laisse \textbf{aucune protection} après le traitement. C'est un traitement purement \cle{curatif}, utilisé quand une maladie bactérienne est déjà déclarée (angine à streptocoque, infection urinaire).

\textbf{La sérothérapie : le secours immédiat mais éphémère.}\par
On injecte un sérum contenant des \cle{anticorps} tout prêts. L'action est \cle{immédiate} : c'est le traitement d'\cle{urgence} (morsure de serpent, suspicion de rage). Mais ces anticorps étrangers sont progressivement détruits par l'organisme : la protection est \cle{courte} (environ 3~semaines). C'est une immunité \cle{passive}, car l'organisme n'a rien fabriqué lui-même.

\textbf{La vaccinothérapie : la protection lente mais durable.}\par
On injecte un vaccin contenant des \cle{antigènes} (microbe tué, atténué ou fragment de microbe). L'organisme réagit en fabriquant \textbf{lui-même} ses anticorps et surtout des \cle{cellules mémoires}. C'est une immunité \cle{active}, lente à s'installer (2 à 3~semaines) mais \cle{durable} (plusieurs années). Son rôle est \cle{préventif} : on vaccine avant la maladie.

\textbf{La séro-vaccinothérapie : le meilleur des deux mondes en urgence.}\par
On associe sérum et vaccin : les anticorps du sérum protègent \textbf{immédiatement} (immunité passive) pendant que le vaccin déclenche la fabrication d'anticorps par l'organisme (immunité active durable). C'est la méthode employée lors d'une plaie souillée de terre chez une personne non vaccinée contre le tétanos.

\begin{propriete}[Synthèse fondamentale]
Seule la \cle{vaccination} confère une immunité \textbf{active et durable}, car elle pousse l'organisme à fabriquer ses propres anticorps et des cellules mémoires. La \cle{sérothérapie} ne fait que \emph{prêter} des anticorps temporairement : elle protège vite mais pas longtemps. L'\cle{antibiothérapie}, elle, ne confère aucune immunité : elle détruit directement la bactérie.
\end{propriete}

\begin{exemplebox}[Des durées de protection très différentes]
La protection conférée par un sérum (sérothérapie) dure environ \textbf{3~semaines}, le temps que les anticorps injectés soient éliminés. En revanche, la protection conférée par le \textbf{vaccin antitétanique} dure environ \textbf{10~ans} : c'est pourquoi un \cle{rappel} est recommandé tous les 10~ans pour entretenir la mémoire immunitaire. Ce contraste illustre parfaitement la différence entre immunité passive (courte) et immunité active (longue).
\end{exemplebox}

\courssub{Comment choisir la bonne méthode ? L'arbre de décision}

Face à un patient, le médecin suit un raisonnement logique, que l'on peut résumer par trois questions :

\begin{enumerate}
\item \textbf{La maladie est-elle déclarée ou y a-t-il seulement un risque ?} (curatif ou préventif)
\item \textbf{Y a-t-il urgence ?} (faut-il une protection immédiate ?)
\item \textbf{Quelle est la nature de l'agent : bactérie ou virus ? Le sujet est-il déjà vacciné ?}
\end{enumerate}

\begin{popfigure}
\begin{tikzpicture}[
  question/.style={draw=popdark, fill=popgoldL, rounded corners=3pt, align=center, font=\small\bfseries, minimum width=4.5cm, minimum height=1cm},
  decision/.style={draw=popdark, fill=popgreenL, rounded corners=3pt, align=center, font=\small, minimum width=3.8cm, minimum height=0.9cm},
  fleche/.style={-{Stealth[length=2.5mm]}, thick, popdark},
  lbl/.style={font=\small\itshape, popdark}
]
\node[question, align=center] (q1) at (0,0){Maladie bactérienne\\déclarée ?};
\node[decision, align=center] (antibio) at (-5.5,-2){\textbf{Antibiothérapie}\\(ex. angine bactérienne)};
\node[question, align=center] (q2) at (3.5,-2){Risque d'infection grave\\ou maladie virale ?};
\node[question, align=center] (q3) at (0.5,-4.2){Urgence ?\\(protection immédiate requise)};
\node[decision, align=center] (vaccin) at (7.5,-4.2){\textbf{Vaccinothérapie}\\(prévention, ex. rougeole)};
\node[question, align=center] (q4) at (0.5,-6.4){Sujet déjà vacciné\\contre cette maladie ?};
\node[decision, align=center] (sero) at (-4,-8.4){\textbf{Sérothérapie}\\(ex. morsure de serpent)};
\node[decision, align=center] (serovacc) at (4.5,-8.4){\textbf{Séro-vaccinothérapie}\\(ex. plaie souillée, tétanos)};

\draw[fleche] (q1) -- node[lbl, above left]{Oui} (antibio);
\draw[fleche] (q1) -- node[lbl, above right]{Non} (q2);
\draw[fleche] (q2) -- node[lbl, above left]{Oui} (q3);
\draw[fleche] (q2) -- node[lbl, above right]{Non} (vaccin);
\draw[fleche] (q3) -- node[lbl, left]{Oui} (q4);
\draw[fleche] (q4) -- node[lbl, above left]{Oui (non vacciné)} (sero);
\draw[fleche] (q4) -- node[lbl, above right]{Non (vacciné)} (serovacc);
\end{tikzpicture}
\legende{Arbre de décision pour choisir la méthode d'aide à l'immunité adaptée à une situation donnée. Les questions (jaune) guident vers la méthode appropriée (vert).}
\end{popfigure}

\begin{methode}[Résoudre un problème de synthèse (méthode en 4 étapes)]
\begin{enumerate}
\item \textbf{Lire attentivement la situation} et repérer l'agent en cause : est-ce une bactérie ou un virus ?
\item \textbf{Déterminer si la situation est préventive ou curative} : la personne est-elle déjà malade, ou cherche-t-on à l'empêcher de le devenir ?
\item \textbf{Déterminer s'il y a urgence} : faut-il une protection immédiate (morsure, plaie souillée) ou peut-on attendre ?
\item \textbf{Choisir et justifier la méthode} : antibiothérapie (bactérie, maladie déclarée), sérothérapie (urgence, pas de vaccin ou sujet déjà immunisé), vaccinothérapie (prévention sans urgence), séro-vaccinothérapie (urgence chez un sujet non vacciné alors qu'un vaccin existe).
\end{enumerate}
\end{methode}

\begin{exoresolu}[Morsure par un chien errant à Yaoundé]
\textbf{Énoncé.} Amina, 13 ans, est mordue à la jambe par un chien errant suspect de rage. Elle n'a jamais été vaccinée contre la rage. Au centre de santé, le médecin lui administre immédiatement un sérum antirabique, puis démarre une série d'injections de vaccin antirabique. (1) Nommer la méthode utilisée. (2) Justifier chacune des deux injections.
\tcblower
\textbf{Solution détaillée.}
\begin{enumerate}
\item La méthode utilisée est la \cle{séro-vaccinothérapie}, car elle associe un sérum et un vaccin.
\item Justification :
\begin{itemize}
\item Le \textbf{sérum antirabique} apporte des anticorps tout prêts qui neutralisent \textbf{immédiatement} le virus de la rage : c'est une immunité \textbf{passive}, indispensable car la rage est mortelle et il y a urgence.
\item Le \textbf{vaccin antirabique} apporte des antigènes qui poussent l'organisme d'Amina à fabriquer \textbf{lui-même} ses anticorps et des cellules mémoires : c'est une immunité \textbf{active}, lente à s'installer mais \textbf{durable}.
\end{itemize}
Ainsi, Amina est protégée tout de suite (grâce au sérum) et durablement (grâce au vaccin). La sérothérapie seule n'aurait protégé que quelques semaines ; la vaccination seule aurait été trop lente face au virus déjà présent.
\end{enumerate}
\end{exoresolu}

\begin{attention}[Erreur fréquente : antibiotique et maladie virale]
Une erreur très répandue consiste à employer le mot « antibiotique » pour « antiviral », ou à réclamer un antibiotique contre une maladie \textbf{virale} (grippe, rougeole, VIH/Sida). Or \textbf{les antibiotiques sont sans effet sur les virus} : ils n'agissent que sur les bactéries. Prescrire un antibiotique contre une maladie virale est non seulement inutile, mais dangereux : cela favorise l'apparition de bactéries résistantes. Au BEPC, ne confonds jamais antibiothérapie (contre les bactéries) et vaccinothérapie (souvent dirigée contre les virus).
\end{attention}

\courssub{Pour aller plus loin (hors programme strict, utile pour la culture scientifique)}

\begin{savaistu}[La variole, seule maladie humaine éradiquée]
Grâce à une immense campagne de \cle{vaccination mondiale} coordonnée par l'Organisation mondiale de la santé (OMS), la \textbf{variole} a été totalement éradiquée : le dernier cas naturel date de 1977 et l'éradication a été officiellement déclarée en \textbf{1980}. C'est à ce jour la \textbf{seule maladie humaine éliminée} de la surface de la Terre. Ce succès illustre la puissance de la vaccination : quand une très grande partie de la population est vaccinée, le microbe ne trouve plus personne à infecter et finit par disparaître.
\end{savaistu}

Deux notions complètent utilement cette leçon, sans être exigibles au BEPC :

\begin{itemize}
\item L'\cle{immunothérapie} désigne, de façon générale, tous les traitements qui agissent sur le système immunitaire pour aider l'organisme à se défendre (la sérothérapie et la vaccinothérapie en sont des exemples).
\item La \cle{couverture vaccinale collective} : lorsqu'une grande proportion d'une population est vaccinée, les microbes circulent difficilement, ce qui protège indirectement même les personnes non vaccinées (nourrissons, personnes fragiles). On parle parfois d'« immunité de groupe ». C'est pourquoi les campagnes de vaccination au Cameroun (contre la rougeole, la poliomyélite, la fièvre jaune) concernent toute la population.
\end{itemize}

\begin{exercice}[À toi de jouer]
Pour chaque situation, indique la méthode d'aide à l'immunité la plus adaptée et justifie ta réponse :
\begin{enumerate}
\item Un élève souffre d'une angine due à un streptocoque (bactérie).
\item Un enfant de 9 mois doit être protégé contre la rougeole, avant toute maladie.
\item Un agriculteur est mordu par un serpent venimeux en allant au champ.
\item Un ouvrier se fait une plaie profonde souillée de terre ; il n'est pas vacciné contre le tétanos.
\end{enumerate}
\end{exercice}

\begin{corrige}[Exercice : À toi de jouer]
\begin{enumerate}
\item \textbf{Antibiothérapie} : la maladie est déclarée et l'agent est une bactérie ; l'antibiotique la détruira directement.
\item \textbf{Vaccinothérapie} : situation préventive, sans urgence ; le vaccin conférera une immunité active durable.
\item \textbf{Sérothérapie} : urgence vitale ; le sérum antivenimeux apporte des anticorps qui neutralisent immédiatement le venin.
\item \textbf{Séro-vaccinothérapie} : urgence (risque de tétanos) chez un sujet non vacciné ; le sérum protège immédiatement et le vaccin assure une protection durable.
\end{enumerate}
\end{corrige}

\begin{aretenir}[L'essentiel de la synthèse]
\begin{itemize}
\item \textbf{Antibiothérapie} : produit chimique, curative, contre les bactéries uniquement, aucune immunité.
\item \textbf{Sérothérapie} : anticorps, curative d'urgence, action immédiate, protection courte, immunité \textbf{passive}.
\item \textbf{Vaccinothérapie} : antigènes, préventive, action lente, protection longue, immunité \textbf{active}.
\item \textbf{Séro-vaccinothérapie} : anticorps + antigènes, urgence chez un sujet non vacciné, protection immédiate \textbf{puis} durable.
\item Pour choisir : maladie déclarée ou risque ? urgence ? bactérie ou virus ? sujet vacciné ?
\end{itemize}
\end{aretenir}

\newpage

\coursec{Activité d'intégration}

\courssub{Objectifs de l'activité}

À l'issue de cette activité d'intégration, tu seras capable de :
\begin{itemize}
\item mobiliser les notions d'\cle{antibiothérapie}, de \cle{sérothérapie}, de \cle{vaccinothérapie} et de \cle{séro-vaccinothérapie} pour analyser des cas cliniques concrets ;
\item interpréter des documents médicaux simples (antibiogramme, courbes de taux d'anticorps, calendrier vaccinal) ;
\item choisir la méthode d'aide au système immunitaire adaptée à une situation donnée ;
\item rédiger et justifier une conclusion argumentée en t'appuyant sur le \cle{principe} de chaque méthode, sa \cle{rapidité} d'action et sa \cle{durée} de protection.
\end{itemize}

\courssub{Situation}

\textbf{Urgence au centre de santé de Ngaoundéré : quel traitement choisir ?}

C'est samedi matin au centre de santé intégré de Ngaoundéré (région de l'Adamaoua). Le médecin de garde, le Dr Amina, reçoit coup sur coup quatre patients. Elle dispose des moyens suivants : des antibiotiques, des sérums thérapeutiques (sérum antivenimeux, sérum antirabique) et les vaccins du Programme Élargi de Vaccination (PEV). Pour chaque patient, elle doit choisir la bonne méthode d'aide au système immunitaire. Tu es stagiaire au centre de santé : aide le Dr Amina à décider, en analysant les quatre dossiers ci-dessous.

\textbf{Dossier 1 — Junior, 8 ans : une angine qui ne passe pas.}

Junior souffre depuis trois jours d'une angine avec forte fièvre (39,5 \textcelsius{}). Le prélèvement de sa gorge, mis en culture au laboratoire, montre la présence d'une bactérie : le streptocoque. On réalise un \cle{antibiogramme} : on dépose sur la culture trois disques imbibés d'antibiotiques différents. Après 24 h, on mesure les auréoles d'inhibition (zones où les bactéries n'ont pas poussé autour de chaque disque) :

\begin{center}
\begin{tabularx}{\linewidth}{|l|X|X|X|}
\hline
\rowcolor{popblueL} Antibiotique testé & Pénicilline (disque A) & Érythromycine (disque B) & Tétracycline (disque C) \\
\hline
Diamètre de l'auréole d'inhibition & 28 mm & 15 mm & 6 mm \\
\hline
Interprétation du laboratoire & Sensible & Sensible & Résistant \\
\hline
\end{tabularx}
\end{center}

\textbf{Dossier 2 — M. Bouba, 45 ans : morsure de serpent dans son champ.}

M. Bouba, paysan de la commune de Belel, a été mordu à la cheville il y a 2 heures par un serpent venimeux alors qu'il travaillait dans son champ de maïs. Il n'a \emph{jamais été vacciné} contre le venin de serpent. Sa jambe gonfle et le venin commence à se diffuser dans son sang. Le centre dispose de sérum antivenimeux (contenant des anticorps antivenin déjà fabriqués).

\textbf{Dossier 3 — Bébé Farida, 6 semaines : la visite de vaccination.}

Farida, en parfaite santé, est amenée par sa maman pour sa première visite de vaccination. Le calendrier du PEV du Cameroun prévoit notamment : BCG et vaccin contre la poliomyélite à la naissance ; à 6, 10 et 14 semaines, le vaccin pentavalent (diphtérie, tétanos, coqueluche, hépatite B, \emph{Haemophilus influenzae}) ; le vaccin contre la rougeole et la fièvre jaune à 9 mois. Farida n'est malade de rien : il s'agit de la \emph{protéger à l'avance}.

\textbf{Dossier 4 — Ibrahim, 14 ans : mordu par un chien errant.}

Ibrahim, élève en classe de 3\textsuperscript{e}, a été mordu à la main par un chien errant suspect de rage. La rage est une maladie virale mortelle dès l'apparition des symptômes, mais le virus met plusieurs semaines à atteindre le système nerveux : on peut encore agir. Le médecin hésite entre trois protocoles. Le laboratoire lui fournit la courbe des taux d'anticorps antirabiques dans le sang d'un patient selon le traitement reçu :

\begin{popfigure}
\begin{center}
\begin{tikzpicture}
\begin{axis}[
width=0.85\linewidth, height=6.2cm,
xlabel={Temps après la morsure (jours)},
ylabel={Taux d'anticorps (unités arbitraires)},
xmin=0, xmax=42, ymin=0, ymax=10,
xtick={0,7,14,21,28,35,42},
ytick={0,2,4,6,8,10},
legend pos=north west,
legend style={font=\small},
grid=major, grid style={popline!50},
]
\addplot[popcurve, popblue] coordinates {(0,0) (2,5) (5,6) (10,4) (20,1.5) (30,0.5) (42,0.2)};
\addlegendentry{Sérum seul}
\addplot[popcurve, poporange] coordinates {(0,0) (7,0.5) (14,2) (21,5) (28,7.5) (35,8) (42,7.8)};
\addlegendentry{Vaccin seul}
\addplot[popcurve, popgreen] coordinates {(0,0) (2,5) (5,6) (10,5.5) (14,5) (21,6.5) (28,8.5) (35,9) (42,8.8)};
\addlegendentry{Séro-vaccination}
\end{axis}
\end{tikzpicture}
\end{center}
\legende{Évolution du taux d'anticorps antirabiques dans le sang d'un patient mordu, selon le traitement : sérum antirabique seul, vaccin antirabique seul, ou association des deux (séro-vaccinothérapie).}
\end{popfigure}

\courssub{Consignes}

Par groupes de quatre élèves, analysez chaque dossier, puis répondez par écrit aux questions suivantes.

\begin{enumerate}
\item \textbf{Dossier 1.} Quelle méthode d'aide à l'immunité convient à Junior ? Justifie ta réponse en rappelant le principe de cette méthode. Quel antibiotique le médecin doit-il prescrire, et pourquoi (appuie-toi sur les mesures de l'antibiogramme) ?
\item \textbf{Dossier 2.} M. Bouba doit-il recevoir un vaccin ou un sérum antivenimeux ? Justifie ton choix en comparant la rapidité d'action des deux méthodes et en précisant ce que contient le sérum.
\item \textbf{Dossier 3.} Farida n'est pas malade : pourquoi la vaccine-t-on quand même ? Explique le principe de la vaccinothérapie et indique si sa protection est immédiate ou différée, courte ou durable.
\item \textbf{Dossier 4.} À l'aide de la courbe, compare l'efficacité dans le temps du sérum seul, du vaccin seul et de l'association des deux. Quel protocole le Dr Amina doit-elle choisir pour Ibrahim ? Justifie en exploitant au moins deux valeurs lues sur le graphique.
\item \textbf{Synthèse.} Construis un tableau comparatif des quatre méthodes (principe, nature de ce qu'on injecte, rapidité d'action, durée de la protection, exemple de situation d'emploi).
\item Rédige une conclusion générale de 8 à 10 lignes expliquant comment le médecin choisit la méthode d'aide à l'immunité selon la situation (maladie déclarée ou non, urgence, protection à long terme).
\end{enumerate}

\courssub{Ressources à mobiliser}

\begin{aretenir}[Les quatre méthodes d'aide au système immunitaire]
\begin{itemize}
\item \textbf{Antibiothérapie} : traitement d'une maladie microbienne (bactérienne) par des \cle{antibiotiques}, substances qui tuent les bactéries ou bloquent leur multiplication. L'\cle{antibiogramme} permet de choisir l'antibiotique efficace : plus l'auréole d'inhibition est grande, plus la bactérie est sensible à cet antibiotique.
\item \textbf{Sérothérapie} : injection d'un \cle{sérum} contenant des anticorps \emph{déjà fabriqués} (par un animal ou un donneur). Action \emph{immédiate} mais \emph{courte} (quelques semaines), car le corps ne fabrique pas ses propres anticorps : c'est une immunité dite \emph{passive}.
\item \textbf{Vaccinothérapie} : injection d'un \cle{vaccin} (microbe tué, atténué ou fragment de microbe) qui pousse l'organisme à fabriquer \emph{lui-même} ses anticorps et à garder en mémoire le microbe. Action \emph{lente} (quelques semaines) mais \emph{durable} (plusieurs années, parfois à vie) : c'est une immunité \emph{active}, utilisée en \emph{prévention}.
\item \textbf{Séro-vaccinothérapie} : association d'un sérum (protection immédiate) et d'un vaccin (protection durable) ; utilisée en urgence contre certaines maladies graves comme la rage ou le tétanos.
\end{itemize}
\end{aretenir}

\begin{attention}[Pièges à éviter]
\begin{itemize}
\item Les antibiotiques sont \emph{inefficaces contre les virus} : ils ne soignent ni la rage, ni la grippe, ni le VIH/Sida. Ils n'agissent que sur les bactéries.
\item Un vaccin ne guérit pas une maladie déjà déclarée et rapide : il agit trop lentement. En urgence (venin, tétanos déclaré), c'est le sérum qui sauve.
\item Ne confonds pas « rapide » et « durable » : le sérum agit vite mais protège peu de temps ; le vaccin agit lentement mais protège longtemps.
\item L'usage abusif des antibiotiques (automédication, traitement interrompu) favorise l'apparition de bactéries \emph{résistantes} : il faut toujours suivre l'ordonnance jusqu'au bout.
\end{itemize}
\end{attention}

\courssub{Démarche et corrigé commenté}

\begin{exoresolu}[Urgence au centre de santé de Ngaoundéré : corrigé]
\textbf{Énoncé.} Pour chacun des quatre dossiers, choisir la méthode d'aide à l'immunité adaptée et justifier la décision médicale.

\tcblower

\textbf{Dossier 1 — Junior : l'antibiothérapie.}

Junior souffre d'une angine \emph{bactérienne} (streptocoque) : la méthode adaptée est l'\cle{antibiothérapie}, car les antibiotiques détruisent les bactéries ou bloquent leur multiplication, aidant ainsi le système immunitaire à éliminer l'infection. Pour choisir le bon antibiotique, on lit l'antibiogramme : plus l'auréole d'inhibition est grande, plus l'antibiotique est efficace contre la bactérie. La pénicilline (disque A) donne la plus grande auréole (28 mm, « sensible »), contre 15 mm pour l'érythromycine et seulement 6 mm pour la tétracycline (« résistant » : la bactérie n'est pratiquement pas gênée). Le Dr Amina prescrit donc la \textbf{pénicilline}.

\textbf{Dossier 2 — M. Bouba : la sérothérapie.}

Le venin se diffuse \emph{déjà} dans le sang de M. Bouba : c'est une urgence. Un vaccin serait inutile ici : il faudrait plusieurs semaines pour que son organisme fabrique ses propres anticorps, et le venin aurait agi bien avant. On injecte donc le \textbf{sérum antivenimeux} : il contient des anticorps antivenin \emph{déjà fabriqués}, qui neutralisent immédiatement le venin. C'est le principe de la \cle{sérothérapie} : action immédiate, mais protection de courte durée (les anticorps injectés sont éliminés en quelques semaines).

\textbf{Dossier 3 — Farida : la vaccinothérapie.}

Farida n'est pas malade : on la vaccine pour la \emph{protéger à l'avance}. Le vaccin introduit dans son organisme des microbes tués ou atténués (ou des fragments de microbes), incapables de provoquer la maladie mais capables de stimuler son système immunitaire : celui-ci fabrique alors ses propres anticorps et garde en mémoire le microbe. Si le vrai microbe arrive plus tard, la défense sera rapide et efficace. La protection met quelques semaines à s'installer (action différée) mais elle est \emph{durable} (plusieurs années). C'est le principe de la \cle{vaccinothérapie}, appliquée ici selon le calendrier du PEV (BCG, poliomyélite, pentavalent, rougeole, fièvre jaune).

\textbf{Dossier 4 — Ibrahim : la séro-vaccinothérapie.}

Lecture de la courbe :
\begin{itemize}
\item \emph{Sérum seul} : le taux d'anticorps monte très vite (environ 5 unités dès le jour 2) mais redescend rapidement (environ 0,5 unité au jour 30) : protection immédiate mais éphémère.
\item \emph{Vaccin seul} : le taux reste quasi nul pendant 7 jours, puis monte lentement (environ 8 unités au jour 35) : protection durable mais trop lente à démarrer, alors que le virus de la rage progresse.
\item \emph{Séro-vaccination} : le taux est élevé dès le début (environ 6 unités au jour 5, grâce au sérum) \emph{et} reste élevé durablement (environ 9 unités au jour 35, grâce au vaccin).
\end{itemize}
Le Dr Amina choisit donc la \textbf{séro-vaccinothérapie} : le sérum antirabique protège Ibrahim immédiatement, le temps que le vaccin antirabique stimule la fabrication de ses propres anticorps, qui assureront une protection durable. C'est le seul protocole qui couvre toute la période de danger.

\textbf{Synthèse (question 5).}

\begin{center}
\begin{tabularx}{\linewidth}{|l|X|X|X|X|}
\hline
\rowcolor{popblueL} Critère & Antibiothérapie & Sérothérapie & Vaccinothérapie & Séro-vaccinothérapie \\
\hline
Ce qu'on administre & Antibiotique & Sérum (anticorps prêts) & Vaccin (microbe tué ou atténué) & Sérum + vaccin \\
\hline
Principe & Détruire les bactéries ou bloquer leur multiplication & Neutraliser immédiatement l'agresseur & Faire fabriquer ses propres anticorps & Combiner les deux actions \\
\hline
Rapidité & Rapide (quelques jours) & Immédiate & Lente (quelques semaines) & Immédiate puis durable \\
\hline
Durée & Le temps du traitement & Courte (quelques semaines) & Longue (années) & Longue \\
\hline
Exemple & Angine bactérienne & Morsure de serpent & PEV (BCG, rougeole\ldots) & Rage, tétanos \\
\hline
\end{tabularx}
\end{center}

\textbf{Conclusion (question 6).} Le choix de la méthode dépend de la situation : si la maladie est bactérienne et déclarée, on utilise l'antibiothérapie guidée par l'antibiogramme ; si l'organisme est en danger immédiat (venin, toxine), on utilise la sérothérapie, seule méthode à action instantanée ; si la personne est en bonne santé, on la protège à l'avance par la vaccinothérapie, qui confère une immunité durable ; enfin, dans les urgences graves où l'on a besoin à la fois d'une protection immédiate et durable (rage, tétanos), on associe les deux dans la séro-vaccinothérapie.
\end{exoresolu}

\courssub{Bilan de l'activité}

Cette activité t'a permis de mettre en évidence les points essentiels de la leçon :
\begin{itemize}
\item il n'existe pas \emph{une} méthode d'aide à l'immunité, mais \emph{quatre}, qui répondent à des situations différentes ;
\item le critère décisif du choix est le couple \cle{rapidité} / \cle{durée} : le sérum agit vite mais protège peu, le vaccin agit lentement mais protège longtemps ;
\item l'\cle{antibiogramme} est un outil précieux pour choisir l'antibiotique efficace et lutter contre l'antibiorésistance ;
\item la lecture de documents (tableaux de mesures, courbes de taux d'anticorps) permet de \emph{justifier} une décision médicale par des données chiffrées, et non par intuition ;
\item enfin, tu as rédigé une conclusion argumentée, compétence exigée au BEPC : toute réponse en SVTEEHB doit être appuyée sur un principe scientifique et, quand c'est possible, sur des valeurs précises.
\end{itemize}

\begin{savaistu}[Le PEV au Cameroun]
Le Programme Élargi de Vaccination, piloté par le Ministère de la Santé publique, offre gratuitement les vaccins essentiels aux enfants dans les centres de santé de tout le pays. Grâce à lui, des maladies autrefois redoutables comme la poliomyélite ont quasiment disparu du Cameroun. La vaccination est l'une des aides à l'immunité les plus efficaces et les moins coûteuses : n'oublie jamais ton carnet de vaccination !
\end{savaistu}

\newpage

\coursec{Méthodes et exercices résolus}

\courssub{Méthode 1 : Identifier le type d'aide à l'immunité dans une situation concrète}

\begin{methode}[Reconnaître antibiothérapie, sérothérapie, vaccinothérapie et séro-vaccinothérapie]
Pour identifier la méthode d'aide au système immunitaire utilisée dans une situation de la vie courante :
\begin{enumerate}
\item \textbf{Repérer ce qui est administré} au malade ou au sujet sain : un antibiotique (pénicilline, streptomycine...), un sérum contenant des anticorps, ou un vaccin (microbe tué ou atténué, toxine inactivée).
\item \textbf{Se demander si le sujet est déjà malade ou non} : si le microbe est déjà présent dans l'organisme, on soigne ; si le sujet est sain, on prévient.
\item \textbf{Distinguer guérir et prévenir} : l'antibiothérapie et la sérothérapie \emph{guérissent} (action immédiate mais de courte durée) ; la vaccinothérapie \emph{prévient} (action lente mais durable, grâce à la mémoire immunitaire).
\item \textbf{Conclure} en nommant la méthode et en justifiant : antibiothérapie = antibiotiques qui tuent ou bloquent les bactéries ; sérothérapie = injection d'anticorps déjà fabriqués ; vaccinothérapie = injection d'un antigène qui pousse l'organisme à fabriquer lui-même ses anticorps ; séro-vaccinothérapie = association des deux (sérum + vaccin).
\end{enumerate}
\begin{center}
\begin{tabularx}{\linewidth}{|l|X|X|X|}
\hline
\rowcolor{popblueL} \textbf{Méthode} & \textbf{Ce qu'on injecte} & \textbf{But} & \textbf{Durée d'action} \\
\hline
Antibiothérapie & Antibiotique & Guérir (bactéries) & Courte \\
\hline
Sérothérapie & Sérum (anticorps) & Guérir vite & Courte \\
\hline
Vaccinothérapie & Vaccin (antigène) & Prévenir & Longue \\
\hline
Séro-vaccinothérapie & Sérum + vaccin & Guérir et prévenir & Courte + longue \\
\hline
\end{tabularx}
\end{center}
\end{methode}

\begin{exoresolu}[Consultation au centre de santé de Bafoussam]
Au centre de santé de Bafoussam, on administre à trois patients des traitements différents :
\begin{itemize}
\item Patient A, atteint d'une angine bactérienne : on lui prescrit de l'amoxicilline (un antibiotique).
\item Patient B, mordu par un serpent : on lui injecte un sérum antivenimeux contenant des anticorps.
\item Patient C, enfant sain : on lui injecte un vaccin contre la rougeole.
\end{itemize}
Pour chaque patient, nomme la méthode d'aide à l'immunité utilisée et justifie ta réponse.
\tcblower
\textbf{Solution détaillée}

\textbf{Patient A : antibiothérapie.} On administre un \cle{antibiotique} (l'amoxicilline) à un malade dont l'organisme contient déjà des bactéries (angine bactérienne). L'antibiotique va tuer les bactéries ou bloquer leur multiplication, ce qui aide le système immunitaire à éliminer l'infection. C'est donc une \textbf{antibiothérapie}, méthode qui \emph{guérit}.

\textbf{Patient B : sérothérapie.} On injecte un \cle{sérum} contenant des \cle{anticorps} déjà fabriqués (contre le venin). Ces anticorps agissent immédiatement en neutralisant le poison, car l'organisme du malade n'a pas le temps de fabriquer ses propres anticorps. C'est donc une \textbf{sérothérapie}, méthode d'action rapide mais de courte durée, qui \emph{guérit}.

\textbf{Patient C : vaccinothérapie.} On injecte à un sujet \emph{sain} un \cle{vaccin}, c'est-à-dire le microbe de la rougeole tué ou atténué (antigène). L'organisme fabrique alors lui-même ses anticorps et garde en mémoire le microbe : en cas de contact ultérieur avec le vrai virus, la défense sera rapide et efficace. C'est donc une \textbf{vaccinothérapie}, méthode qui \emph{prévient} la maladie.

\textbf{Conclusion :} patient A = antibiothérapie ; patient B = sérothérapie ; patient C = vaccinothérapie. Le critère essentiel est la nature de la substance injectée et le but recherché (guérir ou prévenir).
\end{exoresolu}

\courssub{Méthode 2 : Expliquer le principe de l'antibiothérapie}

\begin{methode}[Rédiger l'explication du mode d'action d'un antibiotique]
\begin{enumerate}
\item \textbf{Définir} l'antibiothérapie : méthode qui consiste à administrer des antibiotiques à un malade pour aider son organisme à combattre une infection bactérienne.
\item \textbf{Préciser la cible} : les antibiotiques agissent sur les \emph{bactéries} en les tuant (antibiotiques bactéricides) ou en bloquant leur multiplication (antibiotiques bactériostatiques).
\item \textbf{Citer des exemples} : la pénicilline (découverte par Fleming), la streptomycine, l'amoxicilline.
\item \textbf{Rappeler la limite} : les antibiotiques sont \emph{inefficaces contre les virus} (grippe, rougeole, VIH...).
\item \textbf{Conclure} : affaiblies ou tuées par l'antibiotique, les bactéries sont ensuite éliminées par les défenses immunitaires (phagocytes, anticorps).
\end{enumerate}
\end{methode}

\begin{exoresolu}[Angine et grippe]
Un élève de 3ème déclare : « Mon grand frère a la grippe, je vais lui donner l'antibiotique qui reste dans l'armoire à pharmacie pour le soigner. »
\begin{enumerate}
\item Définis l'antibiothérapie.
\item Ce raisonnement est-il correct ? Justifie ta réponse.
\end{enumerate}
\tcblower
\textbf{Solution détaillée}

\textbf{1. Définition.} L'\cle{antibiothérapie} est la méthode d'aide au système immunitaire qui consiste à administrer des \cle{antibiotiques} (pénicilline, streptomycine, amoxicilline...) à un malade afin de détruire les bactéries responsables de l'infection ou de bloquer leur multiplication.

\textbf{2. Le raisonnement est faux.} La grippe est une maladie causée par un \emph{virus}, et non par une bactérie. Or les antibiotiques n'ont \textbf{aucune action sur les virus} : ils n'agissent que sur les bactéries. Donner un antibiotique à un malade grippé ne le guérira pas. De plus, prendre un antibiotique sans ordonnance et sans nécessité favorise l'apparition de bactéries résistantes, ce qui rendra les antibiotiques moins efficaces plus tard.

\textbf{Conclusion :} l'antibiothérapie ne doit être utilisée que contre les infections \emph{bactériennes}, sur prescription médicale. Contre une maladie virale comme la grippe, elle est inutile et dangereuse pour la collectivité.
\end{exoresolu}

\begin{savaistu}[La pénicilline, une découverte qui a changé la médecine]
C'est en 1928 que le médecin britannique Alexander Fleming remarque par hasard qu'une moisissure, le \emph{Penicillium}, empêche les bactéries de se multiplier autour d'elle dans une boîte de culture. De cette observation naît la \cle{pénicilline}, premier antibiotique utilisé à grande échelle, qui a sauvé des millions de vies. Aujourd'hui encore, au Cameroun, les médecins rappellent qu'il ne faut jamais prendre un antibiotique sans ordonnance ni interrompre un traitement avant la fin : cela favorise l'apparition de bactéries résistantes.
\end{savaistu}

\courssub{Méthode 3 : Expliquer le principe de la sérothérapie}

\begin{methode}[Rédiger l'explication du mode d'action d'un sérum]
\begin{enumerate}
\item \textbf{Définir} la sérothérapie : méthode qui consiste à injecter à un malade un sérum contenant des anticorps déjà fabriqués.
\item \textbf{Expliquer l'origine du sérum} : on injecte d'abord l'antigène (microbe, toxine, venin) à un animal (souvent le cheval) ; celui-ci fabrique des anticorps ; on prélève ensuite son sang et on récupère le sérum riche en anticorps.
\item \textbf{Préciser le mode d'action} : les anticorps injectés neutralisent immédiatement l'antigène (toxine, venin) présent dans l'organisme du malade.
\item \textbf{Caractériser l'action} : rapide (d'où son usage en urgence : morsure de serpent, tétanos, rage) mais de \emph{courte durée}, car l'organisme ne fabrique pas lui-même les anticorps et ne garde pas de mémoire.
\end{enumerate}
\end{methode}

\begin{exoresolu}[Le sérum antitétanique]
Pour obtenir du sérum antitétanique, on injecte à un cheval des doses croissantes de toxine tétanique inactivée. Quelques semaines plus tard, on prélève du sang au cheval et on en extrait le sérum. Ce sérum est injecté à un blessé qui a marché sur un clou rouillé.
\begin{enumerate}
\item Nomme la méthode d'aide à l'immunité utilisée chez le blessé.
\item Explique comment le sérum du cheval a été obtenu.
\item Pourquoi cette méthode est-elle adaptée à une urgence ? Quelle est sa limite ?
\end{enumerate}
\tcblower
\textbf{Solution détaillée}

\textbf{1.} La méthode utilisée est la \textbf{sérothérapie} : on injecte au malade un sérum contenant des anticorps déjà fabriqués.

\textbf{2. Obtention du sérum.} La toxine tétanique inactivée (l'antigène) est injectée au cheval. Le système immunitaire du cheval réagit en fabriquant des \cle{anticorps} spécifiques contre cette toxine. On prélève ensuite le sang du cheval et on en sépare le sérum, qui contient ces anticorps en grande quantité : c'est le sérum antitétanique.

\textbf{3. Avantages et limite.} Cette méthode est adaptée à l'urgence car les anticorps injectés agissent \emph{immédiatement} en neutralisant la toxine tétanique présente dans l'organisme du blessé : on n'a pas le temps d'attendre que son propre organisme fabrique des anticorps (cela prendrait plusieurs jours). Sa limite : la protection est de \emph{courte durée} (quelques semaines), car les anticorps injectés sont progressivement détruits et l'organisme du blessé n'a pas gardé de mémoire immunitaire.

\textbf{Conclusion :} la sérothérapie fournit des anticorps « tout prêts » : action immédiate mais temporaire ; elle guérit mais n'immunise pas durablement.
\end{exoresolu}

\courssub{Méthode 4 : Expliquer le principe de la vaccinothérapie}

\begin{methode}[Rédiger l'explication du mode d'action d'un vaccin]
\begin{enumerate}
\item \textbf{Définir} la vaccinothérapie : méthode qui consiste à injecter à un sujet sain un vaccin pour le protéger contre une maladie.
\item \textbf{Décrire le vaccin} : il contient le microbe tué ou atténué (rendu inoffensif), ou une toxine inactivée, c'est-à-dire un \cle{antigène} incapable de provoquer la maladie.
\item \textbf{Décrire la réponse de l'organisme} : au contact de l'antigène, les lymphocytes fabriquent des anticorps spécifiques ; cette première réponse est lente (quelques jours).
\item \textbf{Insister sur la mémoire immunitaire} : l'organisme conserve des cellules mémoires ; lors d'un contact ultérieur avec le vrai microbe, la réponse est \emph{rapide et intense}, et le microbe est détruit avant de provoquer la maladie.
\item \textbf{Conclure} : la vaccination est une méthode de \emph{prévention} à action durable (plusieurs années, parfois toute la vie) ; c'est le principe des campagnes de vaccination (BCG, rougeole, poliomyélite...) menées au Cameroun.
\end{enumerate}
\end{methode}

\begin{exoresolu}[Pourquoi vacciner les enfants ?]
Lors d'une campagne de vaccination dans un quartier de Yaoundé, une maman demande : « Pourquoi injecter un microbe à mon enfant qui n'est pas malade ? »
\begin{enumerate}
\item Définis la vaccinothérapie.
\item Explique, en trois étapes, ce qui se passe dans l'organisme de l'enfant après la vaccination.
\item Pourquoi dit-on que la vaccination est une méthode de prévention durable ?
\end{enumerate}
\tcblower
\textbf{Solution détaillée}

\textbf{1.} La \cle{vaccinothérapie} est la méthode d'aide au système immunitaire qui consiste à injecter à un sujet sain un \cle{vaccin} (microbe tué ou atténué, ou toxine inactivée) afin de le protéger durablement contre une maladie.

\textbf{2. Ce qui se passe après la vaccination :}
\begin{itemize}
\item \textbf{Étape 1 :} le vaccin introduit dans l'organisme un antigène inoffensif (microbe atténué ou tué) qui ne peut pas provoquer la maladie ;
\item \textbf{Étape 2 :} le système immunitaire reconnaît cet antigène comme étranger et les lymphocytes fabriquent des anticorps spécifiques contre ce microbe ;
\item \textbf{Étape 3 :} l'organisme conserve des cellules mémoires capables de reconnaître ce microbe pendant de longues années.
\end{itemize}

\textbf{3. Prévention durable.} Si l'enfant rencontre plus tard le \emph{vrai} microbe, ses cellules mémoires déclenchent immédiatement une production massive d'anticorps : le microbe est détruit avant de causer la maladie. La protection dure plusieurs années, parfois toute la vie. La vaccination \emph{prévient} donc la maladie au lieu de la guérir.

\textbf{Conclusion :} on vaccine un enfant sain pour « entraîner » son système immunitaire sans risque : c'est la méthode de prévention la plus efficace contre les maladies infectieuses.
\end{exoresolu}

\courssub{Méthode 5 : Expliquer le principe de la séro-vaccinothérapie}

\begin{methode}[Rédiger l'explication d'une injection combinée sérum + vaccin]
\begin{enumerate}
\item \textbf{Définir} la séro-vaccinothérapie : méthode qui consiste à injecter à un malade (ou à un sujet exposé) à la fois un sérum (anticorps) et un vaccin (antigène).
\item \textbf{Justifier le rôle du sérum} : les anticorps injectés agissent \emph{immédiatement} contre le microbe ou la toxine déjà présent(e) : c'est l'action curative d'urgence.
\item \textbf{Justifier le rôle du vaccin} : l'antigène pousse l'organisme à fabriquer ses propres anticorps et à acquérir une mémoire immunitaire : c'est l'action préventive durable.
\item \textbf{Conclure} : la séro-vaccinothérapie associe une protection immédiate (courte durée) et une protection durable ; elle est utilisée par exemple contre le tétanos ou la rage après une exposition au danger.
\end{enumerate}
\end{methode}

\begin{exoresolu}[Morsure par un chien enragé]
Un garçon de Douala est mordu par un chien suspect de rage. À l'hôpital, le médecin lui injecte à la fois un sérum antirabique et un vaccin antirabique.
\begin{enumerate}
\item Nomme cette méthode d'aide à l'immunité.
\item Justifie l'injection du sérum.
\item Justifie l'injection du vaccin.
\end{enumerate}
\tcblower
\textbf{Solution détaillée}

\textbf{1.} On injecte au malade à la fois un sérum (anticorps) et un vaccin (antigène) : c'est la \textbf{séro-vaccinothérapie}.

\textbf{2. Rôle du sérum antirabique.} Le virus de la rage est peut-être déjà entré dans l'organisme du garçon lors de la morsure. Le sérum apporte des \cle{anticorps} déjà fabriqués qui neutralisent \emph{immédiatement} le virus : c'est une action curative d'urgence, indispensable car l'organisme seul mettrait trop de temps à réagir.

\textbf{3. Rôle du vaccin antirabique.} Le vaccin contient le virus de la rage tué ou atténué (antigène inoffensif). Il stimule le système immunitaire du garçon, qui fabrique alors \emph{ses propres anticorps} et acquiert une mémoire immunitaire : la protection devient durable.

\textbf{Conclusion :} la séro-vaccinothérapie combine les avantages des deux méthodes : le sérum agit tout de suite (guérison d'urgence) pendant que le vaccin déclenche l'immunité propre du malade (protection durable). C'est la méthode la plus complète lorsque le danger est déjà présent.
\end{exoresolu}

\courssub{Méthode 6 : Comparer les méthodes et rédiger une conclusion de synthèse}

\begin{methode}[Construire un tableau comparatif et rédiger une conclusion argumentée]
\begin{enumerate}
\item \textbf{Choisir des critères de comparaison} précis : substance administrée, but (guérir ou prévenir), rapidité d'action, durée de la protection, exemples d'application.
\item \textbf{Remplir le tableau} avec un vocabulaire exact : antibiotique / sérum (anticorps) / vaccin (antigène).
\item \textbf{Dégager les points communs} : les quatre méthodes aident le système immunitaire à lutter contre les microbes ; elles nécessitent un avis médical.
\item \textbf{Dégager les différences essentielles} : guérir (antibiothérapie, sérothérapie) versus prévenir (vaccinothérapie) ; action rapide mais brève (sérothérapie) versus lente mais durable (vaccinothérapie).
\item \textbf{Rédiger la conclusion} en 3 phrases : rappel du problème, réponse argumentée avec les notions clés soulignées, ouverture éventuelle (importance de la vaccination et de l'usage raisonné des antibiotiques en santé publique).
\end{enumerate}
\end{methode}

\begin{exoresolu}[Problème de synthèse : choisir la bonne méthode]
Voici quatre situations observées dans un hôpital camerounais :
\begin{itemize}
\item Situation 1 : un enfant atteint de pneumonie bactérienne reçoit un antibiotique.
\item Situation 2 : un accidenté reçoit du sérum antitétanique en urgence.
\item Situation 3 : un nourrisson sain reçoit le vaccin BCG contre la tuberculose.
\item Situation 4 : une personne mordue par un animal enragé reçoit sérum et vaccin antirabiques.
\end{itemize}
\begin{enumerate}
\item Recopie et complète le tableau suivant :

\begin{center}
\begin{tabularx}{\linewidth}{|l|X|X|X|}
\hline
\rowcolor{popblueL} \textbf{Situation} & \textbf{Méthode} & \textbf{Substance injectée} & \textbf{But (guérir / prévenir)} \\
\hline
1 & & & \\
\hline
2 & & & \\
\hline
3 & & & \\
\hline
4 & & & \\
\hline
\end{tabularx}
\end{center}

\item Rédige une conclusion de synthèse comparant la sérothérapie et la vaccinothérapie.
\end{enumerate}
\tcblower
\textbf{Solution détaillée}

\textbf{1. Tableau complété :}

\begin{center}
\begin{tabularx}{\linewidth}{|l|X|X|X|}
\hline
\rowcolor{popblueL} \textbf{Situation} & \textbf{Méthode} & \textbf{Substance injectée} & \textbf{But} \\
\hline
1 & Antibiothérapie & Antibiotique & Guérir \\
\hline
2 & Sérothérapie & Sérum (anticorps) & Guérir \\
\hline
3 & Vaccinothérapie & Vaccin (antigène) & Prévenir \\
\hline
4 & Séro-vaccinothérapie & Sérum + vaccin & Guérir et prévenir \\
\hline
\end{tabularx}
\end{center}

\textbf{2. Conclusion de synthèse.} La sérothérapie et la vaccinothérapie sont deux méthodes d'aide au système immunitaire, mais elles diffèrent par leur principe et leur but. La sérothérapie consiste à injecter des anticorps déjà fabriqués : elle agit immédiatement pour \emph{guérir} un malade, mais sa protection est de courte durée car l'organisme ne garde pas de mémoire. La vaccinothérapie consiste à injecter un antigène (microbe tué ou atténué) qui pousse l'organisme à fabriquer lui-même ses anticorps : son action est lente mais durable, et elle sert à \emph{prévenir} la maladie. Ainsi, on choisit le sérum en urgence devant un danger immédiat, et le vaccin pour protéger durablement les populations, comme le font les campagnes de vaccination organisées au Cameroun.
\end{exoresolu}

\begin{attention}[Les 5 erreurs qui coûtent des points au BEPC]
\begin{enumerate}
\item \textbf{Confondre sérum et vaccin.} Le sérum contient des \emph{anticorps} déjà fabriqués ; le vaccin contient un \emph{antigène} (microbe tué ou atténué). Inverser les deux dans une définition fait perdre tous les points de la question.
\item \textbf{Dire que le vaccin guérit.} La vaccinothérapie \emph{prévient} la maladie chez un sujet sain ; elle ne soigne pas un malade déjà atteint. Écris toujours « prévenir » pour le vaccin et « guérir » pour l'antibiotique et le sérum.
\item \textbf{Affirmer que les antibiotiques tuent les virus.} Les antibiotiques n'agissent que sur les \emph{bactéries} ; ils sont inefficaces contre les maladies virales (grippe, rougeole, VIH/Sida). C'est un piège classique.
\item \textbf{Oublier de justifier la méthode choisie.} Dans une situation concrète, il ne suffit pas de nommer la méthode : il faut justifier par la \emph{substance injectée} et le \emph{but recherché} (guérir ou prévenir).
\item \textbf{Négliger la durée d'action.} Sérothérapie = action rapide mais \emph{courte} ; vaccinothérapie = action lente mais \emph{durable} (mémoire immunitaire). Oublier cette comparaison dans une synthèse fait perdre les points d'argumentation.
\end{enumerate}
\end{attention}

\newpage

\coursec{Exercices}

\courssub{Je vérifie mes connaissances}

\begin{exercice}[QCM : les quatre méthodes d'aide à l'immunité]
Pour chaque question, choisis la (ou les) bonne(s) réponse(s).
\begin{enumerate}
\item L'antibiothérapie consiste à administrer :
\begin{enumerate}
\item[a)] des anticorps fabriqués par un animal ;
\item[b)] des substances capables de tuer les bactéries ou de bloquer leur multiplication ;
\item[c)] des microbes tués ou atténués ;
\item[d)] des vitamines.
\end{enumerate}
\item La sérothérapie utilise :
\begin{enumerate}
\item[a)] des antibiotiques ;
\item[b)] un sérum contenant des anticorps tout prêts ;
\item[c)] des vaccins ;
\item[d)] des globules rouges.
\end{enumerate}
\item La vaccination confère une immunité :
\begin{enumerate}
\item[a)] passive et immédiate ;
\item[b)] active et durable ;
\item[c)] passive et durable ;
\item[d)] active et immédiate.
\end{enumerate}
\item La pénicilline a été découverte par :
\begin{enumerate}
\item[a)] Behring ; \item[b)] Pasteur ; \item[c)] Fleming ; \item[d)] Jenner.
\end{enumerate}
\end{enumerate}
\end{exercice}

\begin{exercice}[Vrai ou faux, à justifier]
Réponds par vrai ou faux, puis justifie ta réponse en une ou deux phrases.
\begin{enumerate}
\item Les antibiotiques sont efficaces contre les maladies virales comme la grippe.
\item La sérothérapie agit rapidement mais sa protection est de courte durée.
\item Le rappel de vaccination permet d'obtenir une réponse immunitaire plus rapide et plus forte.
\item La séro-vaccinothérapie associe l'injection d'un sérum et celle d'un vaccin.
\end{enumerate}
\end{exercice}

\begin{exercice}[Définitions]
Définis en deux ou trois phrases chacun des termes suivants, en précisant pour chacun la substance administrée :
\begin{enumerate}
\item Antibiothérapie ;
\item Sérothérapie ;
\item Vaccinothérapie (vaccination) ;
\item Séro-vaccinothérapie.
\end{enumerate}
\end{exercice}

\begin{exercice}[Immunité active ou passive ?]
Recopie le tableau ci-dessous et place une croix dans chaque case qui convient.
\begin{center}
\begin{tabular}{|l|c|c|c|}
\hline
\rowcolor{popblueL} Méthode & Immunité active & Immunité passive & Aucune immunité \\
\hline
Antibiothérapie & & & \\
\hline
Sérothérapie & & & \\
\hline
Vaccination & & & \\
\hline
Séro-vaccinothérapie & & & \\
\hline
\end{tabular}
\end{center}
Justifie ensuite chacun de tes choix par une phrase.
\end{exercice}

\courssub{Je m'entraîne}

\begin{exercice}[L'expérience de Fleming]
En 1928, Alexander Fleming observe une boîte de Pétri ensemencée de bactéries (staphylocoques) et oubliée sur sa paillasse. Un champignon microscopique, le \emph{Penicillium}, s'y est développé. Autour de la tache de moisissure, on observe une zone claire dépourvue de bactéries ; partout ailleurs, les bactéries ont formé un tapis continu.
\begin{enumerate}
\item Décris précisément les observations faites dans cette boîte de Pétri.
\item Formule une hypothèse expliquant l'absence de bactéries autour du \emph{Penicillium}.
\item Propose une expérience de contrôle permettant de vérifier cette hypothèse.
\item Dégage de cette observation le principe de l'antibiothérapie.
\end{enumerate}
\end{exercice}

\begin{exercice}[Interprétation d'un antibiogramme]
Un enfant de Yaoundé souffre d'une angine bactérienne. Au laboratoire, on réalise un antibiogramme : on dépose sur la culture de la bactérie responsable quatre pastilles imprégnées d'antibiotiques différents (A, B, C, D). Après 24 heures, on mesure le diamètre de l'auréole d'inhibition (zone sans bactéries) autour de chaque pastille. Plus l'auréole est grande, plus la bactérie est sensible à l'antibiotique.
\begin{center}
\begin{tabular}{|l|c|c|c|c|}
\hline
\rowcolor{popblueL} Antibiotique & A & B & C & D \\
\hline
Diamètre de l'auréole (mm) & 6 & 28 & 0 & 18 \\
\hline
\end{tabular}
\end{center}
\begin{enumerate}
\item Quel antibiotique est totalement inefficace contre cette bactérie ? Justifie.
\item Quel antibiotique le médecin doit-il prescrire ? Justifie ta réponse.
\item Que signifie un diamètre d'auréole faible (antibiotique A) ?
\item Explique pourquoi il est dangereux de prendre un antibiotique sans avis médical.
\end{enumerate}
\end{exercice}

\begin{exercice}[La courbe de la vaccination]
On mesure le taux d'anticorps dans le sang d'un enfant vacciné contre la diphtérie. Une première injection (primo-vaccination) est faite au jour 0, puis un rappel au jour 30. Les résultats sont les suivants :
\begin{center}
\begin{tabular}{|l|c|c|c|c|c|c|c|}
\hline
\rowcolor{popblueL} Temps (jours) & 0 & 10 & 20 & 30 & 35 & 40 & 60 \\
\hline
Taux d'anticorps (unités arbitraires) & 0 & 5 & 12 & 8 & 40 & 60 & 55 \\
\hline
\end{tabular}
\end{center}
\begin{enumerate}
\item Trace la courbe du taux d'anticorps en fonction du temps (abscisses : 1 cm pour 5 jours ; ordonnées : 1 cm pour 5 unités). Pour t'aider, un repère gradué est fourni ci-dessous.
\item Compare la réponse après la primo-vaccination et la réponse après le rappel (rapidité, intensité).
\item Explique, en utilisant le rôle des lymphocytes mémoires, pourquoi la deuxième réponse est différente.
\item Quel est l'intérêt du rappel vaccinal ?
\end{enumerate}
\end{exercice}

\begin{popfigure}
\begin{tikzpicture}[scale=0.65]
% Axes
\draw[->, thick] (0,0) -- (13,0) node[right] {Temps (jours)};
\draw[->, thick] (0,0) -- (0,14) node[above] {Taux d'anticorps (u.a.)};
% Graduations abscisses (1cm = 5 jours -> 1 unit = 0.2cm? No, 1cm=5 days. TikZ unit=1cm. So 5 days = 1cm. Max 60 days = 12cm)
% Graduations ordonnées (1cm = 5 u.a. -> 1 u.a. = 0.2cm. Max 60 u.a. = 12cm)
\foreach \x/\xtext in {0/0, 1/5, 2/10, 3/15, 4/20, 5/25, 6/30, 7/35, 8/40, 9/45, 10/50, 11/55, 12/60}
    \draw (\x,0.1) -- (\x,-0.1) node[below] {\xtext};
\foreach \y/\ytext in {0/0, 1/5, 2/10, 3/15, 4/20, 5/25, 6/30, 7/35, 8/40, 9/45, 10/50, 11/55, 12/60}
    \draw (0.1,\y) -- (-0.1,\y) node[left] {\ytext};
% Grid
\draw[popmuted, very thin] (0,0) grid (12,12);
% Points data (x = days/5, y = u.a./5)
% (0,0), (10,5)->(2,1), (20,12)->(4,2.4), (30,8)->(6,1.6), (35,40)->(7,8), (40,60)->(8,12), (60,55)->(12,11)
\foreach \p in {(0,0), (2,1), (4,2.4), (6,1.6), (7,8), (8,12), (12,11)}
    \draw[popblue, fill=popblue] \p circle (0.08);
\end{tikzpicture}
\legende{Repère pour le tracé de la courbe d'évolution du taux d'anticorps après primo-vaccination et rappel.}
\end{popfigure}

\begin{exercice}[L'expérience de Behring]
Behring réalise l'expérience suivante avec le microbe du tétanos :
\begin{itemize}
\item Lot 1 : des cobayes reçoivent une injection de toxine tétanique ; ils meurent tous.
\item Lot 2 : des cobayes reçoivent d'abord la toxine tétanique, puis, quelques heures plus tard, du sérum prélevé sur un mouton préalablement immunisé contre le tétanos ; ils survivent.
\item Lot 3 : des cobayes reçoivent la toxine tétanique puis du sérum d'un mouton non immunisé ; ils meurent.
\end{itemize}
\begin{enumerate}
\item Décris le protocole et les résultats de chaque lot sous forme d'un tableau.
\item Pourquoi le lot 3 est-il indispensable à l'interprétation ?
\item Quelle conclusion peux-tu tirer de ces résultats ?
\item Déduis-en le principe de la sérothérapie et précise le type d'immunité qu'elle confère (active ou passive, rapide ou lente, durable ou temporaire).
\end{enumerate}
\end{exercice}

\courssub{Je me prépare au BEPC}

\begin{exercice}[Problème de synthèse : trois patients à l'hôpital]
Au centre de santé de Bafoussam, le médecin reçoit trois patients le même jour :
\begin{itemize}
\item \textbf{Patient 1} : un élève de 14 ans souffrant d'une angine à streptocoques (maladie bactérienne) ;
\item \textbf{Patient 2} : une femme mordue quelques heures plus tôt par un chien suspecté de rage ;
\item \textbf{Patient 3} : un nourrisson de 9 mois, en bonne santé, qu'il faut protéger contre la fièvre jaune.
\end{itemize}
\begin{enumerate}
\item Pour chaque patient, indique la méthode d'aide au système immunitaire la mieux adaptée (antibiothérapie, sérothérapie, vaccination ou séro-vaccinothérapie).
\item Justifie chaque choix en rédigeant un paragraphe argumenté : nature de la substance administrée, rapidité d'action, durée de la protection, caractère préventif ou curatif.
\item Pour le patient 2, explique pourquoi la vaccination seule ne suffit pas et pourquoi on lui injecte aussi un sérum antirabique.
\item Pour le patient 3, explique pourquoi la sérothérapie serait un mauvais choix pour une protection à long terme.
\end{enumerate}
\end{exercice}

\begin{exercice}[Tableau comparatif des quatre méthodes]
Recopie et complète le tableau ci-dessous en utilisant les mots ou expressions suivants (certains peuvent être utilisés plusieurs fois) : \emph{anticorps tout prêts ; antibiotiques ; microbes tués ou atténués (ou leurs toxines) ; sérum + vaccin ; rapide ; lente (quelques jours) ; courte ; longue ; active ; passive ; préventif ; curatif}.
\begin{center}
\begin{tabularx}{\linewidth}{|l|X|X|X|X|}
\hline
\rowcolor{popblueL} Critère & Antibiothérapie & Sérothérapie & Vaccination & Séro-vaccinothérapie \\
\hline
Substance administrée & & & & \\
\hline
Rapidité d'action & & & & \\
\hline
Durée de la protection & & & & \\
\hline
Type d'immunité & & & & \\
\hline
But (préventif/curatif) & & & & \\
\hline
\end{tabularx}
\end{center}
Réponds ensuite aux questions :
\begin{enumerate}
\item Quelle méthode ne met pas en jeu les anticorps du malade ? Justifie.
\item Quelle méthode procure à la fois une protection immédiate et une protection durable ? Explique pourquoi.
\end{enumerate}
\end{exercice}

\begin{exercice}[Rédaction argumentée : « les antibiotiques soignent la grippe »]
Lors d'une discussion, un élève de 3ème affirme : « Quand j'ai la grippe, ma grand-mère me donne des antibiotiques et je guéris : les antibiotiques soignent donc la grippe. »
\begin{enumerate}
\item Rappelle la nature de l'agent responsable de la grippe.
\item Explique, en t'appuyant sur le mode d'action des antibiotiques, pourquoi l'affirmation de cet élève est fausse.
\item Propose une explication au fait que l'élève guérisse malgré tout après quelques jours.
\item Rédige un court texte argumenté (10 à 15 lignes) destiné à convaincre cet élève de ne plus prendre d'antibiotiques sans ordonnance, en évoquant notamment le danger de l'apparition de bactéries résistantes.
\end{enumerate}
\end{exercice}

\newpage

\coursec{Corrigés des exercices}

\coursec{Leçon N08 : Aide au système immunitaire}

\begin{corrige}[Exercice 1 — QCM : les quatre méthodes d'aide à l'immunité]
\begin{enumerate}
\item \textbf{Réponse b)} : l'antibiothérapie consiste à administrer des substances (les antibiotiques) capables de tuer les bactéries ou de bloquer leur multiplication. Les réponses a) et c) correspondent respectivement à la sérothérapie et à la vaccination.
\item \textbf{Réponse b)} : la sérothérapie utilise un sérum contenant des anticorps tout prêts, fabriqués par un autre organisme (animal immunisé ou donneur).
\item \textbf{Réponse b)} : la vaccination confère une immunité \cle{active} (l'organisme fabrique lui-même ses anticorps et ses lymphocytes mémoires) et \cle{durable} (grâce aux lymphocytes mémoires, la protection dure des années). Elle n'est pas immédiate : la réponse met quelques jours à se mettre en place.
\item \textbf{Réponse c)} : la pénicilline a été découverte par Alexander \cle{Fleming} en 1928. Behring est le père de la sérothérapie, Pasteur celui de la vaccination (rage), Jenner celui du premier vaccin (variole).
\end{enumerate}
\end{corrige}

\begin{corrige}[Exercice 2 — Vrai ou faux, à justifier]
\begin{enumerate}
\item \textbf{Faux.} La grippe est une maladie à \cle{virus}. Or les antibiotiques n'agissent que sur les \cle{bactéries} ; ils sont totalement inefficaces contre les virus.
\item \textbf{Vrai.} Le sérum apporte des anticorps tout prêts qui agissent immédiatement. Mais l'organisme n'a pas fabriqué de lymphocytes mémoires : les anticorps injectés sont éliminés en quelques semaines et la protection disparaît.
\item \textbf{Vrai.} Lors du rappel, les lymphocytes mémoires formés lors de la primo-vaccination reconnaissent immédiatement l'antigène : la réponse secondaire est plus rapide et plus intense (taux d'anticorps plus élevé) que la réponse primaire.
\item \textbf{Vrai.} La séro-vaccinothérapie associe l'injection d'un sérum (anticorps tout prêts, protection immédiate) et celle d'un vaccin (qui déclenche la fabrication d'anticorps et de lymphocytes mémoires, protection durable). Exemple : la rage, le tétanos.
\end{enumerate}
\end{corrige}

\begin{corrige}[Exercice 3 — Définitions]
\begin{enumerate}
\item \textbf{Antibiothérapie} : méthode qui consiste à administrer des \cle{antibiotiques}, substances produites par des champignons microscopiques ou des bactéries (ou fabriquées en laboratoire), capables de tuer les bactéries ou de bloquer leur multiplication. Elle soigne les maladies bactériennes mais ne confère aucune immunité.
\item \textbf{Sérothérapie} : méthode qui consiste à injecter à un malade un \cle{sérum} contenant des \cle{anticorps tout prêts}, fabriqués par un animal immunisé (cheval, mouton) ou par un donneur. Elle confère une immunité passive, immédiate mais de courte durée.
\item \textbf{Vaccinothérapie (vaccination)} : méthode qui consiste à injecter un \cle{vaccin}, c'est-à-dire des microbes tués ou atténués (ou leurs toxines rendues inoffensives, ou des fragments de microbes). L'organisme fabrique alors lui-même ses anticorps et des lymphocytes mémoires : l'immunité est active et durable.
\item \textbf{Séro-vaccinothérapie} : méthode qui associe l'injection d'un \cle{sérum} (anticorps tout prêts, protection immédiate) et celle d'un \cle{vaccin} (protection durable). Elle est utilisée en urgence contre des maladies graves comme la rage ou le tétanos.
\end{enumerate}
\end{corrige}

\begin{corrige}[Exercice 4 — Immunité active ou passive ?]
\begin{center}
\begin{tabular}{|l|c|c|c|}
\hline
\rowcolor{popblueL} Méthode & Immunité active & Immunité passive & Aucune immunité \\
\hline
Antibiothérapie & & & X \\
\hline
Sérothérapie & & X & \\
\hline
Vaccination & X & & \\
\hline
Séro-vaccinothérapie & X & X & \\
\hline
\end{tabular}
\end{center}
\textbf{Justifications :}
\begin{itemize}
\item \textbf{Antibiothérapie} : les antibiotiques tuent directement les bactéries sans faire intervenir le système immunitaire ; ils ne laissent aucune mémoire immunitaire.
\item \textbf{Sérothérapie} : les anticorps sont fabriqués par un autre organisme et simplement reçus ; l'organisme du malade reste passif.
\item \textbf{Vaccination} : l'organisme fabrique lui-même ses anticorps et ses lymphocytes mémoires ; l'immunité est active.
\item \textbf{Séro-vaccinothérapie} : le sérum apporte une immunité passive immédiate, tandis que le vaccin déclenche une immunité active durable.
\end{itemize}
\end{corrige}

\begin{corrige}[Exercice 5 — L'expérience de Fleming]
\begin{enumerate}
\item \textbf{Observations} : dans la boîte de Pétri, les staphylocoques forment un tapis continu partout, sauf autour de la tache de moisissure (\emph{Penicillium}) où l'on observe une zone claire, dépourvue de bactéries.
\item \textbf{Hypothèse} : le \emph{Penicillium} libère dans le milieu une substance qui tue les bactéries ou empêche leur multiplication (cette substance sera appelée \cle{pénicilline}).
\item \textbf{Expérience de contrôle} : on cultive le \emph{Penicillium} seul dans un milieu liquide, on filtre pour éliminer le champignon, puis on dépose le liquide (filtrat) sur une culture de staphylocoques. Si une zone sans bactéries apparaît autour du dépôt, l'hypothèse est vérifiée. Le témoin est une culture de staphylocoques recevant le milieu liquide n'ayant jamais contenu de \emph{Penicillium} : les bactéries doivent s'y développer normalement.
\item \textbf{Principe de l'antibiothérapie} : certaines substances produites par des micro-organismes (les antibiotiques) peuvent détruire les bactéries ou bloquer leur multiplication ; on peut les administrer à un malade pour soigner les maladies bactériennes.
\end{enumerate}
\end{corrige}

\begin{corrige}[Exercice 6 — Interprétation d'un antibiogramme]
\begin{enumerate}
\item L'antibiotique \textbf{C} est totalement inefficace : le diamètre de l'auréole est de \SI{0}{mm}, ce qui signifie que les bactéries se sont multipliées jusqu'au contact de la pastille. La bactérie est \cle{résistante} à cet antibiotique.
\item Le médecin doit prescrire l'antibiotique \textbf{B} : c'est celui qui présente la plus grande auréole d'inhibition (\SI{28}{mm}), donc celui auquel la bactérie est la plus sensible.
\item Un diamètre faible (antibiotique A, \SI{6}{mm}) signifie que la bactérie est peu sensible à cet antibiotique : celui-ci ne bloque que faiblement sa multiplication ; il risque d'être inefficace pour soigner le malade.
\item Prendre un antibiotique sans avis médical est dangereux car : l'antibiotique choisi peut être inefficace contre la bactérie en cause (comme C ici) ; un mauvais usage (dose insuffisante, traitement interrompu) favorise la sélection de \cle{bactéries résistantes}, qui ne pourront plus être tuées par les antibiotiques ; enfin, les antibiotiques sont inutiles et inopérants contre les maladies virales.
\end{enumerate}
\end{corrige}

\begin{corrige}[Exercice 7 — La courbe de la vaccination]
\begin{enumerate}
\item \textbf{Tracé de la courbe} : on place les points de coordonnées (temps ; taux) : (0 ; 0), (10 ; 5), (20 ; 12), (30 ; 8), (35 ; 40), (40 ; 60), (60 ; 55), puis on les relie.
\begin{popfigure}
\begin{tikzpicture}[scale=0.65]
\draw[->, thick] (0,0) -- (13,0) node[right] {Temps (jours)};
\draw[->, thick] (0,0) -- (0,14) node[above] {Taux d'anticorps (u.a.)};
\foreach \x/\xtext in {0/0, 1/5, 2/10, 3/15, 4/20, 5/25, 6/30, 7/35, 8/40, 9/45, 10/50, 11/55, 12/60}
    \draw (\x,0.1) -- (\x,-0.1) node[below] {\xtext};
\foreach \y/\ytext in {0/0, 1/5, 2/10, 3/15, 4/20, 5/25, 6/30, 7/35, 8/40, 9/45, 10/50, 11/55, 12/60}
    \draw (0.1,\y) -- (-0.1,\y) node[left] {\ytext};
\draw[popmuted, very thin] (0,0) grid (12,12);
\draw[popblue, very thick] (0,0) -- (2,1) -- (4,2.4) -- (6,1.6) -- (7,8) -- (8,12) -- (12,11);
\foreach \p in {(0,0), (2,1), (4,2.4), (6,1.6), (7,8), (8,12), (12,11)}
    \draw[popblue, fill=popblue] \p circle (0.08);
\draw[poporange, dashed] (6,0) -- (6,1.6);
\node[poporange, below] at (6,-0.6) {rappel};
\node[popblue] at (3.2,3.6) {réponse primaire};
\node[popblue] at (9.6,9) {réponse secondaire};
\end{tikzpicture}
\legende{Évolution du taux d'anticorps au cours d'une primo-vaccination (jour 0) et d'un rappel (jour 30). La réponse secondaire est plus rapide et plus intense que la réponse primaire.}
\end{popfigure}
\item \textbf{Comparaison} : après la primo-vaccination, la réponse est lente et faible : le taux maximal (12 u.a.) n'est atteint qu'au bout de 20 jours. Après le rappel, la réponse est beaucoup plus rapide (le taux passe de 8 à 40 u.a. en 5 jours) et beaucoup plus intense (maximum de 60 u.a., soit 5 fois plus élevé).
\item \textbf{Explication} : lors de la primo-vaccination, l'organisme fabrique des anticorps et surtout des \cle{lymphocytes mémoires} qui persistent dans le sang. Lors du rappel, ces lymphocytes mémoires reconnaissent immédiatement l'antigène et se multiplient très vite en produisant une grande quantité d'anticorps : c'est la \cle{réponse secondaire}, plus rapide et plus forte.
\item \textbf{Intérêt du rappel} : il augmente fortement le taux d'anticorps et le nombre de lymphocytes mémoires, ce qui assure une protection durable et efficace contre la maladie.
\end{enumerate}
\textbf{Points de vigilance} : ne pas oublier les légendes et les unités des axes ; relier les points dans l'ordre chronologique ; bien distinguer la réponse primaire (lente, faible) de la réponse secondaire (rapide, intense).
\end{corrige}

\begin{corrige}[Exercice 8 — L'expérience de Behring]
\begin{enumerate}
\item \textbf{Tableau du protocole et des résultats} :
\begin{center}
\begin{tabularx}{\linewidth}{|c|X|X|c|}
\hline
\rowcolor{popblueL} Lot & Injection 1 & Injection 2 & Résultat \\
\hline
1 & Toxine tétanique & Aucune & Mort des cobayes \\
\hline
2 & Toxine tétanique & Sérum d'un mouton immunisé contre le tétanos & Survie des cobayes \\
\hline
3 & Toxine tétanique & Sérum d'un mouton non immunisé & Mort des cobayes \\
\hline
\end{tabularx}
\end{center}
\item Le lot 3 est un \cle{lot témoin} : il montre que ce n'est pas le sérum en lui-même qui protège, mais bien le sérum d'un animal \emph{immunisé}. Sans ce lot, on ne pourrait pas attribuer la survie du lot 2 aux anticorps spécifiques du tétanos.
\item \textbf{Conclusion} : le sang d'un animal immunisé contre le tétanos contient des substances (les \cle{anticorps}) capables de neutraliser la toxine tétanique et de sauver un animal intoxiqué.
\item \textbf{Principe de la sérothérapie} : on injecte à un malade un sérum contenant des anticorps tout prêts, fabriqués par un animal préalablement immunisé. Elle confère une immunité \cle{passive} (les anticorps ne sont pas fabriqués par le malade), \cle{rapide} (action immédiate) mais \cle{temporaire} (les anticorps injectés disparaissent en quelques semaines, sans mémoire immunitaire).
\end{enumerate}
\end{corrige}

\begin{corrige}[Exercice 9 — Problème de synthèse : trois patients à l'hôpital]
\begin{enumerate}
\item \textbf{Patient 1} : antibiothérapie. \textbf{Patient 2} : séro-vaccinothérapie (sérum antirabique + vaccin antirabique). \textbf{Patient 3} : vaccination (vaccin antiamaril, contre la fièvre jaune).
\item \textbf{Justifications rédigées} :

\emph{Patient 1} : l'angine à streptocoques est une maladie \cle{bactérienne} déjà déclarée : il faut la soigner (but curatif). On administre des antibiotiques, substances qui tuent les bactéries ou bloquent leur multiplication ; l'action est rapide mais le traitement ne laisse aucune immunité.

\emph{Patient 2} : la rage est une maladie mortelle et le virus est déjà introduit dans l'organisme par la morsure. On pratique la séro-vaccinothérapie : le sérum antirabique apporte des anticorps tout prêts qui neutralisent immédiatement le virus (protection immédiate mais courte), tandis que le vaccin antirabique stimule la fabrication d'anticorps et de lymphocytes mémoires par l'organisme (protection durable). C'est un traitement d'urgence, à la fois curatif et préventif.

\emph{Patient 3} : le nourrisson est en bonne santé ; on cherche à le protéger à l'avance (but préventif). On lui injecte le vaccin contre la fièvre jaune, contenant le virus atténué : son organisme fabrique lui-même ses anticorps et ses lymphocytes mémoires. L'immunité est active, s'installe en quelques jours et dure longtemps.
\item \textbf{Patient 2} : la vaccination seule ne suffit pas car la réponse immunitaire active met plusieurs jours à se mettre en place ; or le virus de la rage peut atteindre le système nerveux avant que les anticorps ne soient fabriqués. Le sérum antirabique fournit des anticorps tout prêts qui agissent immédiatement pour neutraliser le virus, le temps que la vaccination prenne le relais de façon durable.
\item \textbf{Patient 3} : la sérothérapie serait un mauvais choix car elle ne confère qu'une immunité passive de courte durée : les anticorps injectés sont éliminés en quelques semaines et l'organisme ne garde aucune mémoire immunitaire. Le nourrisson ne serait plus protégé ensuite, alors qu'on cherche une protection à long terme.
\end{enumerate}
\textbf{Barème indicatif (sur 10)} : choix corrects des trois méthodes (3 pts) ; justification complète pour chaque patient : substance, rapidité, durée, but (4 pts) ; explication du double intérêt sérum + vaccin pour la rage (1,5 pt) ; argument sur la courte durée de la sérothérapie (1,5 pt).

\textbf{Points de vigilance} : toujours préciser la nature de l'agent (bactérie ou virus) ; distinguer clairement but préventif (avant la maladie) et curatif (maladie déclarée) ; citer les lymphocytes mémoires dès qu'on parle de protection durable.
\end{corrige}

\begin{corrige}[Exercice 10 — Tableau comparatif des quatre méthodes]
\begin{center}
\begin{tabularx}{\linewidth}{|l|X|X|X|X|}
\hline
\rowcolor{popblueL} Critère & Antibiothérapie & Sérothérapie & Vaccination & Séro-vaccinothérapie \\
\hline
Substance administrée & antibiotiques & anticorps tout prêts & microbes tués ou atténués (ou leurs toxines) & sérum + vaccin \\
\hline
Rapidité d'action & rapide & rapide & lente (quelques jours) & rapide (sérum) puis lente (vaccin) \\
\hline
Durée de la protection & courte (aucune immunité) & courte & longue & courte puis longue \\
\hline
Type d'immunité & aucune & passive & active & passive + active \\
\hline
But & curatif & curatif & préventif & curatif et préventif \\
\hline
\end{tabularx}
\end{center}
\begin{enumerate}
\item L'\textbf{antibiothérapie} ne met pas en jeu les anticorps du malade : ce sont les antibiotiques qui détruisent directement les bactéries, sans intervention du système immunitaire. (La sérothérapie non plus ne fait pas fabriquer d'anticorps par le malade, puisqu'elle lui fournit des anticorps tout prêts.)
\item La \textbf{séro-vaccinothérapie} procure à la fois une protection immédiate et durable : le sérum apporte des anticorps tout prêts qui agissent tout de suite (protection immédiate mais courte), tandis que le vaccin déclenche la fabrication d'anticorps et de lymphocytes mémoires par l'organisme (protection longue).
\end{enumerate}
\end{corrige}

\begin{corrige}[Exercice 11 — Rédaction argumentée : « les antibiotiques soignent la grippe »]
\begin{enumerate}
\item La grippe est causée par un \cle{virus} (le virus influenza), et non par une bactérie.
\item Les antibiotiques agissent uniquement sur les \cle{bactéries} : ils détruisent leur paroi ou bloquent leur multiplication. Un virus n'a ni paroi bactérienne ni le mode de multiplication des bactéries : les antibiotiques sont donc totalement inefficaces contre la grippe. L'affirmation de l'élève est fausse.
\item La grippe guérit d'elle-même en quelques jours grâce au système immunitaire de l'élève : ses lymphocytes reconnaissent le virus, fabriquent des anticorps et éliminent les cellules infectées. La guérison n'est donc pas due aux antibiotiques, mais aux défenses immunitaires ; la coïncidence entre la prise d'antibiotiques et la guérison est un piège de raisonnement.
\item \textbf{Texte argumenté (exemple de production attendue)} :

« Tu crois que les antibiotiques t'ont guéri de la grippe, mais c'est faux : la grippe est une maladie à virus, et les antibiotiques n'agissent que sur les bactéries. Si tu as guéri, c'est grâce à ton propre système immunitaire, qui a fabriqué des anticorps contre le virus. Prendre des antibiotiques sans ordonnance est même dangereux : en présence d'un antibiotique mal utilisé (mauvaise dose, traitement interrompu), les bactéries les plus résistantes survivent et se multiplient. On sélectionne ainsi des \cle{bactéries résistantes} que plus aucun antibiotique ne pourra tuer. Le jour où tu auras une vraie maladie bactérienne grave, comme une angine ou une infection, l'antibiotique risque alors d'être inefficace. C'est pourquoi le médecin fait parfois un antibiogramme pour choisir le bon antibiotique. Respecte donc toujours l'avis médical : pas d'antibiotique sans ordonnance, et un traitement suivi jusqu'au bout. »
\end{enumerate}
\textbf{Barème indicatif (sur 10)} : nature virale de la grippe (1 pt) ; mode d'action des antibiotiques et inefficacité sur les virus (2 pts) ; rôle du système immunitaire dans la guérison spontanée (2 pts) ; texte argumenté : idée de résistance bactérienne (2 pts), conseils de bon usage (1 pt), clarté et organisation du texte (2 pts).

\textbf{Points de vigilance} : ne pas écrire que « le virus devient résistant » (c'est la bactérie qui devient résistante aux antibiotiques) ; éviter le raisonnement « j'ai pris X puis j'ai guéri, donc X guérit » ; citer l'antibiogramme comme outil du choix du bon antibiotique.
\end{corrige}

\newpage

\coursec{Fiche bilan}

\begin{popfigure}
\begin{tikzpicture}[mindmap, grow cyclic, every node/.style={concept, font=\small},
  root concept/.append style={concept color=popblue, font=\bfseries\small, minimum size=3.2cm},
  level 1/.append style={level distance=3.6cm, sibling angle=72, font=\scriptsize},
  level 1 concept/.append style={minimum size=2.4cm}]
\node[root concept]{Aide au système immunitaire}
  child[concept color=poporange]{ node {Antibiothérapie : antibiotiques contre les bactéries} }
  child[concept color=popgreen]{ node {Sérothérapie : injection de sérum (anticorps)} }
  child[concept color=poppurple]{ node {Vaccinothérapie : injection de vaccin (antigènes)} }
  child[concept color=poppink]{ node {Séro-vaccinothérapie : sérum + vaccin} }
  child[concept color=popgold]{ node {Immunité : rapide/passive ou lente/active} };
\end{tikzpicture}
\legende{Carte mentale de la leçon : les quatre méthodes d'aide au système immunitaire et le type d'immunité qu'elles procurent.}
\end{popfigure}

\begin{bilan}
\textbf{Définitions clés à connaître par cœur}
\begin{itemize}
  \item \cle{Antibiothérapie} : traitement d'une maladie microbienne à l'aide d'\cle{antibiotiques}, substances qui tuent les bactéries ou bloquent leur multiplication (ex. : pénicilline).
  \item \cle{Sérothérapie} : injection à un malade d'un \cle{sérum} contenant des \cle{anticorps} déjà fabriqués, pour le guérir ou le protéger immédiatement.
  \item \cle{Vaccinothérapie} (vaccination) : injection d'un \cle{vaccin} contenant des \cle{antigènes} (microbes tués ou affaiblis) pour pousser l'organisme à fabriquer lui-même ses anticorps.
  \item \cle{Séro-vaccinothérapie} : injection simultanée d'un sérum (protection immédiate) et d'un vaccin (protection durable).
\end{itemize}

\textbf{Tableau récapitulatif à savoir reproduire}
\begin{center}
\begin{tabularx}{\linewidth}{|l|X|X|X|}
\hline
\rowcolor{popblueL} \textbf{Méthode} & \textbf{Ce qu'on injecte} & \textbf{Type d'immunité} & \textbf{Rapidité / durée} \\
\hline
Antibiothérapie & Antibiotiques & Aucune (ne protège pas) & Action curative, temporaire \\
\hline
Sérothérapie & Sérum (anticorps) & Passive & Immédiate mais courte \\
\hline
Vaccinothérapie & Vaccin (antigènes) & Active & Lente mais durable \\
\hline
Séro-vaccinothérapie & Sérum + vaccin & Passive + active & Immédiate et durable \\
\hline
\end{tabularx}
\end{center}

\textbf{Méthodes express}
\begin{itemize}
  \item Pour choisir la méthode : maladie \emph{bactérienne} déclarée $\rightarrow$ antibiotique ; danger immédiat (morsure, venin, tétanos) $\rightarrow$ sérum ; \emph{prévention} durable $\rightarrow$ vaccin ; urgence + protection longue $\rightarrow$ séro-vaccinothérapie.
  \item Pour rédiger : toujours justifier par le \textbf{type d'immunité} (active = l'organisme fabrique ses anticorps ; passive = anticorps reçus de l'extérieur).
  \item Les antibiotiques sont \textbf{inefficaces contre les virus} (ex. : grippe, VIH) : ne jamais les proposer dans ce cas.
\end{itemize}
\end{bilan}

\begin{aretenir}[Checklist avant le BEPC]
\begin{itemize}
  \item $\square$ Je sais définir l'antibiothérapie et citer un exemple d'antibiotique.
  \item $\square$ Je sais que les antibiotiques agissent sur les bactéries, pas sur les virus.
  \item $\square$ Je sais définir la sérothérapie et dire ce que contient un sérum.
  \item $\square$ Je sais définir la vaccinothérapie et dire ce que contient un vaccin.
  \item $\square$ Je sais définir la séro-vaccinothérapie et citer un cas où on l'utilise (tétanos, rage).
  \item $\square$ Je distingue immunité active (lente, durable) et immunité passive (rapide, courte).
  \item $\square$ Je sais reproduire le tableau comparatif des quatre méthodes.
  \item $\square$ Je sais choisir et justifier la méthode adaptée à une situation concrète.
  \item $\square$ Je rédige mes réponses avec le vocabulaire exact : antigène, anticorps, sérum, vaccin.
  \item $\square$ Je relis ma copie pour vérifier que chaque affirmation est justifiée.
\end{itemize}
\end{aretenir}

\findecours

\end{document}
